% !TEX program = xelatex
% 中文译本 / Chinese translation of:
%   "Almost all zeros of the Riemann zeta function are simple
%    and on the critical line" (paper.tex, v0.92, August 2026)
% 编译: xelatex paper-zh.tex (两遍).
% Overleaf: 菜单 Menu → Compiler 选 XeLaTeX 后 Recompile.
% Compile with XeLaTeX (twice); on Overleaf set Compiler = XeLaTeX.
\documentclass[11pt]{article}
\usepackage[margin=1.1in]{geometry}
\usepackage{amsmath,amssymb,amsthm}
\usepackage{mathtools}
\usepackage{booktabs}
% ---- 中文支持: ctex 优先 (Overleaf / 完整 TeX Live), 本地精简
% ---- 环境自动回退到 fontspec + XeTeX locale 断行 ---------------
\IfFileExists{ctex.sty}{%
  \usepackage[UTF8]{ctex}%
  \usepackage{fontspec}%
}{%
  \usepackage{fontspec}%
  \XeTeXlinebreaklocale "zh"
  \IfFontExistsTF{Noto Sans CJK SC}%
    {\setmainfont{Noto Sans CJK SC}}%
    {\setmainfont{Noto Serif CJK SC}}%
}
\XeTeXlinebreakskip = 0pt plus 1pt minus 0.1pt
\usepackage{xcolor}
\usepackage[colorlinks=true,linkcolor=blue!55!black,citecolor=blue!55!black,urlcolor=blue!55!black]{hyperref}
\numberwithin{equation}{section}
\newtheorem{theorem}{定理}[section]
\newtheorem{lemma}[theorem]{引理}
\newtheorem{proposition}[theorem]{命题}
\newtheorem{corollary}[theorem]{推论}
\theoremstyle{definition}
\newtheorem{hypothesis}[theorem]{假设}
\theoremstyle{remark}
\newtheorem{remark}[theorem]{注记}
\newcommand{\LL}{\Lambda}
\newcommand{\SSS}{\mathfrak{S}}
\newcommand{\e}{\mathrm{e}}
\newcommand{\Si}{\mathrm{Si}}
\title{\bf 黎曼zeta函数的几乎全部零点都是单零点且位于临界线上\thanks{验证状态（请先阅读此条）。\emph{无条件性}：本文中的每一个陈述在数论意义下都是无条件的——没有关于$\zeta$或任何$L$-函数的未证明假设进入任何层级（算术输入是Montgomery的带限配对相关定理、Siegel--Walfisz定理以及Matom\"aki--Radziwi\l{}\l--Tao的平均移位卷积定理；常数是不可有效化的）。下面关于\emph{六个有理常数的验证来源}的等级划分，不涉及任何数学假设。等级划分详见\S\ref{s:verif}和\S\ref{s:tt}。$k\le8$的子塔（标题0.8493/0.9246）处于前驱链的等级：每个被使用的常数都通过两种独立方法在精确有理算术中得到证明。$k=10,12,14$的梯级（标题0.8771/0.9385，0.8964/0.9482，$\mathbf{0.9104/0.9552}$）依赖于\S\ref{s:conv}(vii)中的Toeplitz--迹表示，其基础现已成定理（\S\ref{s:tt}）：模型转移引理、分支相等定理（$m_b(N)N^{b+1}$是单个多项式，对所有$N\ge1$）和奇偶性定理（$L(-N)=(-1)^{b+1}L(N)$）共同将每个新矩确定为精确有理数，由有限多个精确的有限$N$求值确定——这些求值本身已被证明为有理数，通过两个独立实现在预注册的保留点上进行逐位确认，并通过三个独立的$10^{7}$样本测量层进行验证。因此，直到$k=14$的每个梯级所使用的常数都处于证明等级，而密度一的封顶定理（定理~\ref{t:dens1}）仅使用已证明的软输入——决定性、精确对数律以及阶一致工具。本文整体的链等级为\emph{已认证候选者}，非纪录：每个连通常数都被证明是有理数（有理多面体体积的有限带符号和，\S\ref{s:conv}）；$k\le6$塔使用的六个常数已通过这些体积的符号求值证明为精确有理数（$C_5=\tfrac1{36}$，$\{4,2\}=-\tfrac{23}{420}$，$\{6\}=-\tfrac1{126}$，$\{2,2,2\}=\tfrac{32}{105}$，$\{5,2\}=\tfrac18$，$\{2^4\}=\tfrac{1661}{3780}$，$C_7=-\tfrac{17}{360}$，$C_8=\tfrac{157}{4032}$，\S\ref{s:five:conn}，\S\ref{s:seven}；在说明处使用两种独立精确方法）；现在每个被使用的常数都有两个独立确认（\S\ref{s:verif}(d)）——Lean形式化覆盖了恒等式和证书层（内核编译，核心Lean~4，不含任何未完成的证明占位（\texttt{sorry}）；\S\ref{s:verif}(f)），分析层未被形式化，且未经过外部评审。三个写出的条目W1--W3已写出（\S\ref{s:writeouts}），配对层常数已认证为精确有理数（\S\ref{s:five:cert}，附录~\ref{a:cert}），消耗步骤由精确有理对偶证书支持，该证书在无界支撑上有效（引理~\ref{l:dual}）。每一层的精确等级、差异登记表以及程序的验证门（形式化和外部评审先于任何纪录声明）见\S\ref{s:verif}和附录~\ref{a:record}。}}
\author{Hongyi Yang\thanks{多伦多大学。\textit{电子邮箱}：\texttt{hongy.yang@mail.utoronto.ca}}
\and Shihua Yang\thanks{香港大学。\textit{电子邮箱}：\texttt{u3590485@connect.hku.hk}}}
\date{2026年8月\\[4pt]
\small DOI：\href{https://doi.org/10.5281/zenodo.21975236}%
{10.5281/zenodo.21975236}}
\begin{document}
\maketitle

\begin{abstract}
设$N(T,2T)$为$\zeta(s)$的非平凡零点中纵坐标落在$(T,2T]$内的个数，$N_0^{s}$与$N_d$分别为其中在临界线上的单零点数与互异零点数。在标题脚注与\S\ref{s:verif}所述验证分级之下，我们无条件地（常数不可有效化）证明
\[
\lim_{T\to\infty}\frac{N_0^{s}}{N}
=\lim_{T\to\infty}\frac{N_d}{N}=1 :
\]
\emph{在密度意义下，$\zeta$的几乎所有零点都是单零点且位于临界线上}；同样结论对任意固定的本原Dirichlet $L$-函数成立（定理~\ref{t:dens1}）。证明把一座矩阶梯推到其上确界：\cite{C26}的Gabor压缩Weil矩阵的迹矩序列在精确有理算术中求值至第十四阶，具决定性（矩增长至多$(b!)^{2}$，Stieltjes--Carleman），其极限谱测度在原点无原子（精确有限$N$对数律），故所消耗的Christoffel值递减至零。极限处不声明任何有效速率；每个已实例化的梯级都是有效的，最深一级（定理~\ref{t:k14}）给出
\[
\liminf_{T\to\infty}\frac{N_0^{s}}{N}
\;\ge\;1-2\lambda_7\;=\;0.910460\ldots,
\qquad
\liminf_{T\to\infty}\frac{N_d}{N}
\;\ge\;1-\lambda_7\;=\;0.955230\ldots,
\]
其中$\lambda_7$为分子分母各四十位的显式有理数：超过\emph{十分之九}的零点是临界线上的单零点，超过\emph{二十分之十九}的零点互异。算术输入为Montgomery的带限配对相关定理、Siegel--Walfisz定理与Vaughan估计；两条结构定理封闭矩引擎——\emph{分支相等}（$m_b(N)N^{b+1}$是$N$的单一多项式，经全幺模格点多胞体表示与Ehrhart理论证明）与\emph{奇偶性}（$L(-N)=(-1)^{b+1}L(N)$，经Ehrhart--Macdonald自互反）——因此直到$k=14$的每个被消耗常数都是已证明的精确有理数，且任何更高阶只需多项式多次精确有限$N$求值。声明分级逐常数见\S\ref{s:verif}：常数层为已证明；分析链为待外部评审的已认证候选；消耗证书在核心Lean~4中经内核检查。论文附独立的复现与认证包（\url{https://github.com/JoshuaHKU/zeta-density-one-reproduction}）；证伪者注册表（$571$项预注册检查，$61$项触发并转换，全部前向记录）记录了研究轨迹。
\end{abstract}

\medskip
\noindent\textbf{MSC 2020：}11M26（主）；11M06、15B52、
52B20、44A60（次）。\\
\textbf{关键词：}黎曼zeta函数；单零点；临界线；配对相关；
Weil显式公式；Christoffel函数；矩问题；Ehrhart理论；
全幺模矩阵。

\newpage
\tableofcontents
\newpage
\section{引言}\label{s:intro}

$\zeta(s)$的非平凡零点中有正比例位于临界线上，这归功于Selberg \cite{Sel42}；Levinson \cite{Lev74}得到了比例$\tfrac13$，Conrey \cite{Con89}达到了$\tfrac25$，沿Levinson方法的纪录目前为$\kappa>5/12$ \cite{PRZZ}（另见\cite{BCY11,Fen12}）。最近，三分之二定理\cite{C26}通过另一条路径实现：对Montgomery配对相关和的线性代数解读，其中谱信息来自Weil Hermite形式的Gabor压缩$\widetilde G$的迹矩，仅前两矩——在带宽$\lambda\le1$下均可无条件由Montgomery的素数侧求值计算\cite{Mon73,BGST}——就已给出比例$2/3-o(1)$（使用Montgomery--Taylor窗口则为$0.6725$ \cite{Mon75}）。本文将该框架推广到第三至第六矩。由\cite[注1.1]{C26}，仅使用带宽一数据的任何证书不能超过$0.68185$；超越该上限（\cite{C26}第7.5(f)节）的路径是四矩
$\operatorname{tr}\widetilde G^{4}=d\ell_1^{4}(13/4+o(1))$，其价值为$13/18$。

本文在单一文本中攀登此梯，使用同一记号、同一参考文献、同一验证分类账和同一声明等级：证书演算和三矩（\S\ref{s:cert}），逻辑独立的单侧四矩路径达到$0.6916$（\S\ref{s:fm}），精确四矩值$13/18$（\S\S\ref{s:trunc}--\ref{s:four}），第五、六矩及其消耗（\S\S\ref{s:five}--\ref{s:lp}）达到下述定理；然后第七至第十四矩（\S\S\ref{s:seven}--\ref{s:tt}）给出$k=14$证书$0.9104/0.9552$（定理~\ref{t:k14}）；最后是标题定理（\S\ref{s:dens}）：塔的上确界为$1$——矩序列的决定性、原点无原子以及Christoffel极限给出
$\lim N_0^{s}/N=\lim N_d/N=1$
（定理~\ref{t:dens1}），在极限处不声称有效速率，每个有限梯级有效。

\begin{theorem}[主定理；等级：已认证候选者，见\S\ref{s:verif}]\label{t:main}
使用不可有效化常数，
\[
\liminf_{T\to\infty}\frac{N_0^{s}(T,2T)}{N(T,2T)}
\ge \frac{138058329}{162540559}=0.849377\ldots,
\qquad
\liminf_{T\to\infty}\frac{N_d(T,2T)}{N(T,2T)}
\ge \frac{150299444}{162540559}=0.924688\ldots,
\]
且对任意固定的本原Dirichlet $L$函数的零点同样成立。在推导过程中，每个梯级恰好消耗其链所建立的：$m_4\to13/4$双边（因此$13/18-\varepsilon$和$31/36-\varepsilon$）；通过\S\ref{s:fm}中逻辑独立的单侧路径得到$0.6916/0.8458$；仅由$k\le6$子塔（其每个常数都经过符号认证）得到精确无条件梯级
\[
\liminf\frac{N_0^{s}}{N}\ge\frac{2025}{2519},\qquad
\liminf\frac{N_d}{N}\ge\frac{2272}{2519}
\qquad(\text{推论~\ref{c:2025}}).
\]
（在本文研究早期考虑过的四矩加五矩中间梯级已在\S\ref{s:lp}中修复并降级：没有六矩信息时，它在无界谱支撑上不成立——注册条目D6。）
\end{theorem}

定理~\ref{t:main}的依赖结构在\S\ref{s:verif}中明确分级：$13/18$梯级依赖于引理D（\S\ref{s:four}）的经典技术证明和\cite{C26}的消耗层（在\S\ref{s:frame}中回顾）；$2025/2519$梯级额外消耗符号认证常数$C_5=\tfrac1{36}$、$\{4,2\}=-\tfrac{23}{420}$、$\{6\}=-\tfrac1{126}$、$\{2,2,2\}=\tfrac{32}{105}$（\S\ref{s:five:conn}；配对层精确，\S\ref{s:five:cert}）；$0.8493$标题额外消耗第七、八矩层（\S\ref{s:seven}）——符号认证联合常数$\{5,2\}=\tfrac18$、$\{2^4\}=\tfrac{1661}{3780}$、$\{2,2,4\}=-\tfrac{127}{840}$、$\{4,4\}=\tfrac{23}{4536}$、$\{6,2\}=-\tfrac{563}{11340}$，加上$C_7=-\tfrac{17}{360}$和$C_8=\tfrac{157}{4032}$（两者现已通过轨道约化精确运行符号认证，\S\ref{s:seven:conn}）——其余分级条目是三个混合联合常数的第二路径梯级和Lean扩展（\S\ref{s:verif}(d)），而$2025/2519$梯级仍是分级最保守的陈述。本文的任何内容都不涉及Riemann假设本身：比例陈述对$o(N)$的离带零点不敏感；输入（函数方程、显式公式、长度$\le T$的Dirichlet多项式的均值）为那些RH类比为假的客体所满足，且论证只给出下界\cite[\S1.5]{C26}。

\subsection*{与\cite{C26}的关系及新内容}

本文基于\cite{C26}并直接使用其框架而不重复：压缩Weil矩阵、零点侧块结构、尾部估计和前两迹矩均按该文建立的方式使用（\S\ref{s:frame}）。进展有四方面。\emph{(i) 矩梯}：\cite{C26}仅使用$\mu_1\to1$、$\mu_2\to\tfrac43$；本文精确计算$\mu_3\to2$和$\mu_4\to\tfrac{13}4$，并将$\mu_5$--$\mu_8$层在精确有理模型值（$\tfrac{101}{18}$、$\tfrac{640}{63}$、$\tfrac{3439}{180}$、$\tfrac{747361}{20160}$）处求值，其连通Ursell常数通过符号多面体体积积分证明为精确有理数，每一阶的有理层通过重新中心恒等式（\S\S\ref{s:cert}--\ref{s:seven}）以闭式给出。\emph{(ii) 截断技术}：\S\ref{s:trunc}中的$\ell_1$截断单元格分类账结合Siegel--Walfisz/Vaughan闭包，在\cite{C26}中不存在，它使高阶矩能用经典工具计算。\emph{(iii) 消耗}：\cite[引理3.2]{C26}的秩--迹不等式（矩阵形式的$m^{2}\ge2m-1$）被固定矩数据的Christoffel函数（\S\ref{s:lp}）取代：最优对偶多项式是归一化核多项式，其系数为\emph{有理数}，因此消耗是闭式精确证书，在每一度均有效且在无界支撑上成立。\emph{(iv) 验证装置}：预注册证伪检查、只前向的差异登记表和内核编译Lean层（\S\ref{s:verif}）取代了\cite{C26}的传统单一等级展示。常数从$2/3$和$\tfrac56$（互异）上升到$0.8493$和$0.9246$（在所述等级下；$2025/2519$和$2272/2519$在完全符号认证等级下）。

\subsection*{与配对相关路径的关系}

Baluyot、Goldston、Suriajaya和Turnage-Butterbaugh \cite{BGST}证明了Montgomery配对相关方法在弱于RH的假设下也能得到三分之二型结论，而Goldston--Suriajaya \cite{GS}在条件假设下达到了三分之二常数本身；那些是关于\emph{零点}配对相关的条件陈述。本路径不同且无条件：它通过压缩矩阵框架消耗$\LL\star\LL$的$2k$点关联的平均行为——这是素数侧对象，仅由Siegel--Walfisz定理和Vaughan估计控制——对零点没有任何假设。以前仅在这些假设下达到的三分之二阈值，现在无条件地以余量被超过（$2025/2519>4/5>2/3$），而相对于带宽一数据额外消耗的信息正是\S\S\ref{s:cert}--\ref{s:seven}中的第三至第八矩链。

\subsection*{方法概述}

$\widetilde G$的谱测度的矩$m_k$由$k$循环分类账承载：单元格以素数幂模数元组和锁定（偏移）向量为索引，权重$1/(d\ell_1^{k})$，窗口重叠体积提供游走几何。消耗是显式Chebyshev--Markov证书（\S\ref{s:cert}）：谱测度在原点附近的质量仅由矩数据控制，具有精确交换率。三个观察将矩求值化简为经典工具。

\emph{截断（\S\ref{s:trunc}）。}权重$1/\ell_1^{k}$使所有$\ell_1>(\log X)^{B}$的单元格平凡地双边清偿；只剩下对数大小的模数。

\emph{谱框架（\S\ref{s:four}，\S\ref{s:five}）。}锁定方差的$k$循环满足精确恒等式——在$k=4$为离散Parseval恒等式；对每个$k$，\emph{乘积恒等式} $\sum_j N_k(j)^2=\sum_t A_1(t)\cdots A_k(t)$（两个链共享锁定向量当且仅当它们相差一个公共平移）——将其表示为谱$L^2$距离，其每个因子都是\emph{加窗单素数求和}。在对数模数下，每个主项频率都是对数分母的有理数：Siegel--Walfisz定理在双主箱上给出任意对数幂节省，Vaughan估计关闭次主和混合箱，奇异级数梳由下面的局部律精确识别。明显的“孪生困难”锁定关联消解：从不尝试逐锁定求值。

\emph{局部律和锚点折叠（\S\ref{s:local}，\S\ref{s:five}）。}每个循环长度的模$p$结构是一个仿射锁定链，其密度满足四行重合计数闭式并经过穷尽验证。Bell$(5)=52$和Bell$(6)=203$类分类账的五、六矩然后通过机器精确冻结槽恒等式折叠到已认证的四循环锚点上；$m_4$以上唯一真正新的常数是连通Ursell值。

\subsection*{验证纪律和贡献}

本文背后的研究程序以“先证伪”为原则：每个结构性声明都带有预注册数值检查；修正前向记录，从不编辑；声明分级。注册表中61个检查被触发并转换——其中几个直接塑造了本文，包括第38个，它修正了一个已用常数并\emph{提升了}定理（附录~\ref{a:record}，D19）。链的等级、其命名的写出条目以及认证审计结果在\S\ref{s:verif}中陈述；按程序的验证门，形式化和外部评审先于任何纪录声明。本文的数学发展由Claude（Anthropic）——一个大型语言模型——在作者指导的研究程序内完成，作者指定目标、验证纪律和发表决定，并对内容承担全部学术责任。根据现行期刊作者政策，模型不列为作者；其作用在此处及致谢中说明。本文的定理声明以\emph{所述等级}提出：常数层为已证明的精确有理数，分析层为待外部评审的已认证候选——而非已确立的纪录。

\section{框架与记号，遵循\cite{C26}}\label{s:frame}

压缩Weil矩阵$\widetilde G$的构造——Weil Hermite形式、测试族、零点侧块结构、尾部估计和前两迹矩——与三分之二论文\cite{C26}（§§2--5）完全相同。本节回顾记号并陈述矩梯所消耗的\cite{C26}中的五个结果；证明见\cite{C26}，其自身的分析输入是已发表文献\cite{Wei52,IK,Tit86,Mon73,BGST}。我们的贡献从三矩开始（\S\ref{s:cert}）。

\subsection{Weil形式、测试族和$\widetilde G$}\label{s:frame:G}

全文令$l=\log(T/2\pi)$，$\ell_1=l+2\log2-1$（故$N(T,2T)=T\ell_1/2\pi+O(l)$），对固定$\lambda\in(0,1]$：$L=\lambda l$，$X=e^{L}$，$d=\lfloor LT/2\pi\rfloor$。对$f\in C_c^{2}(\mathbb R)$定义$h_f(z)=\widehat f(z)=\int f(u)e^{izu}\,du$；两次分部积分给出
\begin{equation}\label{e:hdecay}
|h_f(x+iy)|\le e^{|y|\Lambda_f}
\min\bigl(\|f\|_1,\|f''\|_1|x+iy|^{-2}\bigr),
\qquad\operatorname{supp}f\subset[-\Lambda_f,\Lambda_f].
\end{equation}
Weil Hermite形式为
$W(f,g)=\sum_{\rho}m_\rho h_f(\gamma_\rho)
\overline{h_g(\gamma_\rho)}$，和遍历互异非平凡零点$\rho=\beta+i\gamma_\rho$（重数$m_\rho$显式；由\eqref{e:hdecay}和$\sum m_\rho|\gamma_\rho|^{-2}<\infty$绝对收敛）。零点多重集在$\rho\mapsto1-\bar\rho$下不变，故$W$为Hermite。令$s=\tfrac12+i\tau$定义素数侧密度
\begin{gather*}
\mu(\tau)=\frac1{2\pi}\mathrm{Re}\,
\frac{\Gamma'}{\Gamma}\Bigl(\frac14+\frac{i\tau}2\Bigr)
-\frac{\log\pi}{2\pi}, \\
\Pi_X(\tau)=\frac{1}{2\pi(\tfrac14+\tau^{2})}
+\frac1\pi\,\mathrm{Re}\,\frac{X^{s}-1}{s}, \\
P_X(\tau)=-\frac1\pi\sum_{n\le X}
\frac{\LL(n)}{\sqrt n}\cos(\tau\log n),
\end{gather*}
并令$\nu_X=\mu+\Pi_X+P_X$。Weil显式公式\cite{Wei52}（采用\cite[定理5.12]{IK}的归一化，应用于$k=f*\tilde g$，$\tilde g(u)=g(-u)$，其变换$h_k=h_fh_{\bar g}$为整函数且具有衰减\eqref{e:hdecay}）给出谱形式
\begin{equation}\label{e:weil}
W(f,g)=\int_{\mathbb R}h_f(\tau)\overline{h_g(\tau)}\,
\nu_X(\tau)\,d\tau,
\qquad\operatorname{supp}f,g\subset[-\tfrac L2,\tfrac L2].
\end{equation}
我们记录$|\Pi_X(\tau)|\le3\sqrt X/(1+|\tau|)$，$|P_X(\tau)|\ll\sqrt X$，$\int_T^{2T}\mu=N(T,2T)+O(l)$以及$\int_T^{2T}\mu^{2}=T\ell_1^{2}/4\pi^{2}+O(l^{2})$（Stirling）。

固定非降$\varrho\in C^{3}(\mathbb R)$满足$\varrho=0$在$(-\infty,0]$，$\varrho=1$在$[1,\infty)$，斜坡宽度$1\le w\le L/8$，以及锥形函数$\phi(u)=\varrho\bigl((L/2-|u|)/w\bigr)$：偶，$0\le\phi\le1$，$\operatorname{supp}\phi=[-L/2,L/2]$，$\phi=1$在$[-L/2+w,L/2-w]$。令$a=\tfrac1L\int\phi^{2}$，$b=\tfrac1L\int\phi^{4}$（故$1-2w/L\le b\le a\le1$），$\Phi=\widehat{\phi^{2}}$，并注意窗口微积分
\begin{equation}\label{e:windowcalc}
(L-2w-|y|)_+\;\le\;g(y):=(\phi^{2}\star\phi^{2})(y)
\;\le\;(L-|y|)_+ .
\end{equation}
临界密度的Gabor系统为
$f_k(u)=\phi(u)e^{-i\tau_ku}$，$\tau_k=T+2\pi k/L$，$0\le k<d$，压缩Weil矩阵为
\[
G_{kl}=W(f_k,f_l)
=\sum_\rho m_\rho\,\widehat\phi(\gamma_\rho-\tau_k)\,
\widehat\phi(\gamma_\rho-\tau_l)
=\int_{\mathbb R}\widehat\phi(\tau-\tau_k)
\widehat\phi(\tau-\tau_l)\,\nu_X(\tau)\,d\tau,
\qquad\widetilde G=G/L,
\]
由任一表达式均为实对称；其归一化迹矩为
$\mu_k=\operatorname{tr}\widetilde G^{k}/(d\ell_1^{k})$。

\begin{lemma}[Gabor系统的Poisson求和]\label{l:gabor}
对所有$\tau,\tau'\in\mathbb R$，
$\sum_{k\in\mathbb Z}\widehat\phi(\tau-\tau_k)
\widehat\phi(\tau'-\tau_k)=L\,\Phi(\tau-\tau')$；特别地
$\sum_{k\in\mathbb Z}
\widehat\phi(\tau-\tau_k)^{2}=aL^{2}$。
\end{lemma}

\begin{proof}
此为\cite[引理2.2]{C26}：临界密度$h=2\pi/L$下的Gabor系统具有无混叠误差的平移不变框架核，适用于任何支撑在长度为$L$的区间内的窗口。
\end{proof}

\subsection{零点侧：惯性、尾部与计数不等式}\label{s:frame:count}

固定$D_0=T^{1/2}$，$I=(T,2T]$，$I'=(T-D_0,2T+D_0]$，将$G=A+E$分解为$\gamma\in I'$内的互异零点（矩阵$A$）和其余部分（$E$）；两者均为实对称。将$I'$内的互异零点分类：$S_1$（简单且在线上；$s_1$），$S_2$（重数>1且在线上；$s_2$），以及离线对$\{\rho,1-\bar\rho\}$（共$p$对），故$\#Z(I')=s_1+s_2+2p$且$N(I')\ge s_1+2s_2+2p$。

\begin{proposition}[块结构]\label{p:inertiaA}
$n_+(\widetilde A)\le s_1+s_2+p$且
$\operatorname{rank}\widetilde A\le\#Z(I')$。
\end{proposition}

\begin{proof}
此为\cite[命题4.1]{C26}：在拉回下的Sylvester惯性——线上零点贡献正$1\times1$块，离线对贡献双曲块，符号为$(1,1)$，且无需求值泛函的独立性。
\end{proof}

\begin{proposition}[尾部]\label{p:tail}
设$A_0$为常数满足$N(t+1)-N(t)\le A_0\log(t+3)$ \cite[定理9.2]{Tit86}。则对$T\ge T_0$，
\[
\|\widetilde E\|_{\mathrm{op}}\;\le\;\theta_0
:=4A_0C_1^{2}X^{1/2}\,\frac{\log(4T)}{D_0^{2}}
\;\ll\;l\,T^{\lambda/2-1},
\qquad C_1=\|\phi''\|_1=\frac{2\|\varrho''\|_1}{w},
\]
且迹范数满足$\|\widetilde E\|_1\le2\theta_0/L$。
\end{proposition}

\begin{proof}
此为\cite[命题4.2]{C26}（锥形函数的$r^{-2}$衰减\eqref{e:hdecay}配合零点密度界）；我们重述显式$\theta_0$，因为\S\ref{s:trunc}的截断清偿和三矩尾部步骤会使用它。
\end{proof}

\begin{proposition}[计数不等式]\label{p:count}
对每个$\theta\ge\theta_0$，
\[
s_1\;\ge\;2\,n_+^{\theta}(\widetilde G)-N(I'),
\qquad
\#Z(I')\;\ge\;n_+^{\theta}(\widetilde G),
\]
因此，利用$N(I'\setminus I)\ll D_0l=o(N)$，
\begin{equation}\label{e:count}
N_0^{s}(T,2T)\ \ge\ 2\,n_+^{\theta}(\widetilde G)
-N(T,2T)-2N(I'\setminus I),
\qquad
N_d(T,2T)\ \ge\ n_+^{\theta}(\widetilde G)
-N(I'\setminus I).
\end{equation}
\end{proposition}

\begin{proof}
此为\cite[命题4.5]{C26}：在阈值$\theta\ge\theta_0$（命题~\ref{p:tail}）下Weyl不等式将$n_+^{\theta}(\widetilde G)$化为$n_+(\widetilde A)$，再结合命题~\ref{p:inertiaA}和重数计数$s_2+p\le\tfrac12(N(I')-s_1)$给出两式。
\end{proof}

这些不等式对条带内任何在$\rho\mapsto1-\bar\rho$下不变的局部有限多重集且满足$N(t+1)-N(t)\ll\log(t+3)$都成立；它们不含算术。所有算术都在迹矩中。

\subsection{前两矩及梯记号}\label{s:frame:moments}

\begin{proposition}[前两矩]\label{p:mu12}
$\operatorname{tr}\widetilde G=aL\,N(T,2T)+O(L\sqrt X)$，且
\[
\operatorname{tr}\widetilde G^{2}
=2\pi bL\int_T^{2T}\mu^{2}
+\frac{T}{\pi}\sum_{n\le X}\frac{\LL(n)^{2}}{n}\,
g(\log n)+O\bigl(Ll\log l\,(l^{2}+X)\bigr).
\]
因此，在端点极限$\lambda\to1^-$，$w/L\to0$下：$\mu_1\to1$且$\mu_2\to\tfrac43$。
\end{proposition}

\begin{proof}
此为\cite[命题5.3和定理5.8]{C26}：迹由Gabor--Poisson折叠得到，二矩由Montgomery在带宽$\le1$下的素数侧配对相关二矩求值得到，在带内无条件\cite{Mon73,BGST}。我们重述它是因为对角和$\sum_{n\le X}(\LL(n)^{2}/n)\,g(\log n)
=\tfrac{L^{3}}{6}(1+O(1/L))$，同样的窗口自相关\eqref{e:windowcalc}，是三矩分类账（定理~\ref{t:m3}）和模型门（\S\ref{s:cert:model}）所消耗的校准。
\end{proof}

\begin{remark}[端点约定]\label{r:endpoint}
端点$\lambda=1$是允许的，因为即使$X\asymp T$，上式的每个次项相对于主项都是$o(1)$；如果读者偏好在整个过程中保留幂节省，可以固定$\lambda<1$然后令$\lambda\to1^-$，得到相同极限\cite[注6.1]{C26}。锥形尾部（\eqref{e:windowcalc}中的$2w/L$亏量）随$w/L\to0$消失；本文所有端点陈述均按此顺序取二重极限（参见\cite[\S7.1]{C26}）。
\end{remark}

由命题~\ref{p:mu12}和\S\ref{s:cert}的分类账，在端点极限下
\begin{equation}\label{e:moments}
\mu_1\to1,\qquad \mu_2\to\tfrac43,\qquad \mu_3\to2,\qquad
\mu_4=\tfrac{197}{60}+R+o(1),
\end{equation}
其中$R$是移位关联类（\S\ref{s:cert:R}），其模型值精确为$-\tfrac1{30}=\tfrac{13}4-\tfrac{197}{60}$。
零点侧消耗是命题~\ref{p:count}的Christoffel函数计数：$\widetilde G$超过消失阈值$\theta_0$（命题~\ref{p:tail}）的特征值个数的下界转换为零计数\eqref{e:count}，而原点附近的质量界仅来自矩数据（\S\ref{s:cert}）。在矩序列$(1,1,\tfrac43,2,\tfrac{13}4)$处，计数精确给出$13/18$和$31/36$；它随每个偶矩单调增加，随奇矩亏量单调减少。

\subsection{重新中心恒等式}\label{s:recentre}

在端点极限下，平坦锥形将Gabor网格精确置于窗口的Nyquist速率，因此框架核在对角线上坍缩（$\Phi(\tau_k-\tau_l)=L\delta_{kl}$，引理~\ref{l:gabor}），$\widetilde G$的平均场块精确为标量。记
\[
P\;:=\;\widetilde G/(L\ell_1)\;-\;I ,
\]
$I$与$P$可交换，因此迹的二项式定理可直接应用：
\begin{equation}\label{e:recentre}
m_k\;=\;\sum_{j=0}^{k}\binom{k}{j}\,\Sigma_j,
\qquad \Sigma_j:=\lim\,\mathbb E\,
\operatorname{tr}P^{j}/d ,
\end{equation}
这是一个\emph{精确}的记账恒等式：中心矩$\Sigma_j$正是\S\S\ref{s:five}--\ref{s:seven}中Bell分类账的完全联合（无单例）类之和，而每一阶的所有含单例类均被二项式系数吸收。\S\S\ref{s:cert}--\ref{s:five}的链反演得到精确中心矩
\begin{equation}\label{e:sigmas}
\Sigma_1=0,\quad \Sigma_2=\tfrac13,\quad \Sigma_3=0,\quad
\Sigma_4=\tfrac14,\quad \Sigma_5=\tfrac1{36},\quad
\Sigma_6=\tfrac{61}{252},
\end{equation}
其中三个非平凡重合经精确验证：$\Sigma_5=C_5$；$\Sigma_6=\{2,2,2\}+\{4,2\}+\{6\}$；以及$m_5$，$m_6$的无锚有理部分$\tfrac{67}{12}$，$\tfrac{39}4$——由\S\ref{s:five}中的完全Bell分类账组装得到——通过一行二项式反演从\eqref{e:recentre}中消去。因此第七、八矩的有理层可\emph{预测}为闭式
\begin{equation}\label{e:rat78}
m_7=\tfrac{685}{36}+\Sigma_7,
\qquad
m_8=\tfrac{217}{6}+8\,\Sigma_7+\Sigma_8,
\end{equation}
且Bell$(7)$/Bell$(8)$分类账组装必须精确落在这些有理数上——预注册门（F-PRED-RAT）通过：$\tfrac{685}{36}+C_7=\tfrac{6833}{360}$和$\tfrac{217}6+8C_7=\tfrac{3221}{90}$再现了独立组装的分类账值。两个无关推导，四个精确一致。

窗口与单元格：窗口$I_m=[X,4X]$和$I_n=(b_2/b_1)I_m$，长度$Y\asymp X$（更一般地$Y\asymp X^{\vartheta}$，$\vartheta\in(\tfrac12,1]$）；素数幂模数$b_1,b_2$，$g=(b_1,b_2)$，$r=b_1/g$，$q=b_2/g$；平移$k\le K\asymp Y/\max(r,q)$，锁定$|j|\le J\asymp X$；$\e(x)=e^{2\pi ix}$；位置格上的加窗素数求和
\[
S_m(\beta)=\sum_{m\in I_m}\LL(m)\,\e(b_2m\beta),
\qquad
S_n(\beta)=\sum_{n\in I_n}\LL(n)\,\e(b_1n\beta).
\]
$N_4(k,j)$表示锁定四元组$(m,\,m-rk,\,n,\,n-qk)$且满足$b_1n-b_2m=j$的$\LL^4$加权计数，$\SSS_4(k,j)\,V(k,j)$为其奇异级数模型，其中$V(k,j)$是单元格的窗口重叠体积——联合窗口中锁定约束下的四元组计数，在消耗归一化后大小$\asymp Y^{2}/(rq\,\ell_1)$——采用四矩单元格分类账（\S\ref{s:fm:R}）的约定，在复现套件中冻结。五情况局部表$\kappa_p(v_p(j);v_p(b_1),v_p(b_2))$定义$\SSS_{b_1,b_2}(j)=\prod_p\kappa_p$，经与剩余计数逐位精确验证，且$\mathbb E_j[\SSS]=1$（\S\ref{s:fm:toolkit}）。

\emph{导入状态。}\cite{C26}是公开预印本；其框架按本节所述精确使用，其自身的分析输入是已发表文献\cite{Wei52,IK,Tit86,Mon73,BGST}。单侧路径的平均移位卷积输入是Matom\"aki--Radziwi\l{}\l--Tao \cite{MRT}，按已发表使用。本文各层的形式化状态在\S\ref{s:verif}(f)中分级：恒等式、局部律和证书层已在核心Lean~4中内核编译；分析层未形式化。

\section{证书层}\label{s:cert}

本节给出将梯子推向\cite[引理3.2]{C26}的秩--迹不等式之外的证书，以及程序的无条件矩记账，全程经机器验证。内容包括：Chebyshev--Markov证书及其精确$\kappa(c)$律；三矩定理；四矩的$197/60$分解及两种记账对账；以及剩余类$R$的两壁结构及其审计装置。

\subsection{Chebyshev--Markov证书与精确交换率}

\begin{lemma}[证书；机器精确验证]\label{l:cert}
设$\nu$是$\mathbb R$上的概率测度，矩为$m_k$。令
\[
Q(x)\;=\;\frac{\bigl(x^2-\tfrac{21}{8}x+\tfrac32\bigr)^2}{(3/2)^2}.
\]
则$Q\ge0$，$Q(0)=1$，$Q\ge1$在$(-\infty,0]$上，$Q$在$[0,\tfrac{21}{16}]$上递减，且
\begin{equation}\label{e:EQ}
\int Q\,d\nu \;=\; \frac{4}{9}\Bigl(m_4-\tfrac{47}{16}\Bigr)
\;+\;R_Q,\qquad
R_Q:=\tfrac{4}{9}\Bigl[-\tfrac{21}{4}(m_3-2)
+\tfrac{633}{64}\bigl(m_2-\tfrac43\bigr)
-\tfrac{63}{8}\,(m_1-1)\Bigr],
\end{equation}
故若$|m_1-1|,|m_2-\tfrac43|,|m_3-2|\le\varepsilon\le1$，则$|R_Q|\le11\varepsilon$。对$0<\theta\le\tfrac1{100}$，
\[
\nu\bigl((-\infty,\theta]\bigr)
\;\le\;\bigl(1+4\theta\bigr)
\Bigl[\tfrac{4}{9}\bigl(m_4-\tfrac{47}{16}\bigr)
+11\varepsilon\Bigr].
\]
若$m_1=1$，$m_2=\tfrac43$，$m_3=2$，$m_4\le\tfrac{13}4$，则$\nu((-\infty,\theta])\le\tfrac5{36}+O(\theta)$，且该界是精确的：支撑在$\{0,a,b\}$上的三原子测度，其中$a=0.840635\ldots$，$b=1.784365\ldots$为方程$x^2-\tfrac{21}8x+\tfrac32=0$的根，达到该界。因此，在$m_4\le\tfrac{13}4+c$下，\S\ref{s:frame}的计数给出精确线性律
\[
\liminf\frac{N_0^{s}}{N}\;\ge\;\frac{13}{18}-\frac{8c}{9},
\qquad
\liminf\frac{N_d}{N}\;\ge\;\frac{31}{36}-\frac{4c}{9}.
\]
\end{lemma}

\begin{proof}
非负性显然，$Q(0)=1$；对$x\le0$，内层二次式$\ge\tfrac32$，故$Q\ge1$；在$[0,a]$上内层二次式从$\tfrac32$递减到$0$且$a>\tfrac1{100}$；直接计算得$Q(\theta)^{-1}\le1+4\theta$对$\theta\le\tfrac1{100}$。展开$Q$并取期望，
\[
\Bigl(\tfrac32\Bigr)^2\!\int Q\,d\nu
=m_4-\tfrac{21}4\,m_3+\Bigl(\tfrac{441}{64}+3\Bigr)m_2
-\tfrac{63}8\,m_1+\tfrac94,
\]
代入$m=(1,\tfrac43,2,m_4)$得$\tfrac49(m_4-\tfrac{47}{16})$；显示的$R_Q$收集一阶偏差，且$\tfrac49(\tfrac{21}4+\tfrac{633}{64}+\tfrac{63}8)\le11$。尖锐性：三原子测度具有矩$(1,\tfrac43,2,\tfrac{13}4)$，由Hankel递推$m_{k+1}=\tfrac{21}8m_k-\tfrac32m_{k-1}$在原子集上得到，且$Q$在$a,b$处双零。线性律由将质量界代入\S\ref{s:frame}的计数得到，如同\eqref{e:EQ}的证明展示：质量$\le\tfrac5{36}+\tfrac{4c}9+O(\theta)$，且$1-2(\tfrac5{36}+\tfrac{4c}9)=\tfrac{13}{18}-\tfrac{8c}9$，$1-(\tfrac5{36}+\tfrac{4c}9)=\tfrac{31}{36}-\tfrac{4c}9$。
\end{proof}

\begin{remark}\label{r:certlp}
引理~\ref{l:cert}的每个恒等式均在复现包（\texttt{certificate\_verification.py}）中符号验证（sympy，精确有理数）并在认证审计中重新运行；根、$5/36$求值、$\kappa(c)$律以及Hankel极值$m_5^{\mathrm{ext}}=\tfrac{177}{32}=5.53125$，$m_6^{\mathrm{ext}}=\tfrac{2469}{256}=9.6445\ldots$精确重现。对每个$c$重新优化证书多项式（线性规划）将破纪录阈值从$c<0.0559$提高到$c<0.0721$（\texttt{moment\_lp\_reopt.py}：$M^{*}=3.3221=\tfrac{13}4\times1.0222$）；\S\ref{s:lp}的矩固定LP是同一机制在全矩序列上运行。
\end{remark}

\subsection{第三矩，无条件}\label{s:cert:m3}

\begin{lemma}[精确$3$循环共振]\label{l:m3res}
设$\lambda\le1$且$n_1,n_2,n_3\le X=T^{\lambda}$为素数幂，满足$|\log(n_1n_2/n_3)|\le C/T$。则对$T\ge T_0(C)$，有$n_1n_2=n_3$；对$3$循环的每个符号模式，类似结论成立。
\end{lemma}

\begin{proof}
$n_1n_2$和$n_3$是正整数，比值$e^{O(1/T)}$；若不等，比值与$1$之差至少$1/n_3\ge T^{-\lambda}$，对大$T$不可兼容，除非$\lambda=1$；边界情况由$|n_1n_2-n_3|\ge1$与$\le Cn_3/T\le C$处理：对$n_3\le T/(2C)$等式再次被强制，而$n_3\in(T/2C,T]$的范围仅通过锥形尾部贡献，在归一化极限中消失（注记~\ref{r:endpoint}）。
\end{proof}

\begin{theorem}[第三矩]\label{t:m3}
对每个固定$\lambda\in(0,1]$和可允许锥形，$\mu_3$在$T\to\infty$时收敛，且在端点极限$\lambda\to1^-$，$\eta\to0^+$下其极限趋于$2$。
\end{theorem}

\begin{proof}[证明（九类分类账；完整细节见存档备忘录\texttt{V1\_completion.md}）]
在Gabor--Poisson折叠（引理~\ref{l:gabor}）后展开$\operatorname{tr}\widetilde G^{3}$对三循环
$\Phi(\tau_1-\tau_2)\Phi(\tau_2-\tau_3)\Phi(\tau_3-\tau_1)$的贡献；尾部修正由迹--H\"older分裂
$\operatorname{tr}(A+E)^3-\operatorname{tr}A^3$控制，其中$\|E\|_{\mathrm{op}}\le\theta_0\ll lT^{\lambda/2-1}$（命题~\ref{p:tail}）且$\|A\|_F^2=O(d\ell_1^2)$：总贡献$o(d\ell_1^3)$。在$\nu_X^{\otimes3}$的九类中，$\nu_X=\mu+\Pi_X+P_X$：（i）含$\Pi_X$的每类在窗口上点态$O(T^{\lambda/2-1})$，可忽略；（ii）单$P$类为$O(\sqrt X\,l^{2}L)$由$\sum\LL(n)n^{-1/2}\le3\sqrt X$；（iii）类$\mu^3$等于$d\ell_1^3(1+O(1/l))$（$\ell_1$归一化精确到$O(1/l)$；在$T=600/38400/10^{8}$处测得$1.0048/1.0014/1.0004$）；（iv）三个$\mu P^2$类中的每一个，由引理~\ref{l:m3res}的$k=2$情形加上对核尾部的Montgomery--Vaughan Hilbert不等式，化简为精确对角和$\sum_{n\le X}(\LL(n)^2/n)\,w(\log n)$，具有与$k=2$时相同的$(L-u)_+$形边缘权重（循环展开的结构常数局部依赖于边和槽，因此由$k=2$情形固定，其总和是Montgomery求值），且$\sum_{n\le X}(\LL^2/n)(L-\log n)=L^3/6\,(1+O(1/L))$：每类贡献$\tfrac13d\ell_1^3(1+o(1))$在端点，三个位置和为$1$；（v）类$P^3$由引理~\ref{l:m3res}强制为精确关系$n_1n_2=n_3$，即单个素数的幂，而$\sum_p\sum_{a+b=c}\log^3p\,p^{-c}=O(1)$使其为$O(Tl)$。求和：$\mu_3\to1+1=2$。交叉检查：同一常数链的$k=2$实例重现Montgomery的$\tfrac13$到四位；模型门$\int O_3C_3=0$精确（\S\ref{s:cert:model}）；有限高度形状$m_3m_1/m_2^2$与模型匹配到$2\%$。两个非标准步骤——命题~\ref{p:tail}的$k=3$尾部适配和局部性加校准固定常数链——在存档备忘录中写出，并按该严谨度使用。
\end{proof}

\subsection{第四矩：$197/60$分类账与类$R$}\label{s:cert:R}

\begin{proposition}[第四矩的可求值类]\label{p:m4classical}
由$k=4$时相同的分类账（$81$类；$\Pi_X$和单$P$类可忽略；配对符号模式由引理~\ref{l:m3res}精确强制；$\{2,2\}$配对权重为显式窗口积分），
\[
\mu_4\;=\;\underbrace{1}_{\mu^4}
\;+\;\underbrace{2}_{\mu^2P^2\ \mathrm{diag}}
\;+\;\underbrace{\tfrac{17}{60}}_{\{2,2\}\ \mathrm{bare\ pairings}}\;+\;R\;+\;o(1)\;=\;\tfrac{197}{60}+R+o(1)
\]
在端点极限下，其中$\tfrac{17}{60}=2\kappa$，$\kappa=\int_{[0,1]^2}xy\,\Sigma_{\mathrm{pair}}O_4
=\tfrac{17}{120}$（已认证$4$循环重叠几何的精确求积），而$R$是移位关联类：受限乘积$N_1=b_1m_1\ne N_2=b_2m_2$（所有腿为素数幂$\le X$）的配对，权重$\LL^4/\sqrt{N_1N_2}$，窗口核$\widehat W(\Delta)=[\sin2T\Delta-\sin T\Delta]/\Delta$在$\Delta=\log N_1/N_2$处，以及显式重叠几何$O_4$（行列不对称；无序单元格权重$Q(\beta_1,\beta_2)+Q(\beta_2,\beta_1)$）；$R$单位为$(1/\pi)\Sigma/(d\ell_1^4)$。$R$的模型值精确为$-\tfrac1{30}=\tfrac{13}4-\tfrac{197}{60}$。
\end{proposition}

\begin{remark}[记账对账；审计轮次25--26]\label{r:books}
程序存档中共存两种记账约定，不可混淆。上面的\emph{裸对角}约定将$\{2,2\}$配对记为$\tfrac{17}{60}$，并将配对级Montgomery平台（$\min(u,1)$对$C_2$的饱和）承载在$R$内部。\emph{正弦模型累积量}约定将$\{2,2\}$记为$2t_{\mathrm{adj}}+t_{\mathrm{opp}}$，连通部分分离。现在锚点已认证为精确有理数（\S\ref{s:five:cert}：$t_{\mathrm{adj}}=\tfrac7{60}$，$t_{\mathrm{opp}}=\tfrac1{30}$），对账精确且对称：
\[
\underbrace{\tfrac{17}{60}}_{\text{裸}}
=\underbrace{\tfrac4{15}}_{2t_{\mathrm{adj}}+t_{\mathrm{opp}}}
+\underbrace{\tfrac1{60}}_{\text{平台}},
\qquad
R_{\mathrm{model}}
=\underbrace{-\tfrac1{60}}_{\text{平台（对-对Wick）}}
+\underbrace{-\tfrac1{60}}_{\text{连通核}}
=-\tfrac1{30},
\]
且$1+2+\tfrac{17}{60}=\tfrac{197}{60}$，$2\cdot\tfrac{17}{120}=\tfrac{17}{60}$精确。存档网格值（$0.2677$，平台$0.0156$，连通$-0.0177$；跨约定测量$-0.01567$在\texttt{m1\_suite.py}中）带有$t_{\mathrm{adj}}$的小网格偏差（引述$0.11717$，精确$\tfrac7{60}=0.11667$；注册条目D7）；精确连通模型值$-\tfrac1{60}=-0.01667$与累积量引擎的独立网格测量$-0.01663$一致，误差$0.2\%$。本文任何地方消耗的总量始终且现在是精确的（$\tfrac{197}{60}$，$-\tfrac1{30}$，$\tfrac{13}4$）；分裂现在也是精确的。以下所有分类账均说明使用哪种约定。
\end{remark}

\subsection{$R$的两壁及其直接测量}\label{s:cert:walls}

本小节内容均为无条件结构或有限高度测量；它是\S\S\ref{s:trunc}--\ref{s:four}证明开发和检查的审计装置。（i）\emph{两个幽灵引理。}平均场四重项为$O(l^{-3})$：$u$-Jacobian精确抵消$1/\sqrt{N_1N_2}$，内层核积分为$-2[\Si(2T/N)-\Si(T/N)]$，在$N\asymp T$处为边界层；簇平方扇区的光滑部分恒为零，因为$\int_{\mathbb R}\widehat W_{\mathrm{in}}\,d\Delta=0$。（ii）\emph{两壁。}因此$R=R_{\mathrm{pp}}+R_{\mathrm{conn}}$：由顶层尺度Poisson对偶格项承载的边缘层（某腿距$X$在$O(1)$内；显式公式机制）和体对象（模型带$s\sim1.5$）。（iii）\emph{分离可观测量，在三个高度认证}：按$e=l-\max u_i$分箱，在$l=5.95/6.64/7.33$处：边缘$-0.0061/-0.0068/-0.0060$（向模型$-0.0156$饱和），体$+0.0004/-0.0012/-0.0015$（在$l\approx6$处符号翻转，单调趋近$-0.0177$）。（iv）\emph{直接测量。}对所有共振四次素因子对进行中间相遇枚举（精确核，类解析，平坦窗口，$\lambda=1$）给出
\[
\begin{array}{lcccc}
l & 4.56 & 5.94 & 7.33 & 8.72\\
R\ (\text{per }d\ell_1^4) & -0.0008 & -0.0057 & -0.0074 & -0.0101
\end{array}
\]
——在每个高度为负，以所有校准门的$1/l$速率趋势；双代码路径审计显示尖锐截断$C=40$低估了$|R|$（在前两个高度全核为$-0.0012$和$-0.0098$），因此这些是量级下界。同一审计认证了分类账常数链在$k=2,3,4$端到端，误差$0.05$--$0.8\%$。$k=1$边缘层的Poisson平稳项以闭式求值（区域1主公式；在三个高度和两个锥形上验证，误差$5$--$18\%$）：正衰减瞬态$\sim c\,l/\ell_1^4$解释了近带后期的符号翻转。

\subsection{模型校准}\label{s:cert:model}

自由费米子累积量引擎通过有序集划分展开$\operatorname{Tr}\log(1+(e^f-1)K)$计算正弦过程的傅里叶场累积量；门：$C_2(v)=\min(|v|,1)$精确，$\int O_3C_3=0$（故$m_3(1)=2$），且$m_4(1)$组装为$3.25103$与$\tfrac{13}4$（网格误差$0.03\%$）。$m_4(1)=\tfrac{13}4$通过三个独立方式固定（累积量组装；GUE模拟$m_{1..4}=1.0000,1.3313,1.9913,3.2213(27)$；由极值常数的Hankel强制）。GUE双窗口Richardson外推测得模型值
\begin{equation}\label{e:guebands}
m_5(1)=5.573\pm0.113,\qquad m_6(1)=10.02\pm0.27,
\end{equation}
（正弦模型/GUE测量——参见差异登记表，附录~\ref{a:record}，条目D2，早期归属来源的更正），且在中心值的六矩证书给出$0.7328/0.8664$：梯子仅通过奇矩的严格正间隙攀登，因为$k\le4$极值已有$m_5=\tfrac{177}{32}$，$m_6=9.6445$。\S\ref{s:five}的组装值位于带\eqref{e:guebands}内。

\section{单侧四矩路径：$0.6916$}\label{s:fm}

本节建立$R$的单侧界，足以超过配对相关天花板$0.68185$，使用经典技术加上\cite{MRT}的平均移位卷积定理。它在逻辑上独立于\S\S\ref{s:trunc}--\ref{s:four}的精确求值，并作为程序首次跨越天花板保留，同时因为其架构（单元格、桥、聚合、清偿分类账）被引理D证明所使用。

\begin{theorem}[单侧四矩界；等级：已认证候选者，附带引理~\ref{l:CL}的数值分析步骤]\label{t:fm}
令$R=R(1)$为命题~\ref{p:m4classical}中的类，则$\limsup_{T\to\infty}R(1)\le\bar R:=0.0066$。因此，由\S\ref{s:frame}的计数在$(1,1,\tfrac43,2,\tfrac{13}4+\Delta)$处，$\Delta=\bar R+\tfrac1{30}=0.039933$，
\[
\liminf\frac{N_0^{s}}{N}\ \ge\ 1-2\Lambda_2(0)=0.691615,
\qquad
\liminf\frac{N_d}{N}\ \ge\ 1-\Lambda_2(0)=0.845807,
\]
其中$\Lambda_2(0)=\dfrac{5/108+\Delta/3}{1/3+4\Delta/3}=0.154193$，且对任意固定的本原Dirichlet $L$函数同样成立。任何$\bar R\le0.0222$都超过$0.68185$天花板；余量为$0.0156$个$R$单位。
\end{theorem}

\noindent（早期草稿打印了$\Lambda_2(0)=0.156293$；认证审计重新计算了闭式，发现所述定理常数$0.6916/0.8458$正确，而该数字是过时的——差异登记表D1，附录~\ref{a:record}。）常数$\bar R$分解为$[\mathrm{core}\le-0.0055]+[\varepsilon_{\mathrm{tail}}\le0.0111]+[\varepsilon_{CL}=0.0010]$，其中$\varepsilon_{CL}$是核心序列的收敛松弛允量（引理~\ref{l:CL}）。

\subsection{工具包：奇异级数、弧恒等式、尾部}\label{s:fm:toolkit}

对$b_1,b_2\ge1$，$j\ne0$，$b_1m_1-b_2m_2=j$的局部密度为$\SSS_{b_1,b_2}(j)=\prod_p\kappa_p$，五情况表经逐位验证；$\mathbb E_j[\SSS]=1$。其弧展开：

\begin{proposition}[弧恒等式]\label{p:S}
$\displaystyle
\SSS_{b_1,b_2}(j)=\sum_{q\ge1}\beta_q\,c_q(j),\qquad
\beta_q=\frac{\mu(q_1)\mu(q_2)}{\varphi(q_1)\varphi(q_2)},\quad
q_i=\frac{q}{(q,b_i)}.$
\end{proposition}

\noindent 在$b_1=b_2=1$时这是\cite[(66)--(67)]{MRT}中使用的经典求值；一般情形由对$\kappa_p$表的$p$局部比较得到。截断尾部满足，对$g=(b_1,b_2)\le(\log M)^{C}$，
\begin{equation}\label{e:X2}
\textstyle\sum_{|j|\le H}\bigl(\SSS-\SSS^{(P)}\bigr)^2(j)\;\ll\;
H\,(\log)^{-B+O(1)},\qquad P=(\log)^{B},
\end{equation}
由$(q_1,q_2)$分级和$\sum_{|j|\le H}c_q(j)^2\ll(H+q)\varphi(q)$。\emph{清偿分类账}（按程序审计纪律保留）：$g>(\log)^{C}$的单元格以已认证成本$0.71\%$权重份额丢弃，$\le3\times10^{-4}$个$R$单位。平凡包络$\sum_{n\le x}\LL(n)\LL(n+h)\ll\SSS(h)x$对$|h|\le x$一致成立，见\cite[推论3.14]{HR}。

\subsection{覆盖区、中带、桥与聚合}\label{s:fm:R}

记$S_i(\alpha)=\sum_{m\sim M_i}\LL(m)\e(b_im\alpha)$。双模数结构由四个存档引理清偿（轮次38--40）：主弧律
$S_i(\alpha)=\frac{\mu(q_i)}{\varphi(q_i)}v_i(\eta)
+O(M_i\exp(-c\sqrt{\log M}))$，其中$q_i=q/(q,b_i)$（$b$一致性免费）；联合近主点的分类（$q\le\min(P^2g,Pb_1,Pb_2)$，有限$b$无关族，其内容为$(q_1,q_2\le P)$分级截断$\SSS^{(P)}$）；伸缩传递（$\gamma_i=b_i\alpha$在\emph{其自身}逼近质量下满足Vaughan估计）；以及平滑领引理对齐Gallagher宽度。清偿后，\cite{MRT}的$b=1$机制可直接应用：

\begin{proposition}[{\cite[定理1.3(i)]{MRT}}，方差形式]\label{p:IM}
对$X^{8/33+\varepsilon}\le H\le X^{1-\varepsilon}$，任意中心$h_0\le X^{1-\varepsilon}$，及任意$A$：
\[
\sum_{|h-h_0|\le H}\Bigl|\sum_{X<n\le2X}\LL(n)\LL(n+h)-\SSS(h)X
\Bigr|^2\;\ll_{A,\varepsilon}\;HX^2(\log X)^{-A}.
\]
\end{proposition}

\noindent 在指数$8/33$处该类的覆盖份额为$0.885$（锚定于存档的$0.797$在指数$\tfrac13$处）。对中带，令$A(b_2,j)=\sum_m\LL(m)\LL((b_2m+j)/b_1)\mathbf1[b_1\mid b_2m+j]$，$MT=\SSS_{b_1,b_2}(j)\,\mathrm{len}/(b_1b_2)$：

\begin{lemma}[色散交换；精确]\label{l:swaps}
$\sum_{b_2,j}A^2$，$\sum A\cdot MT$和$\sum MT^2$精确地是加权孪生和族，分别对应结构化平移$h=qk$（外对$m-m'=rk$），带上的$\SSS$对$\LL$关联，以及$2J\,E_2(b_1,b_2)(\mathrm{len}/b_1b_2)^2(1+o(1))$，其中$E_2=\prod_p\mathbb E_v[\kappa_p^2]$（锚点：$E_2(1,1)=2\prod_{p>2}(1+(p-1)^{-3})$）。
\end{lemma}

\begin{lemma}[桥]\label{l:bridge}
对任意窗口$W$和模数$q$：
$\displaystyle\sum_{k\ne0}\ \sum_{n,n-qk\in W}\LL(n)\LL(n-qk)
=\sum_{a\bmod q}\Bigl[\psi_W(q,a)^2-\sum_{n\in W,\,n\equiv a}
\LL(n)^2\Bigr].$
\end{lemma}

\begin{lemma}[结构化到完全聚合]\label{l:U}
设$W$长度为$Y$，$\mathrm{dev}(h)$为$W$上的（平滑加权）孪生偏差，$Q$为模数集合，平移计数$K_q\le H/q$，$H\le Y^{1-\varepsilon}$。若$\sum_{h\le H}|\mathrm{dev}|^2\ll_{A'}HY^2(\log)^{-A'}$对所有$A'$成立，则对每个$A$
\[
\sum_{q\in Q}\sum_{k\le K_q}|\mathrm{dev}(qk)|^2
\;\ll_A\;HY^2(\log Y)^{-A},
\]
在$Q$上一致——无模数因子损失。
\end{lemma}

\begin{proof}
左边为$\sum_{h\le H}\nu(h)|\mathrm{dev}(h)|^2$，$\nu(h)\le\tau(h)$。在$\nu\le(\log Y)^{C}$处分：通常部分花费$(\log)^{C}$对任意$A'$输入；对稀部分使用包络$|\mathrm{dev}(h)|\le9\SSS(h)Y$及$\sum_{h\le H,\tau(h)>V}\tau(h)\ll H(\log)^3/V$在$V=(\log)^{C}$。取$C=A+4$，$A'=2A+4$。
\end{proof}

\noindent 带消费者仅对平移变量使用一次Cauchy--Schwarz（任何跨族的Cauchy--Schwarz会付出其Poisson地板，测量超出预算$50$--$60$倍，从未使用）。胶层（焊接权重$w_k(n)=\sum_{m\in I(n)}\LL(m)\LL(m-rk)$）通过主项的精确分解、对主弧的胶上Abel求和，以及对次弧的\cite[命题5.4，附录A]{MRT}除数有界包络来闭合。

\subsection{确定性主项与连续极限}\label{s:fm:mains}

写$\SSS_{b_1,b_2}=\bar D+g_P+g_{\mathrm{tail}}$（$\bar D$为精确有限局部均值；$g_P$为$q\le P$振荡部分；$P=40$）。

\begin{lemma}[部分和恒等式]\label{l:PS}
对每个$M\ge1$，
$\displaystyle\sum_{j\le M}(\SSS_{b_1,b_2}-\bar D)(j)=
-\sum_{d\ge1}d\,\bigl\{\tfrac Md\bigr\}\,
\gamma^{\mathrm{free}}_d$，其中$\gamma$为显式Euler局部，对权重$d\min(1,M/d)$绝对收敛。
\end{lemma}

\begin{lemma}[有限$J$求值与一致尾部]\label{l:E3}
对每个有限$J$和$\theta$，
$\sum_{j\le J}\mathrm{tail}_P(j)\,\mathrm{row}_j(\theta)
=\sum_{d\ge1}\gamma_d^{>P}\,S_0(d\theta;\lfloor J/d\rfloor)$
且$S_0(x;K)=\sum_{k\le K}\frac{\sin2xk-\sin xk}{k}$；故$\theta$一致界$\le C_{\mathrm{saw}}\sum_d
|\gamma_d^{>P}|$。任何$J\to\infty$交换都未执行。
\end{lemma}

\begin{lemma}[已认证锯齿常数]\label{l:S0}
$|S_0(x;K)|\le C_{\mathrm{saw}}:=2.10$对所有$x$和$K\ge1$成立。
\end{lemma}

\begin{proof}
$K\le64$：严格网格验证结合Lipschitz界$|\partial_xS_0|\le3K$认证$\max|S_0|\le1.8675$。$K>64$：Dirichlet表示在奇异集附近用正弦积分比较，远离时用第二中值定理，得$\le2.05$；在$K=128,512,4096$处样本为$1.842,1.849,1.851$，从下方趋近$\Si(\pi)=1.85194$。（在认证审计中重新运行：相同值。）
\end{proof}

\noindent 将一致尾部对已认证几何进行区域积分得$\varepsilon_{\mathrm{tail}}\le0.0115$（$T=2400$），$0.0111$（$T=4800$），递减；求值核心（精确有限形式）为$\mathrm{core}(2400)=-0.00549$，$\mathrm{core}(4800)=-0.0064$，对管道锚点（$-0.00545$ vs 套件$-0.00544$）认证。

\begin{lemma}[核心的连续极限；等级：带计算常数的证明形态——见下方等级说明]\label{l:CL}
$\mathrm{core}(T)$在$T\to\infty$时收敛到
$C_{\mathrm{core}}=\sum_{b_1\le b_2}\frac{\LL(b_1)\LL(b_2)}
{b_1b_2}F_\infty(b_1,b_2)$，一个绝对收敛的可计算连续锯齿积分级数，且
$|\mathrm{core}(T)-C_{\mathrm{core}}|\ll(\log T)^{-1}$；该集合坍缩到通用$(1,1)$类（Mertens质量集中在大素数对），其斜率和常数经计算而非拟合：部分和斜率在连续三个十倍程为$0.99684\to0.99962\to0.99996$（经典律，精确为$1$——现已\emph{证明}：Dirichlet级数为$\zeta(s+1)E(s)$且$2C_2E(0)=1$，命题~\ref{p:c0}），常数$c_0=-1.4228$现已\emph{以闭式确定}：$c_0=\gamma-2=-1.4227843\ldots$（Euler--Mascheroni常数$\gamma$；命题~\ref{p:c0}，在$M=2\times10^{7}$处验证到$10^{-5}$）；$\theta$梯$c^{*}=-0.5433$，十倍程间稳定，与存档独立$-0.55(2)$匹配（$c^{*}=\log\pi+\gamma-2-\overline{\mathrm{osc}}$在存档F-QC5门的钉定中，故$c^{*}$的剩余计算部分是平均振荡$\overline{\mathrm{osc}}\approx0.2652$，此处未确定）。连续求积通过$1/l$中的Richardson外推（$l=40,60,100,160$）得到$C_{\mathrm{core}}=-0.0209\pm0.0026$，独立确认存档梯$-0.020(2)$——两条不相交路径误差$4\%$。
\end{lemma}

\noindent\emph{等级说明。}收敛结构（缩放变量Riemann和，普适性坍缩）是分析论证；坍缩级数的斜率和常数值以及$C_{\mathrm{core}}$的求积是计算常数，带有给定带而非认证包络。定理~\ref{t:fm}仅使用保守的$\mathrm{core}\le-0.0055+\varepsilon_{CL}$，$\varepsilon_{CL}=0.0010$，针对测得的递减网格；求积的认证写出是命名开放条目（\S\ref{s:verif}）。代入计算出的$C_{\mathrm{core}}$将得到$\bar R=-0.0098\pm0.0026<0$，从而$0.7031/0.8515$；该升级\emph{在本文任何地方均未使用}——更强烈地：\emph{本节以外的任何结果均不消耗本节任何常数}——定理~\ref{t:main}、\ref{t:k14}和~\ref{t:dens1}通过\S\S\ref{s:trunc}--\ref{s:four}起的精确路径进行，独立于引理~\ref{l:CL}；引理D从\S\ref{s:fm}借用的是其已认证的\emph{架构}（单元格、桥、聚合、清偿分类账），而非任何数值。

\subsection{组装}\label{s:fm:assembly}

在$\lambda=1$，一致全核约定下：
\[
R(1)\;\le\;\mathrm{core}(T)+\varepsilon_{\mathrm{tail}}+o(1)
\;\le\;(-0.0055+0.0010)+0.0111+o(1)\;=\;0.0066+o(1),
\]
其中$o(1)$收集\S\ref{s:fm:R}的覆盖区、带和低区偏差（低区上限在操作点的已认证有限高度价格为$0.0163$，在LP预算$0.0387+0.0177$内）。将$\bar R=0.0066$代入计数得定理~\ref{t:fm}。\qed

\section{截断清偿}\label{s:trunc}

\begin{lemma}[单元格质量界]\label{l:mass}
在锁定和平移范围内一致地，$k$循环单元格的$\LL^{2k}$加权、消耗归一化质量在参数$\ell_1$处为$\ll Y^{2}(\log Y)^{c}\,\ell_1$。
\end{lemma}

\begin{proof}
对固定单元格，每个锁定配置由$(m,k)$以及锁定的剩余数据决定；窗口中的$(m,k)$对计数$\ll YK\ll Y^{2}/\max(r,q)$，每个权重$\ll(\log Y)^{2k}$，锁定约束从$n$范围中选择$1/\ell_1$比例，除数有界重数；消耗归一化（除以分类账体积$V\asymp Y^{2}/(rq\ell_1)$）最多留下一个因子$\ell_1$。数值分类账面测度聚合质量在$\ell_1$上基本平坦（拟合指数$-0.11$），完全在界内。
\end{proof}

\begin{lemma}[重数]\label{l:mult}
具有单元格参数$\ell_1=\ell$的无序素数幂元组数量为$\ll_{\varepsilon}\ell^{\varepsilon}$（除数有界）。
\end{lemma}

\begin{proposition}[截断--消耗兼容性]\label{p:consume}
令$F$为所消耗的$k$循环单元格聚合，$k\ge4$，$F^{(P)}$为其限制$\ell_1\le P$。则
\[
\bigl|F-F^{(P)}\bigr|\;\ll_{\varepsilon}\;
Y^{2}(\log Y)^{c}\,P^{-(k-2)+\varepsilon}
\qquad\text{双边。}
\]
因此，对每个$A_0$存在$B$（任意固定$B>(c+A_0)/(k-2)$即可）使得证明模数$\le P=(\log X)^{B}$的方差界即证明了消耗形式的界，总截断代价$\ll(\log X)^{-A_0}$：实际消耗的仅为模数$\le(\log X)^{B}$。
\end{proposition}

\begin{proof}
权重$\times$质量$\times$重数$=\ell^{-k}\cdot\ell\cdot\ell^{\varepsilon}$对$\ell>P$求和。该界基于绝对值，故双边。存档D1层独立认证$k=4$时平方余项分级（$\mathbb E_j[(F-F^{(P)})^2]\ll P^{-1-\delta}$；尾部局部$t_p^{2}\sim4/p^{2}$）。程序中并存两种截断且不可混淆：\S\ref{s:fm:mains}中$C_{\mathrm{full}}$求积机械内部的锯齿/频率截断（其尾部携带该极限），以及当前的$\ell_1$单元格截断（其尾部平凡小）；审计追踪（存档轮次91--93，条目S1）验证了区分。早期草稿将重数写为$\ell_1$；引理~\ref{l:mult}的除数界是显示指数所使用的。
\end{proof}

\section{局部闭包}\label{s:local}

每一矩阶的局部（模$p$）结构由仿射锁定链控制。固定循环长度$b\ge2$，素数幂$P_0=p^{s}$，单位系数$a_1,\dots,a_b$，自由锁定$j_1,\dots,j_{b-1}\in\mathbb Z/P_0$；定义点$n_{t+1}\equiv a_tn_t+j_t$，最后锁定由循环闭包诱导。写$n_t=A_t(n_1+d_t)$，其中$A_{t+1}=a_tA_t$，$A_1=1$，$d_1=0$；则$d_{t+1}-d_t=j_t/A_{t+1}$是$j_t$的单位倍数。令$N(j)$计数$n_1\in\mathbb Z/P_0$使得所有$b$个点与$p$互素。

\begin{proposition}[重合计数闭包]\label{p:local}
逐点对锁定向量$j$，
\[
N(j)\;=\;P_0\,\frac{p-D(j)}{p},
\qquad
D(j)=\#\{d_t\bmod p:\ t=1,\dots,b\},
\]
且对任何锁定类$\mathcal C$，
\[
\kappa_b(\mathcal C)
=\frac{1-\mathbb E[D\mid\mathcal C]/p}{(1-1/p)^{b}},
\qquad
\mathbb E_j[\kappa_b]=1\ \text{精确。}
\]
禁止类合并是权重$1/(p-1)$的成对事件；三重合并以$O(p^{-2})$进入。
\end{proposition}

\begin{proof}
由于$A_t$是单位，$n_t=A_t(n_1+d_t)$与$p$互素当且仅当$n_1\not\equiv-d_t\pmod p$；全互素条件正好排除$D(j)$个不同剩余$\{-d_t\}$，每个密度$1/p$在$\mathbb Z/P_0$中，得计数。类公式是其对$(1-1/p)^{b}$的平均。均值一恒等式：在均匀锁定下增量独立均匀，故$(d_2,\dots,d_b)$在$(\mathbb Z/p)^{b-1}$上均匀，且$\mathbb E[D]=p(1-(1-1/p)^{b})$，从而$1-\mathbb E[D]/p=(1-1/p)^{b}$。跨越锁定和$\sigma$间隙的合并要求$p\mid\sigma$：对单个单位锁定不可能；对两个单位之和，恰好在第二个变量的一个值上发生，权重$1/(p-1)$。
\end{proof}

机器验证（管道，\texttt{face\_local\_closure}；在认证审计中重新运行，输出相同）：在所有锁定类和$b=4,5,6$，$p\in\{5,7,11,13\}$，$s\in\{1,2\}$的完全枚举和闭式匹配到$2\times10^{-16}$，均值一恒等式精确。四矩表的$\kappa_4$闭式是实例（$\kappa_4(0,0)=1-(4-\tfrac2{p-1})/p$）。该命题提供每一循环长度的奇异级数梳及下面使用的$j$平均分解恒等式。（相应Lean种子的初稿遗漏了$d_t=B_tA_t^{-1}$中的归一化逆；精确枚举代理触发——注册表r103，第29次转换——种子被修正：附录~\ref{a:record}。）

\section{引理D与精确四矩}\label{s:four}

四矩求值$m_4\to13/4$闭合了程序所隔离的估计（“引理D”：平均的、锁定解析的四重关联方差）的消耗截断形式。在其\emph{无限制}形式（模数可达$X$的幂）下，它仍是开放问题，比已发表的Barban--Davenport--Halberstam/色散谱系的对角情形多一个结构性步骤。两个观察从消耗形式中移除了困难：\S\ref{s:trunc}的截断清偿，以及下面的谱框架。

\subsection{框架与Parseval恒等式}

由存档的、机器验证的桥和胶恒等式（\S\ref{s:fm:R}），模数对的$(k,j)$聚合方差由\emph{扩张密度}$\rho(\tau)$承载：令$A_m,A_n$为两个位置格序列的自相关，在周期$N$的分析网格上（$N$为2的幂，超过最大位置的四倍，因此无相关卷绕），
\[
\rho(\tau)\;=\;\sum_{t}A_m(t)\,A_n(t-\tau)
\;=\;\mathrm{irfft}\bigl(|S_m|^{2}|S_n|^{2}\bigr)(\tau),
\]
且锁定方差满足\emph{精确}离散Parseval恒等式
\begin{equation}\label{e:pars}
\sum_{\tau=0}^{N-1}\rho(\tau)^{2}
=\frac1N\Bigl[\mathrm{dens}_0^{2}
+2\!\!\sum_{0<\beta<N/2}\!\!\mathrm{dens}_\beta^{2}
+\mathrm{dens}_{N/2}^{2}\Bigr],
\qquad \mathrm{dens}=|S_m|^{2}|S_n|^{2},
\end{equation}
在管道中逐位验证（相对$10^{-16}$；在认证审计中重新运行）。对$\rho-\mathrm{Model}$应用同一恒等式将方差转化为$\mathrm{dens}$与模型梳的谱$L^2$距离。我们按对$(\gamma_m,\gamma_n)=(b_2\beta,b_1\beta)\bmod1$的弧质量在截断$P^{*}=(\log X)^{B^{*}}$处分类：双主（MM）、混合或双次。

\subsection{三叶}

\begin{lemma}[MM叶]\label{l:MM}
设$\beta$为双主箱：$\gamma_m=a/Q+\eta$，$Q\le P^{*}$，$|\eta|\le$弧宽，$\gamma_n$类似。则对每个$A'$，
\[
S_m(\beta)=\frac{\mu(Q)}{\varphi(Q)}\,I_m(\eta)
+O_{A',B^{*}}\!\bigl(Y(\log Y)^{-A'}\bigr),
\qquad
I_m(\eta)=\sum_{m\in I_m}\e(m\eta),
\]
一致成立，$S_n$同理；常数不可有效化。因此，在MM箱上，
$\bigl|\mathrm{dens}-\mathrm{comb}\bigr|
\ll\mathrm{main}\cdot Y(\log Y)^{-A'}+Y^{2}(\log Y)^{-2A'}$
逐点，其中$\mathrm{comb}$是平方主项之积。
\end{lemma}

\begin{proof}
$S_m(\beta)=\sum_m\LL(m)\e(m\gamma_m)$是在$\gamma_m=a/Q+\eta$处的加窗单素数求和。按模$Q$分拆并应用Siegel--Walfisz定理\cite[\S22]{Dav}（或\cite[推论5.29]{IK}）的形式，结合对相位$\e(m\eta)$的窗口内部分求和，得到显示；常数仅依赖于$(A',B^{*})$，通过Siegel定理不可有效化。逐点梳比较通过展开乘积得到。每箱面（\texttt{face\_sw}）测量六个族上的绝对归一化误差：在$T=2400$处RMS $3.6\times10^{-3}$，在$T=153600$处降至$3.6\times10^{-4}$。
\end{proof}

\begin{remark}[不可有效性在此进入]\label{r:ineffective}
这是链中唯一引入不可有效化常数的点：Siegel定理不给出可有效确定的$T_0(A')$，因此下游每个定理常数都是不可有效化的，而\S\ref{s:faces}中测量的每个有限高度面都带有可计算常数。两类陈述在全文分开分级（见\S\ref{s:faces}的结束注记）；本文其他内容——恒等式、证书、局部律、窗口微积分——均不消耗不可有效化输入。
\end{remark}

\begin{lemma}[次叶与混合叶]\label{l:minor}
对$\gamma$在截断$P^{*}$下为次弧（分母范围$((\log X)^{B^{*}},Y(\log X)^{-B^{*}})$由Dirichlet逼近），$|S(\gamma)|^{2}\ll Y^{2}(\log Y)^{-B^{*}/2+c_0}$由Vaughan估计。因此双次谱质量携带$(\log)^{-B^{*}+2c_0}$节省平方，每个混合箱至少单侧节省；混合和次贡献对谱$L^2$距离为$\ll(\log X)^{-B^{*}/2+c_1}$相对于自然大小。大小记账（像族良好间隔和混合大筛/次上确界分层预算）是经典的，并在面检查中为$24$--$31$倍紧度（\texttt{face\_arcs}）。
\end{lemma}

\begin{lemma}[梳即奇异级数]\label{l:comb}
引理~\ref{l:MM}的谱主梳等于截断奇异级数模型：在$\tau$框架中幸存频率恰好是CRT同余$t\equiv0\ (b_2)$，$t\equiv\tau\ (b_1)$，且局部因子匹配命题~\ref{p:local}的重合计数闭式；奇异级数在$P^{*}$之后的截断尾部为$\ll(\log X)^{-B^{*}+c}$由标准$\sum_{q>P^{*}}\mu^{2}(q)(q,h)/\varphi(q)^{2}$界给出。该识别在附录~\ref{a:comb}中显式展示，并给出了一个模数对的示例；$\tau$分辨模型在16个族和8个高度上以$0.2$--$4\%$的误差跟踪测量轮廓（\texttt{face\_dispersion}）。
\end{lemma}

\subsection{定理}

\begin{theorem}[引理D，截断形式]\label{t:D}
固定$B$，$b_1,b_2\le(\log X)^{B}$为素数幂，窗口长度$Y\asymp X^{\vartheta}$，$\vartheta\in(\tfrac12,1]$。则对每个$A>0$，
\[
\sum_{k\le K}\sum_{|j|\le J}
\bigl|N_4(k,j)-\SSS_4(k,j)V(k,j)\bigr|^{2}
\;\ll_{A,B}\;KJV^{2}(\log X)^{-A}.
\]
因此，结合命题~\ref{p:consume}，$m_4\to13/4$双边。
\end{theorem}

\begin{proof}
取$B^{*}=2(A+c_1+B)$并对$\rho-\mathrm{Model}$应用\eqref{e:pars}。谱$L^2$距离按弧分类分裂。在MM箱上，引理~\ref{l:MM}取$A'=A+2B^{*}$给出逐点相对节省$(\log)^{-A'}$，故$\sum_{\mathrm{MM}}|\mathrm{dens}-\mathrm{comb}|^{2}
\ll(\log)^{-2A'+c}\sum\mathrm{main}^{2}$；由引理~\ref{l:comb}，梳与模型一致到$(\log)^{-B^{*}+c}$相对，被$B^{*}$的选择吸收。混合和双次箱贡献$\ll(\log)^{-B^{*}/2+c_1}$相对由引理~\ref{l:minor}。桥和胶层的$(k,j)$聚合常数已存档且认证（\S\ref{s:fm:R}）；归一化匹配$KJV^{2}$是该处的消费者接口。所有界基于绝对值，故双边；所有常数仅依赖于$(A,B)$且不可有效化。选择$B^{*}$再选择$A'$的顺序不涉及循环。
\end{proof}

\begin{remark}[递推框架路径；证伪者r95]\label{rem:prog}
存档中将方差化简为锁定双模数递推协方差的方法提示了更短论证：将两个递推误差展开为加法特征并用Siegel--Walfisz求值。在模数$(r,q)$上分别执行会产生\emph{不完整}特征和（锁定耦合了模$r$和模$q$的剩余类，对一个模数的$a$和没有对另一个的正交性）；展开必须提升到联合模数$M\asymp rq\le(\log X)^{2B}$，在那里它成为与上面平行的完全和论证。早期草稿呈现了该草图而未提提升；后被谱证明取代，在谱证明中$\tau$变换在整个周期上运行，正交性（引理~\ref{l:comb}）精确。该错误及其修正是向前记录上的证伪者r95。
\end{remark}

\begin{corollary}[$13/18$梯级]\label{c:1318}
使用不可有效化常数，
\[
\liminf_{T\to\infty}\frac{N_0^{s}(T,2T)}{N(T,2T)}\ge\frac{13}{18},
\qquad
\liminf_{T\to\infty}\frac{N_d(T,2T)}{N(T,2T)}\ge\frac{31}{36},
\]
且对任意固定的本原Dirichlet $L$函数同样成立。
\end{corollary}

\begin{proof}
定理~\ref{t:D}和命题~\ref{p:consume}给出$m_4=13/4+o(1)$双边；定理~\ref{t:m3}给出$\mu_3\to2$；证书（引理~\ref{l:cert}，等价于命题~\ref{p:count}在$(1,1,\tfrac43,2,\tfrac{13}4)$处的Christoffel函数计数）精确给出$13/18$和$31/36$，四矩的单调性将$m_4\to13/4$转化为两个下极限界。Dirichlet $L$函数情形同\cite[定理E]{C26}：该框架仅使用函数方程对称性$\rho\mapsto1-\bar\rho$、零点密度界和显式公式，它们对任意固定的本原Dirichlet $L$函数均可用。常数通过引理~\ref{l:MM}不可有效化。
\end{proof}

\section{第五、六矩}\label{s:five}

\subsection{乘积恒等式}

\begin{proposition}\label{p:prod}
对共享一个窗口的$b$循环位置格，任意每格密度$f_i,g_i$：
$\sum_jN_f(j)N_g(j)=\sum_t\prod_iC_i(t)$，
$C_i(t)=\sum_xf_i(x)g_i(x+t)$。特别地，锁定方差是格自相关的逐点乘积，对公共平移求和。
\end{proposition}

\begin{proof}
两个链共享锁定向量当且仅当它们相差一个公共平移；然后链求和分解为位置上的乘积。
\end{proof}

该恒等式（$O(N\log N)$求值）将\S\ref{s:four}的谱论证\emph{逐因子}地传输到每一循环长度：每个因子是加窗单素数求和，因此三叶闭包（对全主箱逐因子Siegel--Walfisz，对任意次因子Vaughan，梳由命题~\ref{p:local}，截断由命题~\ref{p:consume}带更强指数）逐字适用。机器面：孪生奇异级数乘积模型跟踪经验五重积在$0.28$--$2.5\%$（六种模数模式，$T=4800$到$76800$，随高度改善），六重积在$0.11$--$1.29\%$。

\subsection{对层折叠到认证锚点}

Bell$(5)=52$和Bell$(6)=203$类分类账经机器枚举（类表$1/10/15 /10/10/5/1$和$1/15/45/15/20/60/10/15/15/6/1$；五边形弦$10+5$；六边形弦$30+15$；匹配交叉谱$5/6/3/1$）。冻结槽恒等式——单例块重复游走值且不改变重叠范围；机器精确——将每对类化简为四循环几何，其锚点现为\emph{精确有理数}（\S\ref{s:five:cert}）：$t_{\mathrm{adj}}=\tfrac7{60}$，$t_{\mathrm{opp}}=\tfrac1{30}$，$\Phi_4=\tfrac14-2t_{\mathrm{adj}}-t_{\mathrm{opp}}=-\tfrac1{60}$（正弦模型累积量约定，参见注记~\ref{r:books}）。奇块类由局部律消失。因此
\[
m_5=\underbrace{1+\tfrac{10}3+10t_{\mathrm{adj}}
+5t_{\mathrm{opp}}+5\Phi_4}_{=\,67/12\ \text{精确}}+\ C_5,
\qquad
m_6=\tfrac{39}4+6C_5+\{2,2,2\}+\{4,2\}+\{6\},
\]
有理部分无锚（$t$-几何相消：$10t_{\mathrm{adj}}+5\Phi_4+5t_{\mathrm{opp}}=\tfrac54$和$30t_{\mathrm{adj}}+15t_{\mathrm{opp}}+15\Phi_4=\tfrac{15}4$恒等——在认证审计中符号验证并在Lean包中内核编译，\S\ref{s:verif}(f)），其中$\{2,2,2\}=\tfrac{32}{105}$现已\emph{通过二面体放置类精确认证}（命题~\ref{p:t222}；被替代的束$\tfrac{131}{420}$及其修正为注册条目D19），$\{4,2\}=-\tfrac{23}{420}$（带旁观者对的四循环连通对象）和$\{6\}=-\tfrac1{126}$（同样精确认证，\S\ref{s:five:conn}）。先前的写出条目W1--W3已在\S\ref{s:writeouts}中写出。

\subsection{配对层精确认证}\label{s:five:cert}

\begin{proposition}[精确配对常数]\label{p:t222}
二弦锚点是精确有理数
$t_{\mathrm{adj}}=\tfrac7{60}$，$t_{\mathrm{opp}}=\tfrac1{30}$。
六边形的15个三弦匹配分为\emph{五个}二面体放置类，大小分别为$2/3/6/3/1$，交叉数$0/0/1/2/3$，精确值为
\[
T_0=\tfrac3{70}\ (\text{顺序}),\quad
T_0'=\tfrac{17}{420}\ (\text{嵌套}),\quad
T_1=\tfrac1{90},\quad T_2=\tfrac1{180},\quad T_3=\tfrac1{70},
\]
因此
\[
\{2,2,2\}\;=\;2T_0+3T_0'+6T_1+3T_2+T_3\;=\;\tfrac{32}{105}.
\]
交叉数\emph{不是}完全不变量：两个零交叉类具有不同值（$T_0\ne T_0'$），因此本文中每个束均逐放置类求值（注册D19）。
\end{proposition}

\begin{proof}
每个常数是盒上的积分
$(1-\mathrm{spread})_+\cdot\prod\min(|x_i|,1)$对显式游走（Bell分类账的弦模式游走）的积分：一个分段多项式被积函数，其所有扭结超曲面为变量在$\{0,\pm1\}$水平的$\{0,\pm1\}$组合。由于0是游走位置且每个变量是位置或两位置之差，支撑在$[-1,1]^d$内且$\min$-帽从不绑定。然后该积分在\emph{精确}有理算术中通过完全断点枚举的迭代积分计算：内层断点具有单位首项系数，因此所有碰撞和边界结式保持小整数形式，在每个无扭结片上被积函数是次数至多$2/4/8$（在内外变量）的多项式，通过阶$5/7/9$的闭Newton--Cotes公式在有理节点处精确积分。每一片被积分两次（整片和两半）：对次数低于规则阶的多项式，两个精确有理数必须一致；任何遗漏的扭结会在包含片上导致不一致；未发生不一致。$t_{\mathrm{opp}}$和$T_0$另外允许通过象限分解的短闭式求值（附录~\ref{a:cert}），与算法值一致；每个类值由第二种独立精确方法（多面体分解，与迭代积分无算术共享）重现，且轨道不变性通过求值每个二面体类的额外成员进行机器检查。代码和输出：
\texttt{certification/exact\_t222.py}及复现包的\texttt{phase0}审计。
\end{proof}

\begin{remark}[分类修正；注册D19]\label{r:d19}
本研究早期阶段使用了$\{2,2,2\}=5T_0+6T_1+3T_2+T_3=\tfrac{131}{420}$，按交叉数对15个匹配分类，每交叉类求值一个代表；嵌套零交叉几何从未单独求值，存档梯值$0.31170(60)$是由同样的四个代表组装——循环确认。该错误被八矩配对运行的预注册校准门捕获（第38次转换），嵌套类被精确闭包（$T_0'=\tfrac{17}{420}$，三条独立路径），且两个在旧值下偏离$55$--$110\sigma$和$11$--$21\sigma$的模型侧估计量现在读$\sim1\sigma$相对校正值，被回溯记录为独立检测。该修正\emph{降低}$M_6$并\emph{提升}定理：先前发表的界（$0.7962/0.8981$）消耗了一个较弱但为真的上界，仍然有效。
\end{remark}

\begin{remark}[锚点修正]\label{r:anchor}
存档锚点$t_{\mathrm{adj}}=0.11717$（轮次69--71网格）由命题~\ref{p:t222}修正为$\tfrac7{60}=0.11\overline{6}$：在扭结处的$+0.43\%$矩形法则偏差。下游无变化：$m_5$和$m_6$有理层无锚（上面），注记~\ref{r:books}的对账变为精确。先前猜测的三弦类有理识别现为定理。注册条目D7、D8。
\end{remark}

\subsection{写出条目W1--W3}\label{s:writeouts}

审计分类账中作为命名写出条目的三项在此处理；其预注册面在\texttt{certification/writeouts\_verification.py}中（全部通过；一条面在构造期间触发并被转换——附录~\ref{a:record}）。

\begin{lemma}[W2：自由槽分解]\label{l:W2}
在任意$b$循环单元格分类账中，若槽的锁定无约束则精确分解：单元格聚合等于收缩的$(b-1)$循环聚合乘以自由槽的窗口质量$\sum_{m\in W}\LL(m)$，且$\sum_{m\in W}\LL(m)=|W|\bigl(1+O(e^{-c\sqrt{\log X}})\bigr)$。
\end{lemma}

\begin{proof}
对自由槽偏移的锁定和是无限制的，因此双重和通过Fubini分解；在游走几何侧这是冻结槽恒等式（单例块重复游走值且不改变重叠范围——机器精确，\texttt{face\_o5\_gate}/\texttt{face\_b6\_frame}）。质量求值是素数定理带经典误差项\cite[\S18]{Dav}；消耗归一化将其除去。面验证了实族上的精确Fubini恒等式以及高度处的PNT质量（$X=2\times10^{5}$处$1.00041$）。
\end{proof}

\begin{lemma}[W3：三对分解]\label{l:W3}
六矩分类账的$\{2,2,2\}$类聚合等于三个对主项乘积乘以联合重叠体积，在消耗归一化中偏差聚合$\ll_A KJ'V'^2(\log X)^{-A}$。
\end{lemma}

\begin{proof}
六个位置形成三条不相交仿射对链。$p$局部密度通过CRT和命题~\ref{p:local}在三链上分解（不相交槽；合并修正是已计入局部律的成对事件），故奇异级数模型是三条孪生级数的乘积。阿基米德耦合仅通过联合窗口重叠（游走几何$O_6$）进入，即命题~\ref{p:t222}的已认证精确层。与对主项乘积的偏差由三次应用$k=2$实例的谱传输控制：由命题~\ref{p:prod}，解析方差在三链上分解，每个因子是加窗单素数求和，且\S\ref{s:four}的三叶闭包对每个因子适用，指数为命题~\ref{p:consume}。面在$X=6\times10^{4}$处以中位数$1.9\%$（四种平移模式）跟踪乘积模型。
\end{proof}

\begin{proposition}[W1：解析多维梳]\label{p:W1}
对具有$b-1$个锁定坐标的$b$循环，解析扩张密度
$\rho(\tau_1,\dots,\tau_{b-1})$满足多维梳识别：模型梳支撑在乘积CRT同余上（每坐标一个同余，如引理~\ref{l:comb}），局部因子为命题~\ref{p:local}的$\kappa_b$值及模数上的孪生级数因子；且解析谱$L^2$距离与梳的界满足定理~\ref{t:D}，指数相同。
\end{proposition}

\begin{proof}
由命题~\ref{p:prod}，解析方差由联合分析网格上的乘积$\prod_i|S_i|^2$承载，且多$\tau$变换在每个坐标上完全，因此$(b-1)$维Parseval恒等式逐坐标成立（在面中精确验证，$1.2\times10^{-16}$）。\S\ref{s:four}的频率分类逐因子适用：在全主箱上每个因子满足引理~\ref{l:MM}；任何次因子由引理~\ref{l:minor}闭包；每坐标的梳识别是引理~\ref{l:comb}，其CRT同余现在每坐标运行（每坐标完全正交，因为每个变换覆盖其全周期）。定理~\ref{t:D}的论证然后逐字重复，带有$b-1$个分类指标；截断指数随$k$增强（命题~\ref{p:consume}）。先前所述的括号（上面聚合面，下面解析四循环情形）被包含：聚合实例是$\tau$求和情形，解析四循环情形是引理~\ref{l:comb}本身。解析梳面（实际$5$循环族）直接显示CRT类对比。
\end{proof}

\subsection{连通值}\label{s:five:conn}

$C_5=\int O_5C_5^{\mathrm{Urs}}=\tfrac1{36}$，
$\{4,2\}=-\tfrac{23}{420}$和$\{6\}=-\tfrac1{126}$，
\emph{在精确有理算术中证明}并在机器精度下独立确认。

\emph{符号认证。}直接路径——命题~\ref{p:t222}的分段精确积分从三维推广到四维和五维——逐项闭包$\{4,2\}$，但不能闭包$\{6\}$：片数在积分层次上倍增，单个五维项投影到数百年精确算术。闭包通过表示变换实现。由于
\[
\mathrm{ov}(P)\;=\;\bigl(1-(\max_i p_i-\min_i p_i)\bigr)_+
\;=\;\mathrm{Leb}\{\,t\in\mathbb R:\ t\le p_i\le t+1\ \forall
i\,\},
\]
\S\ref{s:conv}带符号和的每一项\emph{就是}一个有界有理多面体的体积，
\[
\int_{\mathbb R^d}\mathrm{ov}_A(x)\,\mathrm{ov}_B(x)\,dx
\;=\;\mathrm{Vol}_{d+2}\bigl\{(x,t,s):\ t\le A_i(x)\le t+1,\
s\le B_j(x)\le s+1\bigr\},
\]
每个面系数在$\{0,\pm1\}$中（有界性因为0在两个位置族中，迫使$t,s\in[-1,0]$；$\{4,2\}$变体的权重$\min(|v|,1)$再提升一维）；代价从每级片数的\emph{乘积}转移到面格的大小。每个体积通过精确面递归
$\mathrm{Vol}_n(P)=\tfrac1n\sum_i\mathrm{dist}(c,H_i)\,
\mathrm{Vol}_{n-1}(F_i)$求值，安排为将每个面投影到坐标超平面以抵消无理面归一化——计算全程保持在有理算术中——并在紧约束集上记忆化。所有$1310$项（$C_5$的$150$项，$\{4,2\}$的$78$项，$\{6\}$的$1082$项）作为精确分数求值，每项存档在复现包中，其带符号和为
\[
C_5=\tfrac1{36},\qquad \{4,2\}=-\tfrac{23}{420},\qquad
\{6\}=-\tfrac1{126}
\]
精确。门全部通过：经同一代码路径的纯四循环返回$\Phi_4=-\tfrac1{60}$（两种精确方法）；三维回归重现命题~\ref{p:t222}的每个精确值（$\tfrac3{70},\tfrac1{90},\tfrac1{180},\tfrac1{70}, \tfrac7{60},\tfrac1{30}$）；两种符号方法在双方计算的每项上逐分式一致；项构造（位置族、掩码、符号、权重）经数值交叉检查与梯引擎一致；$\{4,2\}$的三个放置类结果为$-\tfrac2{315}$，$-\tfrac1{1260}$，$-\tfrac1{252}$，以重数$6/6/3$重组为$-\tfrac{23}{420}$。

\emph{数值确认。}七阶分组中点梯（$C_5$的$dv$降至$0.00125$，$\{6\}$降至$0.0125$，$\{4,2\}$旁观者对变体降至$0.0015625$；GPU加速引擎，支持剪枝、项trie和否定对称性，在数学上与参考引擎相同并在三种架构上逐位交叉验证）收敛，相邻行Romberg漂移$2\times10^{-16}$，收敛到已证明分数；在误差带内有理重构\emph{由$Q^{2}\ge1/(2\varepsilon)$准则唯一}，落在相同值上。可信度锚点：相同管道在相同运行中重现$b=3$的精确值$0$和$b=4$的$\Phi_4=-\tfrac1{60}$至$10^{-18}$；原始梯误差比相对已认证分数收敛到纯$h^{2}$值$4.000$；且$\{4,2\}$\emph{命中其预注册候选}（$-\tfrac{23}{420}$，在精细梯运行前从三个粗略参考梯注册），逐分量。七循环常数以少一阶识别为$C_7=-\tfrac{17}{360}$（仅识别等级；其符号闭包被委托，并在$k=7,8$层以所述等级消耗，\S\ref{s:seven}）。混合奇块类$\{3,3\}$\emph{消失}：恒为零——在测试的每个网格（两种协议，四种分辨率，所有三个放置类）上为精确浮点零，而非小数字；算术侧奇块消失的模型面在\S\ref{s:five}中证明；它先前在$m_6$预算中携带的剩余允量已退役（注册D17）。一个观察但不消耗：通过等级$6$认证和识别的常数——$\tfrac13$，$\tfrac7{60}$，$\tfrac1{30}$，$-\tfrac1{60}$，$\tfrac{32}{105}$，$\tfrac1{36}$，$-\tfrac{23}{420}$，$-\tfrac1{126}$，$\{4,2\}$分量和五个六边形弦类——分母整除$1260$，加上$C_7=-\tfrac{17}{360}$则全部分母整除$2520=\mathrm{lcm}(1..9)$；但该律在等级$8$\emph{失效}：\S\ref{s:seven}的$\{2^4\}$轨道值分母高达$90720=2^5\cdot3^4\cdot5\cdot7$，且束$\tfrac{1661}{3780}$有$3780\nmid2520$——这正是小分母重构在那找不到候选的原因（注册D21）。
（早期端点网格求值已退役——注册条目，第32次转换。）
涨落由逐因子方差引理控制（定理~\ref{t:D}的传输；解析多维梳的展示由聚合面和解析四循环情形界定——现为命题~\ref{p:W1}）。端到端面：组装
$m_5=\tfrac{101}{18}=5.6\overline{1}$和$m_6=\tfrac{640}{63}=10.15873\ldots$均位于正弦模型带\eqref{e:guebands}内。

\emph{等级说明。}由\S\ref{s:conv}的维数定理，每个连通常数都是有理数；上面的符号认证精确计算了三个已用常数，它们以\emph{证明}等级进入\S\ref{s:lp}——精确有理计算机辅助，同命题~\ref{p:t222}的等级（注册D18）。先前的识别条件已解除；数值识别（机器精度Romberg极限，唯一性保证重构，双精确锚定校准在同一管道中，逐分量$\{4,2\}$预注册命中）作为独立确认，与符号值逐位一致。精确等级见\S\ref{s:verif}(d)。

\subsection{卷积演算、有理性和中点协议}\label{s:conv}

连通常数具有第二个结构性表示，既固定其算术性质又指示数值协议。记$W$为单位方窗自相关，$B$为正弦核盒；累积量被积函数由这两个轮廓的迭代卷积组装。三个结果（推导和验证门在复现包中，\texttt{cyclic\_cumulant}和\texttt{midpoint\_ladder}引擎）：

\emph{(i) 迹/卷积恒等式。}在$b=2$时连通积分以闭式求值，$\int O_2C_2=1-\|W\star\widehat B\|^{2}=\tfrac13$精确——这是下面所有内容的校准门（F-TRACE）；$b=3,4$实例重现命题~\ref{p:t222}的精确配对值至$3\times10^{-14}$。

\emph{(ii) 划分-循环累积量公式（F-CYC）。}对$b$循环连通类，
\[
C_b(v)\;=\;\sum_{P\vdash[b]}(-1)^{|P|-1}
\sum_{\text{cyclic orders of }P}
\operatorname{ov}\bigl(\text{merged-}w\text{ walk}\bigr),
\]
内和对块的不同循环排列进行，$\operatorname{ov}$为合并游走的重叠体积——在$b=3,4,5$处与直接Ursell被积函数验证至$3\times10^{-14}$。

\emph{(iii) 维数与有理性。}在(ii)的每个$(P,\sigma)$项中，游走约束消去一个积分变量（$(m-1)$维），所得区域是附录~\ref{a:cert}(A1)的单位系数超平面族$\mathcal H$切割的有限多面体并，带有有理系数多项式被积函数。因此每一项都是有理数，$C_5$、$\{4,2\}$和$\{6\}$亦然——产生$t_{\mathrm{adj}}=\tfrac7{60}$和$\{2,2,2\}=\tfrac{32}{105}$闭式的同一论证。剩余常数与命题~\ref{p:t222}的区别仅在于精确积分的片数，而非种类；\S\ref{s:five:conn}的多面体路径驯服了该片数并完成了所有三个常数的精确积分，关闭了条目N1（\S\ref{s:verif}）。

\emph{(iv) 长球迹恒等式。}令$A$表示限制在$[0,1]$上的正弦核积分算子$K(x,y)=\mathrm{sinc}(\pi(x-y))$（单位密度的经典长球算子）。则对每个$m\ge1$，
\[
\operatorname{tr}A^{m}
\;=\;\int_{\mathbb R^{m-1}}
\bigl(1-\mathrm{span}(0,p_1,\dots,p_{m-1})\bigr)_+
\prod_{i=1}^{m}\mathrm{sinc}(\pi v_i)\,dv ,
\]
右端正是本文重叠加权的纯$S$循环多面体积分。\emph{证明}：写$\operatorname{tr}A^m=\int_{[0,1]^m}\prod K(x_i,x_{i+1})\,dx$，代入增量$v_i=x_{i+1}-x_i$，注意$\int_{\mathbb R}\prod_k\mathbf 1[p_k\in[t,t+1]]\,dt =(1-\mathrm{span})_+$——重叠权重\emph{就是}硬窗口的平移平均。$\square$（在$m=3$处数值验证至所有打印数字。）因此整个纯矩循环层在每一阶都由单一经典谱承载；连通常数仅通过重合（Ursell）装饰与之不同，专用探测显示该装饰\emph{不是}该谱的谱多项式泛函。

\emph{(v) 非交叉闭式演算。}对于弦图\emph{无交叉}（允许嵌套）的配对类，类值坍缩为一维卷积代数：令$g(t)=\int_t^{t+1}|x|\,dx=t^2+t+\tfrac12$在$[-1,0]$上，则顺序$\tfrac b2$弦轨道等于$\int_{-1}^{0}g(t)^{b/2}dt$，嵌套弦作为一维迭代卷积复合。四个精确确认：$b=6$顺序$=\tfrac3{70}$和嵌套$=\tfrac{17}{420}$（命题~\ref{p:t222}的两个零交叉类）；$b=8$顺序$=\tfrac{83}{5040}$（存档$\{2^4\}$轨道）；以及$b=10$顺序值$\tfrac{73}{11088}$，\emph{在查询存档前由闭式预测}并精确找到（最大的$\{2^5\}$轨道）。伴随的\emph{跨度分解}将初等范围扩展到低宽度的交叉情形：将每个卦限的重叠因子化为显式分段线性跨度，将所有五个$b=6$类（包括交叉值$\tfrac1{90}$，$\tfrac1{180}$，$\tfrac1{70}$）化为Dirichlet型多面体积分，重新推导了命题~\ref{p:t222}的校正束$\{2,2,2\}=\tfrac{32}{105}$的闭式：这是D19常数的第三条完全独立路径。该演算还确定了分母现象的来源：$\int\mathrm{sinc}^{2m}$及其相关量是中心B样条值，其阶乘族分母产生全文观察到的随等级变化的$\mathrm{lcm}$模式（\S\ref{s:outlook}(e)）。一般（交叉）情形的内在复杂度参数是图的\emph{交叉宽度}——非交叉扇区坍缩为闭式，而完全交叉轨道保持环境维数；这在\S\ref{s:outlook}(c)中量化。

\emph{(vi) 对数行列式恒等式：原点无原子。}塔的测量来自精确有限模型：Haar随机$N\times N$酉矩阵的特征角$\theta_j$处Fourier系统$W_{jm}=e^{im\theta_j}$（$m=0,\dots,N{-}1$）的Gram矩阵$\widehat G_N=WW^{*}/N$，其期望谱矩重现正弦模型塔。其对数的行列式在每一个有限$N$处精确可计算：$\det\widehat G_N=|\Delta(e^{i\theta})|^{2}/N^{N}$，$\Delta$为Vandermonde，Morris积分
$\mathbb E\,|\Delta|^{2s}
=\Gamma(1+(1+s)N)\big/\bigl(\Gamma(2+s)^{N}\Gamma(1+N)\bigr)$
给出，在$s=0$处微分，
\[
\mathbb E\Bigl[\tfrac1N\textstyle\sum_i\log\lambda_i\Bigr]
= H_N-1-\log N,
\qquad
\operatorname{Var}\Bigl[\tfrac1N\textstyle\sum_i\log\lambda_i\Bigr]
= \psi'(N{+}1)-\tfrac{\psi'(2)}N ,
\]
对每个$N$精确（$N=1$时两侧为零）。
因此极限谱测度$\nu$满足
\[
\int\log x\,d\nu\;=\;\lim_{N}\,(H_N-1-\log N)\;=\;\gamma-1
\]
——Euler常数在程序中的第二次精确出现（连续求积常数N2为$\gamma-2$）——并且定量地，\emph{原点无原子}：由于$\log^{+}x\le(x-1)_{+}$，且由Cauchy--Schwarz对精确第二中心矩$\Sigma_2=\tfrac13$（\eqref{e:sigmas}）有$\tfrac1N\sum_i(\lambda_i-1)_{+}\le\sqrt{\Sigma_2}$，负部满足$\tfrac1N\sum_i\log^{-}\lambda_i\le \tfrac1{\sqrt3}+1-\gamma+o(1)$，故由Markov
\[
\nu([0,\varepsilon])\;\le\;
\frac{1/\sqrt3+1-\gamma}{\log(1/\varepsilon)}
\qquad\Longrightarrow\qquad \nu(\{0\})=0 .
\]
两个恒等式作为程序每个谱池的常设接受门（门F-LOGDET）：在三个存档池（$N=128,256,512$；$800/250/60$样本）上，均值恒等式在$-0.08/+0.24/-0.48$标准误差内，方差恒等式在峰度校正误差棒内为$-1.05/+1.53/+0.10$。事实上同一Morris积分给出\emph{整个}累积量生成函数在每个有限$N$处的精确表达式：对$s>-1$，
\[
\log\mathbb E\exp\Bigl(s\,{\textstyle\sum_i}\log\lambda_i\Bigr)
=\log\Gamma(1+(1+s)N)-N\log\Gamma(2+s)-\log\Gamma(1+N)
-sN\log N,
\]
因此每个累积量闭式，
$\kappa_j(\sum_i\log\lambda_i)
=N^{j}\psi^{(j-1)}(N+1)-N\,\psi^{(j-1)}(2)$，且极限自由能
$\Lambda(s)=\lim_N\frac1N\log\mathbb E\,e^{s\sum\log\lambda_i}
=(1+s)\log(1+s)-s-\log\Gamma(2+s)$。所有$j\ge2$阶累积量都是广延的，因此对数行列式服从中心极限定理，具有精确方差常数$\Lambda''(0)=2-\pi^{2}/6$并以指数浓度集中——无原子陈述的浓度形式（门F-CGF：存档池上$s=\pm0.1$处的倾斜均值和三阶累积量在一个标准误差内重现闭式）。其对程序视野的推论见\S\ref{s:outlook}(f)。

\emph{(vii) Toeplitz--迹表示：精确有限$N$矩。}由于$WW^{*}$和$W^{*}W$共享谱，$\widehat G_N$（(vi)）与随机Toeplitz矩阵$G'_{mn}=t_{m-n}$等谱，其中$t_k=\frac1N\operatorname{tr}U^{k}$为CUE迹。展开$\operatorname{tr}(G')^{b}$对指标游走求和，在每个有限$N$处精确得到
\[
m_b(N)
=\frac1{N^{b+1}}\!\!\sum_{k_1+\dots+k_b=0}\!\!
\bigl(N-\mathrm{span}(0,S_1,\dots,S_{b-1})\bigr)_{+}\,
N^{\#\{k_i=0\}}\;
\mathbb E\!\!\prod_{k_i\neq0}\!\operatorname{tr}U^{k_i},
\]
$S_j$为部分和——这是贯穿本文的重叠权重$(1-\mathrm{span})_+$的离散孪生——且迹矩在每个$N$处通过Schur正交性精确，
$\mathbb E\,[p_\mu\overline{p_\nu}]
=\sum_{\lambda\vdash K,\,\ell(\lambda)\le N}
\chi^{\lambda}(\mu)\chi^{\lambda}(\nu)$，Diaconis--Shahshahani范围即无长列扇区。两个推论有两行证明。\emph{消失引理}：非零权重迫使$\mathrm{span}<N$，故逐阶$|k_i|<N$且对每个$(+\cdots+,-)$模式有$K<N$，将这些配置置于正交扇区，混合圈型积分为零。\emph{每个有限$N$处的闭式}：
\[
m_2(N)=\tfrac43-\tfrac1{3N^{2}},\qquad
m_3(N)=2-\tfrac1{N^{2}},\qquad
\Sigma_3(N)\equiv0,
\]
由$\mathbb E|\operatorname{tr}U^{k}|^{2}=\min(k,N)$和初等格点求和得到。完全和（有理算术；两个独立实现，直接求和与多重集聚合动态规划，逐位匹配）的精确求值给出$m_b(N)$到$b=10$，且值组织为$1/N^{2}$的多项式，次数$\lfloor b/2\rfloor$，其常数项\emph{重新推导，独立于本文所有位置空间方法}，塔值$\tfrac{13}4$，$\tfrac{101}{18}$，$\tfrac{640}{63}$，$\tfrac{3439}{180}$，$\tfrac{747361}{20160}$——特别是D19校正六矩的第四条独立路径，拟合是未知的并显式排除了被替代的$\tfrac{4297}{420}$——每个识别通过在保留点$N$值上的逐位确认（例如$m_9(7)$的十二位分子有理数和保留点$m_{10}(8)$，均在计算前预测）。$1/N^{2}$系数是精确有限尺寸律，例如$\Sigma_2(N)=\tfrac13-\tfrac1{3N^{2}}$到$\Sigma_8(N)=\tfrac{633}{2240}-\tfrac{3371}{2016}N^{-2}
+\tfrac{12247}{2880}N^{-4}-\tfrac{5545}{1008}N^{-6}
+\tfrac{277}{105}N^{-8}$，并在$N=128$和$N=192$处$10^{7}$样本校准运行中重现观察到的系统偏差，覆盖所有六个已知中心矩（十二次比较，均在统计误差内）：测量偏差转化为定理。$m_b(N)N^{b+1}$是否对每个$b$都是$N$的多项式（观察到$b=10$；八个保留点确认，加上$b=9$处顶部系数消失确认次数律）是首次观察到时表示中唯一开放的结构问题；\S\ref{s:tt}现在将其完全证明（定理~\ref{t:poly}）：答案是肯定的，对每个$b$和每个$N\ge1$，具有奇偶性$L(-N)=(-1)^{b+1}L(N)$（定理~\ref{t:parity}）。扇区隔离实验已将其形状尖锐化：将和写为Wick扇区加特征修正扇区（Schur正交性），每个扇区从$b=5$起分别是一个\emph{周期二}准多项式（偶和奇分支在$b=5$处通过$N=17$的保留点验证），且两扇区的奇奇偶性差异在总和上精确相消——这一相消现在由\S\ref{s:tt}的多面体表示解释：扇区分裂是特征基的伪像，组装对象根本没有剩余类。

\emph{中点协议。}(i)--(iii)的数值梯揭示了一个网格病理，记录为注册表第32次转换（附录~\ref{a:record}，D10）：在扭结被积函数上端点网格在$b\ge6$且粗步长时符号翻转，而中点网格——从不在扭结集$\mathcal H$上采样——单调稳定。本文所有连通常数采用分组中点网格梯求值，在$b=4$处针对命题~\ref{p:t222}的精确值校准（在梯自身带内一致），早期端点网格采用已退役。

\section{第七、八矩}\label{s:seven}

\subsection{分类账与中心矩折叠}\label{s:seven:ledger}

Bell$(7)=877$和Bell$(8)=4140$类分类账由两个独立实现机器枚举两次（一个直接规范化二面体模式，一个在$D_k$下进行轨道分解，带有可除性检查$|\mathrm{orbit}|\,\big|\,|D_k|$——该检查在计算任何值之前捕获了一个枚举器的去重错误；注册D20）。每个含$3$块的类消失（\S\ref{s:five}的奇块消失，逐类验证）；每个含单例的类通过冻结槽恒等式折叠到较低锚点，并被\S\ref{s:recentre}重新中心恒等式\eqref{e:recentre}的二项式系数逐阶吸收。幸存者恰为无单例层：
\[
\Sigma_7=\{5,2\}+C_7,
\qquad
\Sigma_8=\{2^4\}+\{2,2,4\}+\{4,4\}+\{6,2\}+C_8,
\]
放置类清单（全部二面体，两次机器审计）：$\{5,2\}$：$21$个划分为$3$个大小为$7$的轨道；$\{2^4\}$：$105$个匹配为$17$个轨道；$\{2,2,4\}$：$210$个为$22$个；$\{4,4\}$：$35$个为$7$个；$\{6,2\}$：$28$个为$4$个（大小$8/8/8/4$）。有理层是闭式\eqref{e:rat78}，分类账组装精确重现它们（门F-PRED-RAT：两个独立推导，\S\ref{s:recentre}）。

\subsection{联合常数，精确认证}\label{s:seven:conn}

\begin{proposition}[等级$7$和$8$的精确联合常数]\label{p:joints78}
逐放置类在精确有理算术中求值，通过面递归多面体积分器（第二种方法：迭代Newton--Cotes断点积分；两者无算术共享）：
\[
\{5,2\}=7\Bigl(\tfrac5{504}+\tfrac1{360}+\tfrac{13}{2520}\Bigr)
=\tfrac18,
\qquad
\{2^4\}=\tfrac{1661}{3780},
\]
\[
\{2,2,4\}=-\tfrac{127}{840},\qquad
\{4,4\}=\tfrac{23}{4536},\qquad
\{6,2\}=-\tfrac{563}{11340},
\]
$\{2^4\}$值为$17$个精确轨道值的加权和（每轨道存档；分母高达$90720$）。纯连通常数，通过同一积分器在轨道约化下（$685$个$D_7$轨道，$9366$项；$6027$个$D_8$轨道，$94586$项；带符号轨道权重记账经$b=5,\dots,8$与原始项和验证）：
\[
C_7=-\tfrac{17}{360},
\qquad
C_8=\tfrac{157}{4032}.
\]
因此，所有输入现经符号认证，
\[
\Sigma_7=\tfrac18-\tfrac{17}{360}=\tfrac7{90},
\qquad
\Sigma_8=\tfrac{307}{1260}+\tfrac{157}{4032}=\tfrac{633}{2240},
\]
并由\eqref{e:rat78}，
$m_7=\tfrac{685}{36}+\tfrac7{90}=\tfrac{3439}{180}$，
$m_8=\tfrac{217}6+\tfrac{28}{45}+\tfrac{633}{2240}
=\tfrac{747361}{20160}$。
\end{proposition}

\emph{验证谱系。}$\{5,2\}=\tfrac18$ \emph{命中其预注册候选}（在符号运行前从模型侧测量$0.124944$注册；符号结果与候选偏差：零），且此后每个轨道获得了完整的GPU梯（$3/3$轨道在$1.1\times10^{-6}$内）——这是故意冗余，因为$C_8$事件表明候选命中本身不是保证。$\{2^4\}$轨道值与符号运行前计算的独立三阶中点梯一致，$17/17$类（最大偏差$1.3\times10^{-7}$来自$h^4$外推；预注册梯带乐观了$3$倍——注册D21）。$\{2,2,4\}$，$\{4,4\}$和$\{6,2\}$目前仅依赖面递归路径（验收套件$22/22$相对每个先前认证常数，轨道不变性检查通过）；其独立第二路径GPU梯已返回并通过（F-V-JOINTS，$33/33$轨道在$3\times10^{-6}$内；\S\ref{s:verif}(d)）。纯常数的委托符号运行已\emph{返回}：$C_7$的运行（$685$个轨道，$208$秒在$24$核上）精确落在识别值$-\tfrac{17}{360}$上——一个早两版注册的预测，现为第三次预注册命中——而$C_8$的运行（$6027$个轨道，$3811$秒在$200$进程上）返回$\tfrac{157}{4032}$，\emph{偏离}预注册候选$\tfrac7{180}$达$4.96\times10^{-5}$：在容差门（$3\times10^{-3}$）内，因此角点移动，证书重新优化（\S\ref{s:lp}）；候选偏差前向记录（D22），分母$4032=2^6\cdot3^2\cdot7$重复了等级$8$对$2520$律的打破。返回值的模型侧交叉检查：$\Sigma_8=\tfrac{633}{2240}=0.28259$相对测量$0.28249$（偏差$9.9\times10^{-5}$）；$m_7$，$m_8$相对CUE模型为$1.6\times10^{-4}$，$1.7\times10^{-3}$。

\subsection{传输元定理W$_\infty$}\label{s:seven:transports}

第七、八矩的算术侧使用与\S\S\ref{s:trunc}--\ref{s:five}相同的四个$k$一致工具——乘积恒等式（命题~\ref{p:prod}，任意$b$），重合计数局部律（命题~\ref{p:local}，任意循环长度），截断清偿（命题~\ref{p:consume}，指数$P^{-(k-2)}$随$k$ \emph{增强}），以及\S\ref{s:four}的三叶谱闭包。我们不逐一写出本文四个类的类比W1--W3，而是\emph{一次性地}写出传输，对块结构一致；每一矩阶的每个联合类都是实例，且每增加一阶的分析边际成本为零。

\begin{theorem}[W$_\infty$：无单例类的逐因子传输]\label{t:winfty}
设$b\ge2$且$\mathcal B=\{B_1,\dots,B_r\}$是$b$循环槽的划分，每块大小$\ge2$且无大小为$3$。将每块视为$b$循环单元格分类账中的仿射锁定链（块内锁定自由；块间仅通过共享窗口相互作用）。则类聚合等于链主项乘积乘以联合重叠体积$O_b$，在消耗归一化中偏差聚合
\[
\Bigl|\,\mathrm{Agg}(\mathcal B)\;-\;\prod_{i=1}^{r}
\mathrm{Main}(B_i)\cdot O_b\,\Bigr|
\;\ll_{A,B}\;KJ'V'^2(\log X)^{-A}
\qquad\text{对每个 }A>0 .
\]
\end{theorem}

\begin{proof}
\emph{(i) 局部分解。}块的位置格不相交，因此模$p$密度通过CRT跨块分解；在每块内它是重合计数闭式（命题~\ref{p:local}，任意链长），跨块合并修正是已计入局部律的成对事件（权重$1/(p-1)$；三重合并$O(p^{-2})$）。奇异级数模型因此是每链梳的乘积。含$3$块的类恒消失（\S\ref{s:five}的奇块消失及其模型面），因此排除不失一般性。

\emph{(ii) 阿基米德耦合仅通过$O_b$。}块间唯一的阿基米德相互作用是联合窗口重叠——游走几何$O_b$——即\S\ref{s:five:cert}，\S\ref{s:seven:conn}的已认证精确层。

\emph{(iii) 逐因子方差传输。}由命题~\ref{p:prod}（对任意每格密度成立），解析方差是链自相关乘积对公共平移求和；多$\tau$变换的每个坐标在其全周期上运行，因此$(r{-}1)$维Parseval恒等式逐坐标成立，如命题~\ref{p:W1}。频率分类逐因子进行：在全主箱上每个因子是满足引理~\ref{l:MM}的加窗单素数求和；任何次因子由引理~\ref{l:minor}闭包；每坐标的梳识别是引理~\ref{l:comb}，每坐标的CRT同余具有完全正交性。

\emph{(iv) 截断与常数。}$\ell_1$截断由命题~\ref{p:consume}清偿，指数$P^{-(b-2)+\varepsilon}$随$b$增强；$(k,j)$聚合常数的桥和胶层以及$KJ'V'^2$归一化接口是\S\ref{s:fm:R}的存档认证层，与块结构无关。所有界对绝对值成立，故双边；所有常数依赖于$(A,B)$且通过Siegel定理不可有效化。无步骤依赖于$b$或$\mathcal B$，除显示参数外。
\end{proof}

\begin{corollary}[实例]\label{c:winst}
W3（引理~\ref{l:W3}）是$\mathcal B=\{2,2,2\}$情形；本文消耗的类是进一步实例\emph{W4} $=\{5,2\}$，\emph{W5} $=\{2^4\}$，\emph{W6} $=\{2,2,4\},\{6,2\}$，\emph{W7} $=\{4,4\}$；且第九、十阶分类账（$\{7,2\}$，$\{5,4\}$，$\{5,2,2\}$；$\{8,2\}$，$\{6,4\}$，$\{5,5\}$，$\{6,2,2\}$，$\{4,4,2\}$，$\{4,2,2,2\}$，$\{2^5\}$）无需进一步写出。纯类$\mathcal B=\{b\}$不通过定理~\ref{t:winfty}路由：其模型值直接输入，其涨落由逐因子方差引理控制，如$k=5,6$（\S\ref{s:five:conn}）。
\end{corollary}

\begin{theorem}[第七、八矩]\label{t:m78}
在端点极限下，以命题~\ref{p:joints78}所述等级，
\[
m_7\;\to\;\tfrac{685}{36}+\Sigma_7\;=\;\tfrac{3439}{180},
\qquad
m_8\;\to\;\tfrac{217}6+8\Sigma_7+\Sigma_8\;=\;\tfrac{747361}{20160},
\]
双边，等级为\S\ref{s:verif}层(b)：有理层由重新中心恒等式和冻结槽折叠（均为精确），无单例类由定理~\ref{t:winfty}（实例W4--W7，推论~\ref{c:winst}）带已认证联合常数，涨落控制由定理~\ref{t:D}的逐因子方差传输带命题~\ref{p:consume}在$k=7,8$的截断指数。
\end{theorem}

\begin{proof}
$k=7,8$的迹展开运行于\S\ref{s:seven:ledger}的Bell分类账。含$3$块的类由奇块消失消失（\S\ref{s:five}）；含单例的类由W2（引理~\ref{l:W2}）精确分解并折叠到较低锚点——这些贡献的总和是二项式有理层，等于$\tfrac{685}{36}$和$\tfrac{217}6+8\Sigma_7$-修正，由\eqref{e:recentre}且分类账组装精确重现闭式（F-PRED-RAT）。无单例类是推论~\ref{c:winst}的W$_\infty$实例和纯类$C_7$，$C_8$；其模型值为命题~\ref{p:joints78}的认证常数，其偏差逐因子传输如同$k=5,6$层（\S\ref{s:five:conn}，“涨落”），带更强截断指数。尾部和阈值步骤$k$无关（命题~\ref{p:tail}）。交叉检查：组装$m_7=19.1055\ldots$和$m_8=37.0714\ldots$位于模型测量带（$19.09\pm0.09$，$37.1\pm0.3$）内，并高于/在固定$m_1,\dots,m_6$上的校正矩锥顶点（$19.0712$，$36.7622$）内。
\end{proof}

\section{第九至第十四矩：Toeplitz--迹路径}\label{s:tt}

$k=10,12,14$证书所需的六个进一步中心矩由\S\ref{s:conv}(vii)的表示产生，现在我们将其作为引擎运行。通过Vandermonde--Gram和迹-Toeplitz矩阵的等谱性，$m_b(N)=\mathbb E\int x^{b}d\nu_N$在每个有限$N$处是精确有理数（乘以$N^{b+1}$为\emph{整数}——现为定理，推论~\ref{c:struct}(i)，并在全部八十二个存档值上确认——冻结表七十九个加三个奇数阶扩展点），且计算值在$b=14$的每一阶组织为$1/N^{2}$的多项式，次数正好$\lfloor b/2\rfloor$。每一阶由$N=2,\dots,\lfloor b/2\rfloor+2$处的精确值确定，并在\emph{预注册保留点}$N=\lfloor b/2\rfloor+3$处确认，每个情形在计算前预测并逐位匹配；在战役期间次数是测量而非假设的——现为推论~\ref{c:struct}(ii)的定理——（在$b=12$时次数$5$的拟合在采用次数$6$前由其自身预测失败排除，在$b=14$时注册次数$6$对$m_{14}(9)$的预测与真值一致到六位小数，但作为精确分数不同——只有精确算术能区分）。两个独立实现（直接求和；多重集聚合带跨度动态规划，从数学规范独立编写）在每一重叠点上逐位一致。所得精确层，每层满足其预注册重新中心锚点（F-PRED-RAT）并在$N=128/192/256$处匹配三个独立$10^{7}$样本测量层，通过精确有限尺寸律$\Sigma_j(N)$（三十九次比较，全部在三个标准误差内）：
\[
\Sigma_9=\tfrac{52207}{302400},\quad
\Sigma_{10}=\tfrac{1333891}{3326400},\quad
\Sigma_{11}=\tfrac{128291}{340200},\quad
\Sigma_{12}=\tfrac{1092211019}{1556755200},
\]
\[
\Sigma_{13}=\tfrac{45789263}{52390800},\qquad
\Sigma_{14}=\tfrac{27183066233}{18162144000},
\]
等价地$m_9=\tfrac{11166011}{151200}$，
$m_{10}=\tfrac{83443081}{554400}$，
$m_{11}=\tfrac{852071287}{2721600}$，
$m_{12}=\tfrac{1033020076559}{1556755200}$，
$m_{13}=\tfrac{240004263497}{167650560}$，
$m_{14}=\tfrac{85585542088667}{27243216000}$。完整的$1/N^{2}$多项式、七十九个精确格点值（加三个奇数阶扩展点）以及十八门套件存在于复现包（\texttt{constants/tt/}）中。

三个结构事实强化了纪录；前两个现为定理（下方推论~\ref{c:struct}），第三个仍为经验观察。\emph{整数性与分母}：$m_b(N)N^{b+1}\in\mathbb Z$始终（推论~\ref{c:struct}(i)），极限分母对\emph{每个}$b$整除$(b{+}1)!$（推论~\ref{c:struct}(iii)），更锐的Lagrange界$D(b)$对$b\le13$验证——\S\ref{s:outlook}(e)的分母现象现在在所有阶都有已证天花板。\emph{单调定价}：精确Christoffel序列满足$n\lambda_n(0)$对$n\le7$严格递增（$0.2500,0.2778,0.2942,0.3012,0.3070,0.3108,0.3134$），等价于每个新正交值$p_n(0)^2$低于运行平均值——这\emph{正是}每个精确梯级略低于其测量带定价的原因。
该路径的基础是三个陈述——模型转移、分支相等、奇偶性——我们现在证明，它们依赖单一共享机制：由行列式过程\cite{Sos00}的循环簇结构（每个有限$N$处精确），迹$\operatorname{tr}U^{k_i}$的每个联合累积量是循环顺序上的带符号和，积分为$\int\prod S_N(\theta_i-\theta_{\sigma i}) \prod e^{ik_i\theta_i}$；将每个$S_N$展开为其Fourier和，角度积分逐边强制守恒，积分坍缩为\emph{格点盒计数}$(N-\mathrm{span}_\sigma(k))_+$——正好是(vii)的跨度权重。因此$m_b(N)N^{b+1}$是集合划分和循环顺序上的有限和，乘积分段多项式跨度计数的格点求和。

\begin{lemma}[模型转移]\label{l:transfer}
对每个$b$，$\lim_{N\to\infty}m_b(N)$存在且等于由\S\ref{s:conv}的重叠加权簇演算定义的正弦模型矩$m_b$。
\end{lemma}

\begin{proof}
每个（划分，循环顺序）项是格点Riemann和：缩放$k=Nv$，计数$(N-\mathrm{span})_+/N$逐点且紧致一致收敛到$(1-\mathrm{span})_+$（相同壁的分段线性函数），重叠权重将所有差变量限制在紧多面体$\mathrm{span}\le1$内，被积函数一致有界；控制收敛加Riemann求和给出每一项收敛到对应的连续重叠加权盒卷积积分，逐项即正弦模型的F-CYC定义表达式（正弦核循环在连续情形中是同一Fourier盒卷积）。项数有限。八个实例（$b\le8$）与独立证明的位置空间常数逐位一致。
\end{proof}

\begin{lemma}[平移提升：每个循环项是一个格点多面体]\label{l:lift}
固定集合划分$P=\{B\}$的$b$循环边，共$r$块，每块有循环顺序$\sigma_B$。写$k_i=m_i-m_{i+1\bmod b}$，对$S\subseteq\{0,\dots,b-1\}$，$\partial S=\sum_{i\in S}(e_i-e_{i+1\bmod b})$，使得$\sigma_B$部分和为$P_S=\partial S\cdot\vec m$，其中$S$为$\sigma_B$前缀集。则$m_b(N)\,N^{b+1}$的$(P,\vec\sigma)$项等于其Soshnikov系数乘以格点$(\vec m,\vec c)\in\mathbb Z^{b+r}$的个数，满足
\[
0\le m_i\le N-1,\qquad
\partial B\cdot\vec m=0\ \ (\forall B),\qquad
0\le c_B+\partial S\cdot\vec m\le N-1\ \ (\forall\text{ prefix }S\subseteq B),
\]
包含空前缀。
\end{lemma}

\begin{proof}
固定$\vec m$，计数逐块分解。对一块，循环累积量的守恒delta是等式行$\partial B\cdot\vec m=0$；给定它，平移计数恒等式
$\#\{c_B\in\mathbb Z:0\le c_B+P_S\le N-1\ \forall S\}
=(N-\operatorname{span}_{\sigma_B}(k))_+$精确重现盒计数（$c_B$在正时长度为$N-\mathrm{span}$的整数区间上变化）。单例块在同一公式下化简为$N\delta_{k_i,0}$，即精确一阶累积量——无特殊情形。注意提升带来的好处：$\operatorname{span}$的$\operatorname{argmax}/\operatorname{argmin}$切换壁（这是区间边界行并破坏层状结构）被$c$坐标吸收；无需腔分解。
\end{proof}

\begin{lemma}[提升壁系统的全幺模性]\label{l:tu}
对每个$(P,\vec\sigma)$，引理~\ref{l:lift}的约束矩阵——行$e_i$，$\partial B$，以及$e_{c_B}+\partial S$在坐标$(\vec m,\vec c)$上——是全幺模的。
\end{lemma}

\begin{proof}[证明（Ghouila--Houri，顶部锚定交替符号）]
由Ghouila--Houri行准则，只需对任意行子集$R$赋予符号使带符号和的所有分量在$\{0,\pm1\}$中。$R$的边界行按块分裂为嵌套链（$\sigma_B$的选定前缀，其中平凡行$\partial B$算作完整块的第二个拷贝）。\emph{在链内}，按包含排序并从最大集向下交替符号，最大集取$+1$；在平凡行$\partial B$与提升行$e_{c_B}+\partial B$的平局处，平凡行取较早符号。每条边$e$位于链的上部，因此带符号指示和$f_B=\sum\varepsilon_S\chi_S$在每条边上取$+1,-1,+1,\dots$的前缀和，即$f_B\in\{0,1\}^{B}$。\emph{跨链}块不相交，故$f=\sum_Bf_B\in\{0,1\}^{b}$全局，边界行带符号和的$\vec m$部分为循环差分$Df$，分量$f_j-f_{j-1}\in\{0,\pm1\}$。对$c_B$项：提升行的符号构成完整交替序列减去平凡行所取符号（可能有）；交替总量在$\{0,1\}$内，平凡行在平局时取$+1$（或为链顶），故$c_B$项在$\{-1,0,1\}$内。最后每个选定单位行$e_i$取符号$-\operatorname{sign}((Df)_i)$（若零则任意）。带符号和的所有分量都在$\{0,\pm1\}$中。
\end{proof}

\begin{theorem}[分支相等：一个多项式，所有$N\ge1$]\label{t:poly}
对每个$b$，$m_b(N)\,N^{b+1}$是$N$的单个多项式，次数至多$b+1$，对所有整数$N\ge1$成立。
\end{theorem}

\begin{proof}
固定$(P,\vec\sigma)$项。引理~\ref{l:lift}的每个约束形如$0\le\langle a,x\rangle\le N-1$或$\langle a,x\rangle=0$，因此该项等于$\operatorname{ehr}_Q(N-1)$，即\emph{固定}多面体$Q=\{x:0\le\langle a,x\rangle\le1,\ \partial B\cdot x=0\}$在伸缩$N-1$下的Ehrhart计数函数。由引理~\ref{l:tu}和右端数据的整数性，$Q$的每个顶点是整点：$Q$是格点多面体，因此$\operatorname{ehr}_Q$在每个$n\ge0$处是多项式——无剩余类，无阈值。其次数为$\dim Q$：等式行$\{\partial B\}$的秩正好为$r-1$（支撑不相交的$\{\chi_B\}$在循环差分的核中仅与常数相交，且$\sum_B\partial B=0$），且点$(\tfrac12,\dots,\tfrac12)$严格满足每个不等式，故$\dim Q=(b+r)-(r-1)=b+1$。将有限多项与其$N$无关Soshnikov系数相加得结论。
\end{proof}

\begin{theorem}[奇偶性]\label{t:parity}
定理~\ref{t:poly}的每个$(P,\vec\sigma)$项，从而总和，满足
$L(-N)=(-1)^{b+1}L(N)$作为多项式。
\end{theorem}

\begin{proof}
全1向量$w$（在$\vec m$和$\vec c$中）满足每个不等式行的$\langle a,w\rangle=1$——边界向量$\partial S\cdot\mathbf 1=0$，所以单位行给出$1$，提升行给出$c$分量$1$——且$\partial B\cdot w=0$。因此$x\mapsto x-w$将$nQ$的相对内部格点双射到$(n-2)Q$的格点：$\operatorname{ehr}_Q^{\circ}(n)=\operatorname{ehr}_Q(n-2)$（自反）。Ehrhart--Macdonald互反给出$\operatorname{ehr}_Q(-n)=(-1)^{\dim Q}\operatorname{ehr}_Q^{\circ}(n)$，$\dim Q=b+1$；组合后，项$T(N)=\operatorname{ehr}_Q(N-1)$满足$T(-N)=\operatorname{ehr}_Q(-(N+1))=(-1)^{b+1}\operatorname{ehr}_Q(N-1)=(-1)^{b+1}T(N)$。块数$r$已在$\dim Q=b+1$中抵消：奇偶性对每个划分形状相同，这正是在真累积量扫描中观察到的逐项局域化（$b=4$处15/15项，$b=5$处52/52，零违规）。
\end{proof}

\begin{corollary}[矩阶梯得证]\label{c:pin}
由奇偶性，$m_b(N)N^{b+1}$至多有$\lfloor(b{+}1)/2\rfloor+1$个未知系数。存档精确值（已证明有理数：Schur正交性，双实现）在每$b\le14$至少达到该数量，在每偶$b$严格过剩，且在剩余扩展点命中后，奇$b$也过剩；因此多项式唯一确定，且引理~\ref{l:transfer}将其首项系数识别为正弦模型矩。上面显示的$m_9,\dots,m_{14}$值为完全证明等级的精确有理数，$k=10,12,14$证书使用的每个常数由此确定。
\end{corollary}

\begin{corollary}[整数性、次数律、分母]\label{c:struct}
对每个$b\ge1$与$N\ge1$：(i) $m_b(N)\,N^{b+1}\in\mathbb Z$；
(ii) $m_b(N)$是$1/N^{2}$的多项式，次数至多$\lfloor b/2\rfloor$（在每个已计算阶$b\le14$处取到）；
(iii) 极限分母满足$\operatorname{den}(m_b)\mid(b{+}1)!$。
\end{corollary}

\begin{proof}
(i) 由引理~\ref{l:lift}，对象是格点计数的有限带符号和。(ii) 记$L_b(N)=m_b(N)N^{b+1}$。由定理~\ref{t:parity}，$L_b$的非零系数只出现在次数$\equiv b+1\pmod2$处。此外$L_b(0)=0$：逐项地，$T(0)=\operatorname{ehr}_Q(-1)=(-1)^{b+1}\operatorname{ehr}_Q^{\circ}(1)$，而$Q$在膨胀一处的相对内部不含格点（任何严格单位约束$0<m_i<1$对整数为空）。故出现的最低次数$\ge1$，且当$b$为奇数时$\ge2$（奇偶性），于是$m_b(N)=L_b(N)/N^{b+1}$是$1/N^{2}$的多项式，次数对偶数$b$至多$b/2$、对奇数$b$至多$(b-1)/2$。(iii) 由(i)与$L_b(0)=0$，次数$\le b+1$的多项式$L_b$在所有非负整数处取整数值，故其在二项基$\binom N0,\dots,\binom N{b+1}$下的坐标为整数，其首项系数——即$m_b$——属于$\tfrac1{(b+1)!}\mathbb Z$。
\end{proof}

(vii)中观察到的孤立扇区的周期二行为是基的伪像——特征分解的倍系数壁$2u=N$——而组装的格点多面体表示根本不具备该性质；扇区相消因而被解释而非仅仅观测。以上证明在进入本文前经过了程序的先证伪纪律：引理~\ref{l:tu}的符号构造在\emph{所有}行子集上对$b\le5$的所有$587$个壁系统进行机械验证；全子式扫描在$b\le7$的$21{,}171$个去重$\vec m$空间系统中发现零违规，在$b=8$的$197{,}357$个系统中（两腿采样覆盖，存档说明）以及引理~\ref{l:tu}的$42{,}271$个提升系统中（$b\le7$穷举；$b=7$的$37{,}633$个系统以GPU筛选加对每个被标记子式的精确整数Bareiss裁决完成）发现零违规；引理~\ref{l:lift}的完整合并项族——Fubini规模$75/541/4{,}683$于$b=4/5/6$——通过了自$N{=}1$起带保留点的单多项式拟合、逐项奇偶性以及逐位重现$m_b(N)N^{b+1}$的带符号装配（两个仅依系数规格独立编写的实现，每侧一个，且在$b=4$处CPU/GPU交叉核对）；真累积量M\"obius层在$b=4/5/6$的全部$15/52/203$个分拆上通过逐项奇偶性并带有其自身的逐位装配门；奇$b$剩余点作为已确定拟合的预测进行预注册然后计算——$m_9(8)$、$m_{11}(9)$与$m_{13}(10)$逐位命中（$b{=}13$，$N{=}10$时$3{,}243{,}531$项），使每个$b\le14$的表都拥有至少一个点的拟合余量。$\Sigma_9$的面递归第二方法已返回：$C_9=\tfrac{27649}{302400}$逐位命中其预注册目标（$60{,}739$个带符号术语轨道，$1{,}091{,}670$个F-CYC项，三机分片，自检和和与重构通过；目标由已认证$\{7,2\}=-\tfrac{4313}{12600}$的$\{9\}$联合恒等式预先固定），因此$\Sigma_9$现在有两个独立证明——分支相等链和面递归；此前未完成的委托双实现点$m_{13}(4)$，$m_{14}(3)$已逐位返回。上述均不涉及关于$\zeta$的任何假设。

\begin{theorem}[$k\le14$塔；已证明——见引理~\ref{l:transfer}，定理~\ref{t:poly}，\ref{t:parity}，推论~\ref{c:pin}]\label{t:k14}
在上面列出的矩数据$(1,m_1,\dots,m_{14})$处，$8\times8$ Hankel矩阵及其平移均正定，核多项式$Q_7^{*}$具有严格交替有理系数（注记~\ref{r:stieltjes}），Christoffel值为精确有理数
\[
\lambda_7(0)=
\tfrac{352633869846878511557783511830740995191}{7876602339133293193971616991853147607579}
=0.04476979\ldots,
\]
没有有理化代价。因此，
\[
\liminf_{T\to\infty}\frac{N_0^{s}}{N}\;\ge\;1-2\lambda_7
\;=\;
\tfrac{7171334599439536170856049968191665617197}{7876602339133293193971616991853147607579}
\;=\;0.9104604105\ldots,
\]
\[
\liminf_{T\to\infty}\frac{N_d}{N}\;\ge\;1-\lambda_7
\;=\;
\tfrac{7523968469286414682413833480022406612388}{7876602339133293193971616991853147607579}
\;=\;0.955230205\ldots,
\]
且中间梯级给出$1-2\lambda_5=0.877197\ldots$，$1-\lambda_5=0.938598\ldots$（$k=10$）和$1-2\lambda_6=0.896403\ldots$，$1-\lambda_6=0.948201\ldots$（$k=12$），其中$\lambda_5=\tfrac{46970100247159}{764967228211380}$和$\lambda_6=
\tfrac{13166900320841109317259245}{254195527518153210548497708}$。
\end{theorem}

证书是盲排演过的：预注册窗口（标题$\approx0.9101$，互异$\approx0.9551$，由带$\Sigma_{13},\Sigma_{14}$在其精确值存在前计算）事后打开并命中。同一排演量化了在此深度带级消耗是无用的——$\pm0.004$的$\Sigma_{13}$带会使标题摆动$7.6\times10^{-2}$——因此精确性不是路径的奢侈而是其关键。

\section{消耗与定理~\ref{t:main}的证明}\label{s:lp}

在$(m_1,\dots,m_4)=(1,\tfrac43,2,\tfrac{13}4)$双边固定（定理~\ref{t:m3}，定理~\ref{t:D}，\eqref{e:moments}）以及\S\ref{s:five:conn}和\S\ref{s:seven}的认证值下，剩余约束为精确有理数：
\[
m_5\;\ge\;M_5:=\tfrac{101}{18},
\quad
m_6\;\le\;M_6:=\tfrac{119}{12}+\Sigma_6=\tfrac{640}{63},
\quad
m_7\;\ge\;M_7:=\tfrac{3439}{180},
\quad
m_8\;\le\;M_8:=\tfrac{747361}{20160}.
\]
所有带贡献由精确认证退役，消耗为闭式：

\begin{lemma}[消耗步骤是Christoffel函数]\label{l:christoffel}
设$\nu$是$[0,\infty)$上的概率测度，Hankel矩阵$(H_n)_{ij}=m_{i+j}$，$0\le i,j\le n$，非奇异。则
\[
\nu(\{0\})\;\le\;\lambda_n(0)
:=\min\{\textstyle\int Q^2\,d\nu:\ \deg Q\le n,\ Q(0)=1\}
\;=\;\frac1{e_0^{\mathsf T}H_n^{-1}e_0},
\]
由归一化核多项式$Q^*=H_n^{-1}e_0/(e_0^{\mathsf T}H_n^{-1}e_0)$达到，其系数在矩为有理数时\emph{有理}。此外（符号引理）$Q^*$的所有根在支撑的内部，因此系数交替，故$P=(Q^*)^2=\sum y_jx^j$有$\operatorname{sign}(y_j)=(-1)^j$：证书始终从下方消耗奇矩、从上方消耗偶矩——正好是分类账提供的单侧数据，每度零代价。
\end{lemma}

\begin{proof}
$\nu(\{0\})=\int1_{\{0\}}\,d\nu\le\int Q^2\,d\nu$对任意$Q$满足$Q(0)=1$，因为$Q^2\ge0$且$Q(0)^2=1$；在$q_0$归一化系数向量上最小化二次型$q^{\mathsf T}H_nq$，一个Lagrange步骤给出所述闭式，最小化子为$H_n^{-1}e_0$归一化——显然有理。若$g(x)=\prod(x-r_i)$且所有$r_i>0$，则$g$的系数严格交替，故贡献$(g^2)_j$的每项$g_ag_b$带有符号$(-1)^j$无相消；$Q^*$的根由正交多项式理论在支撑间交错（它们是再生核$K_n(x,0)$的零点），因此为正。
\end{proof}

\begin{remark}[符号引理的矩数据形式]\label{r:stieltjes}
符号引理仅需矩\emph{数据}：若Hankel矩阵及其平移$(m_{i+j+1})$正定——即Stieltjes条件，对塔的精确有理算术验证（前导平移子式$1,\tfrac29,\tfrac{23}{1296},\tfrac{27857}{45722880}$；门F-STIELTJES在认证链中）——则正交多项式具有正零点，$K_n(x,0)=\sum_jp_j(0)p_j(x)$的$x^m$系数每项都带有符号$(-1)^m$，\emph{与$j$无关}：$Q_n^*$及$(Q_n^*)^2$的交替性逐项成立，无相消，在塔的每一当前和未来阶同时成立。因此认证链的逐度符号验证是交叉检查而非承重步骤。
\end{remark}

\begin{remark}[转换系数是紧的]\label{r:tight2}
从谱质量到零点计数的转换使用重数分类账：$m$重零点贡献精确$(m{-}1)$维核，因此互异$\ge1-\lambda$，且由于$(m-1)/m\ge\tfrac12$，简单$\ge1-2\lambda$——因子$2$来自全双重点极值。该因子无法从框架承载的信息改进。在矩之外，Nyquist对角坍缩（\S\ref{s:recentre}）使$\widetilde G$的归一化对角线在每个零点处等于1，因此每个双重点贡献精确$\bigl(\begin{smallmatrix}1&1\\1&1\end{smallmatrix}\bigr)$块；压缩到（相互正交的）对和向量并应用带Cauchy交错的力量迹不等式给出约束族$\sum_{\mathrm{top}\,p}\lambda_i^{\,j}\ge2^{j}p$对所有$j\ge1$，$p$为双重点数。在Christoffel极值测度（原点处原子$\lambda_6$加六个Gauss节点，由$k\le12$矩精确计算）下，该族的每个成员都是\emph{非活跃}的——顶部数量尾部均值超过阈值$+5.6\%$（$j=1$）到$+43\%$（$j=6$），余量随$j$和$k$扩大——因此全双点构型与整个信息集一致，系数$2$是紧的。在固定矩数据下对标题的任何未来改进必须来自测度侧（更小$\lambda_k$）而非转换层；配对窗口族扫描（八个方向，包括一阶$m_2$下降方向）同样发现平坦窗口$\lambda$最优，因此框架的两个设计层不携带残余松弛。
\end{remark}

\begin{lemma}[精确对偶证书，度六]\label{l:dual}
在$(m_1,\dots,m_4)$固定且$(M_5,M_6)=(\tfrac{101}{18},\tfrac{640}{63})$下，
\[
Q_3^*(x)\;=\;1-\tfrac{8232}{2519}\,x+\tfrac{7368}{2519}\,x^2
-\tfrac{1932}{2519}\,x^3 ,
\]
且对任意概率测度$\nu$在$[0,\infty)$上——无支撑界——满足$m_1=1$，$m_2=\tfrac43$，$m_3=2$，$m_4=\tfrac{13}4$，$m_5\ge M_5$，$m_6\le M_6$：
\[
\nu(\{0\})\;\le\;\int (Q_3^*)^2\,d\nu
\;\le\;\lambda_3(0)\;=\;\frac{247}{2519}\;,
\qquad
1-2\lambda_3\;=\;\frac{2025}{2519},\qquad
1-\lambda_3\;=\;\frac{2272}{2519},
\]
全部精确：单侧替换有效因为$(Q_3^*)^2$的系数交替（引理~\ref{l:christoffel}），角点求值是有限有理计算（核心Lean~4核检查，\texttt{lean/RhGate/Certificate.lean}，以及独立\texttt{certification/certify84.py}）。无求解器、网格、原子提取或有理化；早期修订的历史原子，到四位小数，是在被替代的$M_6$下$Q_3^*$分子的根（附录~\ref{a:cert}(A6)）。
\end{lemma}

\begin{corollary}[耗尽的$k\le6$塔]\label{c:2025}
使用不可有效化常数，无条件地在$k\le6$链的等级下（每个使用常数符号认证），
\[
\liminf_{T\to\infty}\frac{N_0^{s}}{N}\ \ge\ \frac{2025}{2519},
\qquad
\liminf_{T\to\infty}\frac{N_d}{N}\ \ge\ \frac{2272}{2519},
\]
且这些是六矩数据的\emph{精确}最优值：没有仅使用$m_1,\dots,m_6$的证书能做得更好。
\end{corollary}

\begin{lemma}[精确对偶证书，度八]\label{l:dual8}
在$(m_1,\dots,m_6)$如上固定且$(M_7,M_8)=(\tfrac{3439}{180},\tfrac{747361}{20160})$下，度$4$核多项式$Q_4^*$具有交替符号的有理系数，且对任意概率测度$\nu$在$[0,\infty)$上满足所述单侧数据，
\[
\nu(\{0\})\;\le\;\lambda_4(0)\;=\;\frac{12241115}{162540559},
\qquad
1-2\lambda_4\;=\;\frac{138058329}{162540559}\;=\;0.849377\ldots,
\qquad
1-\lambda_4\;=\;\frac{150299444}{162540559}.
\]
该角点在矩锥内部（在固定$m_1,\dots,m_6$下顶点$m_7=19.0712\ldots$，$m_8=36.7621\ldots$），且在$(M_7,M_8)$处连续：在分类账变量中，标题仅需要$\Sigma_7\ge0.074844$（在$\Sigma_7=\tfrac7{90}$处余量$3.8\%$）和$\Sigma_8\le0.290506$（在认证$\Sigma_8=\tfrac{633}{2240}$处余量$2.8\%$），$\Sigma_7$单调\emph{递增}尽管$m_8$中有$+8\Sigma_7$项——因此两个新和仅需单侧且精确到百分之几。
\end{lemma}

定理~\ref{t:main}通过\S\ref{s:frame}的消耗机制得出：引理~\ref{l:dual8}结合定理~\ref{t:m78}在所述等级给出标题；推论~\ref{c:2025}是全认证梯级；仅$m_4$给出$13/18$，$31/36$（推论~\ref{c:1318}，证书$Q$引理~\ref{l:cert}，首项系数为正——也无支撑界）。Dirichlet $L$函数情形逐字相同。所有常数不可有效化。\qed

\begin{corollary}[仅符号梯级]\label{c:signs}
若三个连通常数仅满足$C_5\ge0$，$\{4,2\}\le0$和$\{6\}\le0$（且$C_5\le0.1$），则由附录~\ref{a:cert}(A5)的重新优化精确证书，
\[
\liminf\frac{N_0^{s}}{N}\;\ge\;0.7403,\qquad
\liminf\frac{N_d}{N}\;\ge\;0.8701 .
\]
符号认证是比现存的精确认证（\S\ref{s:five:conn}）粗得多的目标；完整收益曲线见附录~\ref{a:cert}(A5)。
\end{corollary}

\begin{remark}[修复的中间梯级；注册D6]\label{r:rung}
本研究早期阶段引用了四、五矩链的中间梯级（$0.7295/0.8647$），“不消耗六矩信息”。认证审计发现该梯级按所述不成立：仅有$m_5$的\emph{下}界且无$m_6$（或支撑）信息时，原点质量问题在无界支撑上根本没有改善——高度$\Lambda$处的质量$\varepsilon=c/\Lambda^{4}$可使$m_5$任意增大，对$m_1,\dots,m_4$成本可忽略，审计在支撑范围$8/20/60/200$处的LP给出$0.73144\to0.72580\to0.72340\to0.72257\to\tfrac{13}{18}$。任何对偶证书若仅有五矩下界，需要负的$x^5$系数被正的$x^6$系数支配，即六矩（或支撑）信息是\emph{必要的}。因此该梯级被降级：它以所述形式仅在$m_6$上限下成立，此时被引理~\ref{l:dual}包含。完整证书不受影响（$y_6>0$）。
\end{remark}

\medskip
\noindent\emph{合并审计注记：带模型。}审计首先运行了\emph{反}相关变体（$m_5$的$\delta$与$m_6$符号相反），得到$0.78972$；上面的相关形式是正确的记账，因为$C_5$是同时移动两个约束的单个未知数。两个运行均在注册表（D3）中记录，以便标题对带模型的敏感性有记录：在（错误、双重保守的）反相关记账下标题本为$0.7897/0.8949$。现在每个带已退役且输入精确（\S\ref{s:five:conn}），记账问题不再出现；此注记作为前向记录的一部分保留。

\section{密度一极限：塔的上确界为一}\label{s:dens}

每个偶阶$2k$的塔给出无条件梯级$\liminf N_0^{s}/N\ge1-2\lambda_k(0)$，且定理~\ref{t:poly}使每个进一步梯级成为多项式成本计算。本节证明梯级的上确界为$1$：矩序列$(m_b)$是确定的，其测度在原点无原子，因此Christoffel值$\lambda_k(0)$趋于零。不声称速率——该陈述是封顶极限，而非有效进度表；有效版本是\S\ref{s:outlook}的硬边计划。我们首先将每阶输入作为一个定理独立陈述（\S\ref{s:schema}），然后证明软三元组并取上确界（\S\ref{s:soft}）。

\subsection{阶模式：每个梯级都是定理}\label{s:schema}

定理~\ref{t:dens1}对每个$k$消耗$k$阶梯级不等式。我们不将其作为证明段落留下，而是作为定理陈述并写出其证明，使用本文中已陈述的一般阶的统一组件；阶仅作为显示参数出现，不需要$k$的一致性（注记~\ref{r:nounif}）。

\begin{lemma}[每阶的纯类传输]\label{l:pure}
对每个$b\ge4$，完全锁定的$b$循环类聚合传输到其模型常数：在消耗归一化中，
\[
\bigl|\operatorname{Agg}(\{b\})-C_b\cdot O_b\bigr|
\;\ll_{A,B}\;KJ'V'^{2}(\log X)^{-A}
\qquad\text{对每个 }A>0,
\]
其中$C_b$是$b$链模型簇常数——一个多项式成本的精确有理数（引理~\ref{l:transfer}，定理~\ref{t:poly}）。对$b=3$聚合消失（$\Sigma_3=0$，\S\ref{s:five}的奇块消失）。
\end{lemma}

\begin{proof}
这是定理~\ref{t:winfty}的单因子（$r=1$）实例，因子为整条链；我们运行相同四步。\emph{(i) 局部律。}完全锁定$b$链的模$p$密度是命题~\ref{p:local}的重合计数闭式，适用于任意循环长度；无跨块合并。\emph{(ii) 阿基米德层。}游走几何$O_b$是已认证精确重叠层（\S\ref{s:five:cert}，\S\ref{s:seven:conn}），锁定链在该几何上的模型值是$b$链簇（盒卷积）常数$C_b$，其存在性和对\emph{每个}$b$为精确有理数是引理~\ref{l:transfer}和定理~\ref{t:poly}。\emph{(iii) 方差传输。}由命题~\ref{p:prod}（任意每格密度），解析方差是链自相关对公共平移求和；定理~\ref{t:D}的论证逐字重复，带有$b-1$个分类指标（\S\ref{s:writeouts}）：$(b-1)$维Parseval恒等式逐坐标成立，每个变换覆盖其全周期，全主箱是满足引理~\ref{l:MM}的加窗单素数求和，任何次因子由引理~\ref{l:minor}闭包，梳识别每坐标由引理~\ref{l:comb}进行，每坐标CRT完全正交。\emph{(iv) 截断与常数。}$\ell_1$截断由命题~\ref{p:consume}清偿，指数$P^{-(b-2)+\varepsilon}$随$b$增强；$(k,j)$聚合常数和归一化接口是\S\ref{s:fm:R}的存档认证层，与链长无关；尾部和阈值步骤与阶无关（命题~\ref{p:tail}）。所有界双边；所有常数依赖于$(A,B)$和阶$b$，通过Siegel定理不可有效化。
\end{proof}

\begin{theorem}[阶模式]\label{t:schema}
对每个$b\ge2$，在端点极限下，算术迹矩双边收敛到模型常数$m_b(T)\to m_b$，等级为\S\ref{s:verif}层(b)。因此，对每个$k\ge2$，
\[
\liminf_{T\to\infty}\frac{N_0^{s}}{N}\;\ge\;1-2\lambda_k(0),
\qquad
\liminf_{T\to\infty}\frac{N_d}{N}\;\ge\;1-\lambda_k(0),
\]
其中$\lambda_k(0)$是矩数据$(1,m_1,\dots,m_{2k})$的Christoffel值。隐常数可能依赖于$k$且不可有效化。
\end{theorem}

\begin{proof}
六个组件，每个已在一般阶陈述。\emph{(1) 分类账分解（精确）。}$\widetilde G^{\,b}$的迹展开在重合类的Bell分类账上运行；重新中心恒等式\eqref{e:recentre}是每一阶的精确算子二项式恒等式，且每个含单例的类由W2（引理~\ref{l:W2}，任意$b$循环）精确分解并被二项式系数逐阶吸收。剩余的是阶$\le b$的无单例类。\emph{(2) 三块消失。}含$3$块的无单例类在定理~\ref{t:winfty}(i)--(ii)的分解后携带$3$链主项作为因子，且该常数为零（$\Sigma_3=0$精确；算术侧逐类在\S\ref{s:five}，模型面F-CYC）。\emph{(3) 多块类}（所有块大小$\ge2$，无大小$3$，$r\ge2$）：定理~\ref{t:winfty}，在一般$b$和$\mathcal B$下证明。\emph{(4) 纯类}：引理~\ref{l:pure}。\emph{(5) 模型组装。}模型侧类值组装为$\Sigma_b$，并通过\eqref{e:recentre}组装为$m_b$：这是F-CYC簇演算，其每个常数存在且对每个$b$是精确有理数（引理~\ref{l:transfer}，定理~\ref{t:poly}）；步骤(2)--(4)将算术聚合等同于这些精确值，因此算术矩收敛到$m_b$，双边。\emph{(6) 消耗。}零点侧计数是命题~\ref{p:count}在矩数据$(1,m_1,\dots,m_{2k})$处的Christoffel函数计数，在每个偶矩上单调，带与阶无关的尾部和阈值装置（命题~\ref{p:tail}）；在每个固定$k$处给出显示不等式，正如$k\le7$在定理~\ref{t:main}和~\ref{t:k14}中实例化。无步骤依赖于$k$，除显示参数外。
\end{proof}

\noindent\emph{等级。}上述每个分析步骤是\S\ref{s:verif}层(b)链——模式不引入新等级或新装置；其内容是被消耗装置的集合是有限且固定的，阶仅作为参数进入。实例化梯级$k\le7$仍是有效的、完全写出的情形。

\subsection{决定性、原子与Christoffel极限}\label{s:soft}

\begin{lemma}[矩增长]\label{l:growth}
对每个$b\ge1$和每个$N$，$m_b(N)\le(b!)^2$；因此$m_b\le(b!)^2$。
\end{lemma}

\begin{proof}
将$\mathbb E\prod_{i}\operatorname{tr}U^{k_i}$展开为$b$循环槽的集合划分$P$，每块为联合累积量。对行列式过程，块$B$的迹统计量的累积量是对$B$槽的合并划分$Q$的和，其截断关联函数是核链循环排列的带符号和\cite{Sos00}；对迹测试函数，每个链积分通过\S\ref{s:tt}的Fourier机制坍缩为$[0,N]$内的格点盒计数，且块在其频率守恒（$\partial B\cdot\vec m=0$）时消失。三个计数完成证明。\emph{系数质量}：每块的（合并划分，循环排列）对数至多$\mathrm{Bell}(|B|)\,(|B|-1)!$，跨块$\prod_B\mathrm{Bell}(|B|)\le\mathrm{Bell}(b)$（不相交合并是特定合并），而$\sum_P\prod_B(|B|-1)!=b!$（按圈型计数的排列）；所以总系数质量至多$\mathrm{Bell}(b)\cdot b!\le(b!)^2$。\emph{守恒维数}：$r-1$个独立等式行将$\vec m$限制在至多$N^{\,b-r+1}$个格点（引理~\ref{l:tu}使解格类坐标；初等消元也足够）。\emph{盒因子}：每块$r$贡献至多$N$。组装
\[
m_b(N)\;\le\;\frac{N^{\,b-r+1}\cdot N^{\,r}}{N^{\,b+1}}
\times(\text{系数质量})
\;\le\;\mathrm{Bell}(b)\cdot b!\;\le\;(b!)^2 .
\]
（存档精确值远低于界——$m_{14}=3141.5\ldots$相对$(14!)^2\approx7.6\times10^{21}$——且组装系数质量本身经检查低于$(b!)^2$且比值递减至$b\le8$；该界的唯一作用是Carleman，阶乘平方增长足够。）
\end{proof}

\begin{proposition}[决定性与弱收敛]\label{p:det}
序列$(m_b)_{b\ge0}$是$[0,\infty)$上唯一概率测度$\nu$的矩序列（Stieltjes确定），且期望谱测度弱收敛，$\nu_N\Rightarrow\nu$。
\end{proposition}

\begin{proof}
由引理~\ref{l:growth}，
$\sum_b m_b^{-1/(2b)}\ge\sum_b\bigl((b!)^2\bigr)^{-1/(2b)}
\asymp\sum_b e/b=\infty$：Stieltjes--Carleman条件成立，因此$[0,\infty)$上的矩问题至多有一个解。测度$\nu_N$是紧的（$m_1(N)=1$精确），且对每个$b$，函数$x^b$对$\nu_N$一致可积（$m_{b+1}(N)\le((b{+}1)!)^2$一致）；由引理~\ref{l:transfer}每个矩收敛$m_b(N)\to m_b$，因此$(\nu_N)$的每个子序列弱极限具有矩$(m_b)$，由决定性等于同一测度$\nu$：全序列收敛。
\end{proof}

\begin{proposition}[原点无原子]\label{p:noatom}
$\nu(\{0\})=0$；事实上$\int\log^{-}x\,d\nu\le2-\gamma$。
\end{proposition}

\begin{proof}
\S\ref{s:conv}(vi)的有限$N$对数律精确：$\int\log x\,d\nu_N=H_N-1-\log N$。由于$\log^{+}x\le x$且$m_1(N)=1$，
\[
\int\log^{-}x\,d\nu_N
=\int\log^{+}x\,d\nu_N-\int\log x\,d\nu_N
\;\le\;1+\bigl(1+\log N-H_N\bigr)\;\le\;2-\gamma
\]
对$N$一致（$H_N-\log N\downarrow\gamma$）。函数$\log^{-}$在$[0,\infty)$上非负且下半连续（在$0$处值$+\infty$），因此由Portmanteau定理和命题~\ref{p:det}，
$\int\log^{-}x\,d\nu\le\liminf_N\int\log^{-}x\,d\nu_N
\le2-\gamma<\infty$。若在$0$处有原子则积分为无穷。
\end{proof}

\begin{lemma}[Christoffel极限 \cite{Akh,Sim98}]\label{l:chr}
对确定矩问题，Christoffel值随$k$递减并收敛到点质量：
$\lambda_k(0)\downarrow\nu(\{0\})$。
\end{lemma}

\begin{theorem}[密度一：几乎全部零点简单且位于临界线上；等级为\S\ref{s:verif}层(b)]\label{t:dens1}
无条件地，
\[
\lim_{T\to\infty}\frac{N_0^{s}(T,2T)}{N(T,2T)}=1,
\qquad
\lim_{T\to\infty}\frac{N_d(T,2T)}{N(T,2T)}=1,
\]
且对任意固定本原Dirichlet $L$函数的零点同样成立。
\end{theorem}

\begin{proof}
由定理~\ref{t:schema}，对每个$k$，$\liminf_T N_0^{s}/N\ge1-2\lambda_k(0)$。左边是独立于$k$的单个数，因此它支配$\sup_k(1-2\lambda_k(0))$——可数个真命题的上确界，无极限交换（注记~\ref{r:nounif}）。由引理~\ref{l:chr}和命题~\ref{p:det}、~\ref{p:noatom}，$\sup_k\bigl(1-2\lambda_k(0)\bigr)=1-2\nu(\{0\})=1$。对$N_d$以$1-\lambda_k$同样论证，对任意固定本原Dirichlet $L$函数亦然，其链相同（\S\ref{s:frame}，定理~\ref{t:schema}）。
\end{proof}

\begin{remark}[为何不需要$k$的一致性]\label{r:nounif}
极限不涉及交换或外推，将逻辑以两行显示值得。对每个\emph{固定}$k$，梯级
$\liminf_T N_0^{s}/N\ge1-2\lambda_k(0)$是一个独立有限定理：$T\to\infty$在固定$k$下取，其常数可依赖于$k$（且不可有效化）而无妨。左边是独立于$k$的单个数；因此它支配每个梯级，即
$\liminf_T N_0^{s}/N\ge\sup_k(1-2\lambda_k(0))$——可数个真命题的上确界，非极限交换，无需$k$一致误差控制。该上确界等于$1$是引理~\ref{l:chr}和命题~\ref{p:det}、~\ref{p:noatom}——决定性、精确对数律和经典矩论——其中没有有限阶数值出现：存档矩表在定理~\ref{t:dens1}的证明中不起作用，只在其有效实例化梯级中起作用。特别地，沿途不可能发生主项与误差项的倒置：固定$k$时解析误差为$O_k((\log X)^{-A})$，在$k$移动之前即已消亡；而上确界一步根本不含解析误差项。（引理~\ref{l:growth}的$(b!)^{2}$是\emph{极限矩本身}的增长界，供Carleman判据消耗——它是决定性的燃料，不是误差估计。）
\end{remark}

\begin{remark}[声明与未声明]\label{r:dens}
该定理是密度陈述，无有效速率：它不给出$\varepsilon$之外的$k(\varepsilon)$，也不断言所有零点简单或位于线上。有效进度表——$k$须多大才达到$1-\varepsilon$——受$\lambda_k(0)$的衰减控制，经验上$\lambda_k(0)\approx0.33/k$（\S\ref{s:outlook}）；将该速率证明为定理是硬边计划，本节故意绕过：决定性、精确对数律和盒计数多面体的互反性都是软输入，均已证明。消耗的分析层为\S\ref{s:verif}层(b)，定理的等级正是该层的等级；不涉及关于$\zeta$的假设。
\end{remark}

\section{数值验证}\label{s:faces}

独立管道（\texttt{pipeline/}；\texttt{numpy}，\texttt{scipy}）在几分钟内重新运行每个门。中心高度表（十六个互素模数族）：

\begin{center}
\begin{tabular}{lcccccccc}
$T$ & 2400 & 4800 & 9600 & 19200 & 38400 & 57600 & 76800 & 153600\\
\hline
色散（$10^{-3}$） & 8.39 & 5.19 & 4.16 & 4.06 & 3.87 & 1.84 & 1.29 & 1.04\\
跟踪（\%） & 1.5 & 1.0 & 0.7 & 0.9 & 0.8 & 0.4 & 0.2 & 0.2\\
\end{tabular}
\end{center}

\noindent（全范围几何平均每倍增$0.706$：纯平均律$2^{-1/2}$）。进一步面：局部闭包精确（零违规，七个配置）；Parseval精确；每箱Siegel--Walfisz面在范围上下降十倍；Bell(5)/Bell(6)门精确；五、六重积面在所有高度通过；良好间隔零违规；混合大小控制$24$--$31$倍高度稳定；消耗LP网格稳定。注册表：$478$个预注册检查，$436$个通过，$58$个触发并转换，全部前向记录（第$31$--$35$个转换属于修正常数通过、Lean编译和谱路径审计；第$36$--$41$个属于模型边缘审计、$\{2,2,2\}$分类修正及其后果、以及$C_8$候选偏差，\S\ref{s:verif}(f)，附录~\ref{a:record}，D19--D22）。

2026年8月16日的认证审计重新运行：全矩链门套件在$T=9600$（全部通过，$72$s：色散$4.163\times 10^{-3}$，跟踪中位数$0.73\%$，SW RMS $1.599\times10^{-3}$，紧度$24.8$--$31.3\times$，五重$0.06$--$1.29\%$，六重$0.01$--$0.73\%$，LP在退役端点常数下为$0.79454/0.89727$两网格；修正常数重跑给出$0.79472/0.89736$）；四矩恒等式套件在$T=9600$（全部通过，$48$s）；\S\ref{s:fm}的常数套件（\texttt{m1\_suite.py validate}：$6/6$局部因子，$100/100$向量筛，锚点和跨约定差重现；核心$-0.00545$相对存档$-0.00544$；$\alpha$截断$-0.00491$）；锯齿认证（\texttt{sawtooth\_cert}：网格$1.8675$，样本$1.8419/1.8494/1.8512$，证书$2.100$）；连续求积门（\texttt{quadrature\_cert}：几何尾比$0.500$认证，$c_0=-1.4228$十倍程间稳定）；以及\S\ref{s:cert}的证书套件（\texttt{certificate\_verification.py}精确；\texttt{moment\_lp\_reopt} $0.0721$；\texttt{mu3\_ledger\_gates}；\texttt{cumulant\_engine}门$C_2$精确，$\int O_3C_3=0$）。消耗LP的独立重实现（不同代码，HiGHS后端）重现了\S\ref{s:lp}中记录的每个梯级。

第二次认证通过（同日）增加了：精确配对层积分（\texttt{exact\_t222.py}：命题~\ref{p:t222}的所有值，双方案相同有理数，闭式交叉检查，及猜测有理数的匹配）；精确对偶证书（\texttt{certify\_lp.py}：展开，两个带角点，最终不等式，全部精确）；写出面（\texttt{writeouts\_verification.py}：自由槽Fubini精确，窗口质量$1.00041$，三对积模型中位数$0.9814$，解析Parseval $1.2\times10^{-16}$，解析CRT对比——一个面在构造期间触发并转换，注册表第30个）；以及支持范围探测（注记~\ref{r:rung}）。

修正常数通过重新运行：分组中点网格梯用于$C_5$，$\{4,2\}$，$\{6\}$（\S\ref{s:conv}；检查点引擎，$b=4$校准相对命题~\ref{p:t222}精确）；修正角点的精确对偶证书（\texttt{certify\_lp.py}：展开，两个带角点，最终不等式，全部精确，最坏角点$w_0=\tfrac{1153107070889}{11233957316589}$）；卷积迹门（F-TRACE：$b=2$精确$\tfrac13$，$b=3,4$至$3\times10^{-14}$）；以及——程序中首次——在维护者工具链上对恒等式层的真正Lean内核编译（\S\ref{s:verif}(f)）。

识别通过（8月17日）运行加速梯引擎（支持剪枝$171\times$在$b=6$，项trie，融合核，否定对称性——在数学上与参考引擎相同，逐门验证并跨V100/消费级GPU/256线程CPU路径逐位验证）到每个常数的七个Romberg梯级：$C_5$到$dv=0.00125$（$3200^4$网格），$\{6\}$到$dv=0.0125$（$320^5$），$\{4,2\}$旁观者对变体到$dv=0.0015625$（$2560^4$；总墙时$234$s），$C_7$到$dv=0.025$（$160^6$）。在精细运行前预注册的三个预测门（F-RAT，来自粗梯的$h^{2}$模型拟合）全部确认——对$\{4,2\}$，精细梯\emph{精确命中预注册分数}$-\tfrac{23}{420}$；相同管道在相同运行中以$10^{-18}$重现精确锚点（$b=3$为$0$，$b=4$为$-\tfrac1{60}$）。

符号认证通过（8月17日）关闭N1：三个常数的每一项都提升为有理多面体体积，并通过\S\ref{s:five:conn}的精确面递归求值，$1310$个精确分数每项存档，完整$\{6\}$和在$9.1$单核小时内完成。（分段路径作为独立第二方法保留，在相同$\{6\}$项集上投影到$\sim10^{8}$小时——多面体提升在$b=6$处是$\sim10^{9}$倍的算法约化。）预注册证伪者F-SYM-N1——符号和与识别分数之间的任何不匹配——未触发：符号值在三个常数及所有三个$\{4,2\}$放置分量上与数值重构逐位一致。

\begin{remark}[有效面，无效定理]
面是有限高度测量，带有有效常数；定理常数无效（Siegel）。面验证精确恒等层、模型识别以及证明预测的定性平均行为；它们不能且不验证渐近性。两类证据全文保持区分。
\end{remark}

\section{验证状态}\label{s:verif}

\emph{声明等级：已认证候选者。}链=经典输入（Siegel--Walfisz \cite{Dav,IK}，Vaughan，奇异级数截断，\S\ref{s:fm}的\cite{MRT}）+存档认证引理（桥，胶，聚合，D1截断，消费者接口，锚点证书：\S\ref{s:fm}）+ \S\S\ref{s:cert}，~\ref{s:local}--\ref{s:five}的机器验证恒等式 + \S\ref{s:five:conn}的符号认证连通常数 + \S\ref{s:frame}的消耗层。按层分级：

\medskip
\noindent\emph{阅读指南（标题定理，各消耗什么，何等级）。}机器验证与精确有理层(a)和人类解析层(b)在全文中物理隔离：每条定理的标题携带其等级，希望集中审查解析链的评审者可以只关注下列(b)层条目，并对照\S\S\ref{s:trunc}--\ref{s:four}与\S\ref{s:dens}的逐引理结构进行审计。
\begin{itemize}
\item 定理~\ref{t:main}（$0.8493/0.9246$）：精确四矩（层(b)分析链）+ 符号认证$k\le8$常数（层(a)）+ 消耗LP（精确）。独立于\S\ref{s:fm}和\S\ref{s:tt}。
\item 定理~\ref{t:k14}（$0.9104/0.9552$）：上述 + Toeplitz--迹矩$m_9..m_{14}$，现已证明精确有理数（分支相等和奇偶性定理，\S\ref{s:tt}；证明等级有限表，带双实现和预注册保留点）。
\item 定理~\ref{t:dens1}（密度一）：阶一致算术装置（层(b)）+ \S\ref{s:dens}的软三元组（决定性、无原子、Christoffel极限——已证明）+ 无数值表且无速率。
\item 定理~\ref{t:fm}（$0.6916/0.8458$）：一个计算机辅助梯级（引理~\ref{l:CL}数值分析步骤）；隔离——\S\ref{s:fm}外部不被任何东西消耗。
\end{itemize}
\medskip

\emph{(a) 机器精确或精确有理，在认证审计中重新验证。}证书代数（引理~\ref{l:cert}）和$\kappa(c)$律；局部闭包（命题~\ref{p:local}）；Parseval和乘积恒等式；Bell分类账和冻结槽恒等式；无锚$67/12$和$39/4$；$197/60$分类账算术及现在完全精确的双记账对账（注记~\ref{r:books}）；通过二面体放置类的精确配对层常数$t_{\mathrm{adj}}=\tfrac7{60}$，$t_{\mathrm{opp}}=\tfrac1{30}$，$\{2,2,2\}=\tfrac{32}{105}$（命题~\ref{p:t222}，附录~\ref{a:cert}）；符号认证连通常数$C_5=\tfrac1{36}$，$\{4,2\}=-\tfrac{23}{420}$，$\{6\}=-\tfrac1{126}$（\S\ref{s:five:conn}；注册D18）；精确对偶证书（引理~\ref{l:dual}）——消耗步骤不涉及LP求解器、网格或支撑界；锯齿常数$C_{\mathrm{saw}}=2.10$（网格加解析认证）。

\emph{(b) 已认证候选者分析层（预Lean，预评审）。}九类$\mu_3$分类账（两个非标准步骤在存档备忘录中写出）；截断清偿；引理D的谱证明（定理~\ref{t:D}）——其次/混合箱大小记账经典但在写出等级携带并经过面检查预算——及其到$k=5,6$的逐因子传输，包括写出的W1--W3（\S\ref{s:writeouts}），它们以该层等级证明；\S\ref{s:fm}的覆盖区/带/胶链。冻结链前有五轮对抗审计（存档轮次r93，r95，r100，r102）。

\emph{(c) 写出条目。}已处理：W1（命题~\ref{p:W1}），W2（引理~\ref{l:W2}），W3（引理~\ref{l:W3}），每个带预注册面。W2完整（初等）；W1和W3以层(b)等级证明，因为它们消耗相同的谱传输。

\emph{(d) 认证值。}N1状态：\emph{已关闭}。所有连通常数已证明有理（\S\ref{s:conv}），三个已用常数已\emph{证明}——每个是带每项存档的精确有理多面体体积计算，双方逐分式一致，所有校准门通过（\S\ref{s:five:conn}）——以证明等级消耗，即命题~\ref{p:t222}的等级。数值识别（机器精度Romberg极限，相邻行漂移$2\times10^{-16}$，唯一性保证有理重构，双精确锚定校准在同一管道中，跨架构逐位一致，预注册$h^{2}$预测门，以及$\{4,2\}$的预注册命中并逐分量识别）作为独立确认。
$k=7,8$层（\S\ref{s:seven}）逐项分级：$\{5,2\}$，$\{2^4\}$，$C_7$，$C_8$为证明等级——$\{5,2\}$，$\{2^4\}$和$C_7$有两个独立确认（预注册命中；$17/17$梯一致；精确命中注册识别），$C_8$通过轨道约化符号运行，模型侧一致到$10^{-4}$；$D_8$对称前提审计\emph{完成}（六个轨道，每轨道三个成员，相同有理数），原始项蛮力复制故意停止以将计算资源重新分配给九矩层——记录决定：$C_8$的谱系是符号运行加模型侧A/B/C检查，而非第二条独立符号路径）；$\{2,2,4\}$，$\{4,4\}$，$\{6,2\}$为证明等级，其第二路径GPU梯\emph{已返回并通过}（门F-V-JOINTS：$22/22$，$7/7$，$4/4$轨道在修订$3\times10^{-6}$带内；首次比较因轨道序对齐导致两个类不匹配，由其自身模式捕获并用量子化规范键修复——教训，\emph{内容规范键和对每次跨工具链比较的断言匹配计数}，前向记录）——$k\le8$链的常数层在双路径证明等级完成。定理~\ref{t:main}的标题承载这些等级中最弱者，连同W$_\infty$传输（定理~\ref{t:winfty}，实例W4--W7）的层(b)等级；推论~\ref{c:2025}不消耗其中任何，处于完整$k\le6$认证等级。附录~\ref{a:cert}(A5)的\emph{证书族}保留为分级插头记录。N2（部分处理）：求积的斜率精确证明为$1$，其常数以闭式确定$c_0=\gamma-2$（命题~\ref{p:c0}），几何尾部已认证（r60）；剩余计算的是$c^{*}$内部的平均振荡$\overline{\mathrm{osc}}$和普适性坍缩速率写出。仅保守收费$\varepsilon_{CL}$在本文任何地方消耗（引理~\ref{l:CL}）。N2完全认证将定理~\ref{t:fm}升级为$0.7031/0.8515$。

\emph{(e) 导入。}\cite{C26}：\S\ref{s:frame}的框架（矩阵，块结构，尾部，计数，前两矩）——本程序公开预印本，未经独立同行评审；其自身的分析输入为已发表文献。直接消耗的已发表文献：\cite{Wei52}和\cite[定理5.12]{IK}（显式公式），\cite[定理9.2]{Tit86}（零点密度），\cite{Mon73,BGST}（素数侧配对相关求值，带内无条件），\cite{MRT}，\cite{HR}，\cite{Dav}。

\emph{(f) 形式化状态。}
恒等式层现已在维护者机器上\emph{编译并核检查}：核心Lean~4（工具链\texttt{leanprover/lean4:v4.33.0}，无\texttt{mathlib}），包\texttt{lean/RhGate/}，构建日志包含。编译内容：无锚组装恒等式$m_4=\tfrac{13}4$，$m_5$-部分$=\tfrac{67}{12}$，$m_6$对层$=\tfrac{39}4$作为全称量化有理恒等式（任意锚点$t_{\mathrm a},t_{\mathrm o}$——\texttt{Anchors.lean}）；消耗算术$\tfrac{13}{18}$，$\tfrac{31}{36}$；以及$b=4$，$p=5$处的重合计数局部律通过全部$125$个锁定类的完全核枚举（\texttt{LocalClosure.lean}——该定理\emph{无公理}；其他仅使用\texttt{propext}，\texttt{Classical.choice}，\texttt{Quot.sound}）。不含任何未完成的证明占位（\texttt{sorry}）。编译本身触发了一个注册检查并成为第$33$次转换：种子文件的证明使用了\texttt{ring}和\texttt{decide}在$\mathbb Q$上，但核心Lean中均不可用（\texttt{ring}是\texttt{mathlib}策略；核心\texttt{Rat}算术不是核归约的，因此\texttt{decide}在\texttt{instDecidableEqRat}上停滞）；所有五个受影响证明用核心\texttt{grind}策略修复，陈述不变（附录~\ref{a:record}）。同日第二编译（8月16日，20:09）覆盖证书层：\texttt{Certificate.lean}（去分母展开恒等式引理~\ref{l:dual}，非负性，归一化，两个角点求值$w_0$，角点选择符号——因此\emph{选择}绑定角点本身由核检查——标题不等式，附录~\ref{a:cert}(A6)的幺正展开，精确$\{2,2,2\}$算术及实例化锚点恒等式；第三编译（8月17日）覆盖先前修订——\texttt{Certificate.lean}在($M_5=\tfrac{101}{18}$，$M_6=\tfrac{12809}{1260}$，标题$\tfrac{7962}{10000}$/$\tfrac{8981}{10000}$)，十三个证书层定理，干净构建，公理审计如所述；本修订的证书层——校正$\{2,2,2\}$束，$M_6=\tfrac{640}{63}$，有理核多项式$Q_3^*$，$Q_4^*$及精确标题$\tfrac{2025}{2519}$，$\tfrac{138058329}{162540559}$——已编写（\texttt{Certificate84.lean}）并\emph{排队进行同机重编译}）和\texttt{LocalClosure2.lean}（$b=4$，$p=7$和$b=5$，$p=5$局部律：$343$和$625$个锁定类，通过核心\texttt{decide}在$\mathbb N$上——两者\emph{无公理}）。该编译第二次触发注册表（第$34$次转换）：\texttt{grind}的环/线性核心将平方$(\cdot)^2$视为不透明原子，无法关闭完全平方非负性；修复是\emph{核心}有序环引理\texttt{Lean.Grind.OrderedRing.sq\_nonneg}（项模式，陈述不变，仍无mathlib）。预编译种子\texttt{SOSCertificate.lean}——其标题常数携带退役端点网格角点——从包中移除，以防止共存的$0.7946$/$0.7947$标题定理；其仍有效内容吸收到\texttt{Certificate.lean}中，草稿保存在程序存档中。每个陈述在精确有理算术（\texttt{certify\_lp.py}）或完全\texttt{Python}枚举（零偏差）中代理验证。Lean认证的是有理算术层本身：矩固定、测度论Chebyshev--Markov步骤和分析链仍按层(b)--(e)分级。截断、谱和消耗\emph{分析}仍未被形式化；分析层的形式化——依赖\cite{C26}框架——等级待定。伴随\cite{C26}的形式化包与此存储库分离；本文编译覆盖正是上面所述。

\emph{(g) 门。}按程序的验证门，纪录声明遵循（i）N1分数的符号验证和$\{4,2\}$闭包——\emph{现已处理}（\S\ref{s:five:conn}，(d)上）——和可选N2，（ii）恒等式和证书层的编译Lean形式化，队列Q1--Q6（两者现已编译，\S\ref{s:verif}(f)；分析层保留），以及（iii）外部评审。本文未使用Riemann假设、零点密度输入、配对相关猜想或Kloosterman技术。本文内容不涉及RH本身。

\section{展望：速率、结构与墙}\label{s:outlook}

塔已被推至上确界（定理~\ref{t:dens1}），展望不再关乎更高的比例——那个问题已经关闭。剩下的事情分三类：\emph{有效性}（已实例化梯级逼近极限的速率）、\emph{结构}（测度$\nu$及其常数的闭式）与\emph{墙}（这一切与Riemann假设之间的分界，在(c)条末尾的插值地图中精确标出）。我们逐条记录，附手头的定量证据。

\emph{(a) 完成当前链。}指定闭包中，第一个——三个连通分数的符号验证，将条目N1从识别移至证明——已在本文完成（\S\ref{s:five:conn}）；剩余的是分析层的形式化和外部评审。它们改变等级，不改变常数。

\emph{(b) 完成$k=8$等级，及第九、十矩。}$k=8$梯级在\S\ref{s:verif}(d)所述等级关闭；剩余的是Lean扩展和评审。\emph{下一个}梯级正在构建，其侦察是定量的。在精确侧，第九、十阶分类账的十一个无单例类中，五个已认证：
\[
\{2^5\}=\tfrac{10531}{13860},\quad
\{4,2,2,2\}=-\tfrac{208771}{498960},\quad
\{2,2,5\}=\tfrac{2263}{5040},\quad
\{5,4\}=-\tfrac{257}{10080},\quad
\{6,2,2\}=-\tfrac{120139}{498960},
\]
第一个现已带完整每轨道第二路径梯（$79/79$轨道在$2.1\times10^{-6}$内；其最大轨道是\S\ref{s:conv}(v)的盲预测$\tfrac{73}{11088}$），第二个由两个不同并行化主机返回相同，第三个仅单认证路径（其梯排队）。\S\ref{s:conv}(v)的闭式还\emph{预注册}第十二、十四分类账中的两个有理目标——顺序轨道$\tfrac{523}{192192}$和$\tfrac{119}{102960}$，该族的首批分母被$13$整除，正是(e)的B样条机制预测——等待独立面递归确认（门F-SEQ-12/14，前向记录）。在测量侧，精确有限模型的高统计运行（$N=128$，$2\times10^{5}$ CUE样本；六个已知中心矩$\Sigma_2,\dots,\Sigma_8$均在$0.4$标准误差内重现，预分配$\Sigma_7$门为$+0.1$）给出
\[
\Sigma_9=0.172618(61),\qquad
\Sigma_{10}=0.400778(108),
\]
重新中心闭式$m_9=\tfrac{495107}{6720}+\Sigma_9$，$m_{10}=\tfrac{199427}{1344}+10\Sigma_9+\Sigma_{10}$由独立Bell$(9)$/Bell$(10)$分类账组装精确重现（门F-PRED-RAT）。Toeplitz--迹引擎此后在证明等级关闭了所有六个进一步层——值、证据矩阵、$k\le14$证书以及支撑它们的分支相等和奇偶性定理是\S\ref{s:tt}的内容；委托双实现点$m_{13}(4)$和$m_{14}(3)$已逐位返回。此前的剩余是冗余而非证明：$C_9$的面递归已逐位命中其预注册目标$\tfrac{27649}{302400}$（\S\ref{s:tt}），给$\Sigma_9$第二个独立证明；第十、十二层联合类仍是已证明常数的进一步第二方法。

\emph{(c) 塔的视界。}矩路径不能停在任意固定比例——排斥力迫使$w_0(k)\to0$当所有矩被固定——但每梯级的边际增益现在受LP极值的隐含矩与正弦模型矩之间的失配控制，精确数据现已量化：精确Christoffel序列给出$n\lambda_n(0)$递增，$7\lambda_7=0.3134$，二点外推$A\approx0.330$，$B\approx0.114$在$n\lambda_n\approx A-B/n$——与早期$n=10..15$处池测量$A\approx0.30$有张力，最可能因为这些测量携带了现已成为精确定理的有限$N$矩偏差；刷新带级定价为$k=16\to\approx0.921$，$k=18\to\approx0.930$，$0.95$在$k\approx26$，$0.99$在$k\approx120$（外推律）。由于定理~\ref{t:poly}和~\ref{t:parity}对\emph{每个}$b$成立，任何未来梯级一诞生即为证明等级：只需$\lfloor(b{+}1)/2\rfloor+1$个精确有限$N$求值——多项式成本——而非新结构论证。有了定理~\ref{t:dens1}，塔的目的地不再有疑问——只有其进度表：这里仍开放的是\emph{速率}$\lambda_k(0)\approx0.33/k$，即$\nu$在$0$处的硬边密度和Máté--Nevai--Totik常数，其证明路径是下面的重求和计划。为记录该计划的结构提示：对数谱律积分到累积量生成函数$\Lambda(s)=(1{+}s)\log(1{+}s)-s-\log\Gamma(2{+}s)$，且$e^{\Lambda(b)}=(1{+}b)^{1+b}e^{-b}/(1{+}b)!$恰好是\emph{临界} Poisson Galton--Watson树的大小偏置总后代（Borel）序列——$\nu$的对数谱侧是一个临界分支对象，暗示$\nu$在该族内的精确识别。

\medskip
\noindent\emph{插值地图，以及本文没有做的事。}值得精确记录定理~\ref{t:dens1}相对黎曼猜想所处的位置，因为本框架把这段距离刻画得很精确。方法有两个旋钮：所消耗素数侧数据的带宽$\lambda$（Montgomery归一化）与矩阶$k$：$(\lambda{=}1,k{=}2)$给出$2/3$；$(\lambda{=}1,k{=}14)$给出$0.9104$；$(\lambda{=}1,k{=}\infty)$即定理~\ref{t:dens1}——密度一，\emph{带限平均世界的上确界}，现已达到；而Weil判据（对\emph{全部}检验函数的逐个正性）就是黎曼猜想本身。剩余的距离不是一个百分比，而是输入种类的更换，且框架精确定位了它。距离线$\delta>0$的零点通过其函数方程镜像对进入带限形式，三线（对数凸性）约束压制了它在单尺度$g\star\bar g$形式中的效应；其全额指数权重$T^{2\lambda\delta}$只能被\emph{双尺度}检验结构恢复——而这恰是配对相关函数$F(\alpha)$的$\alpha>1$域，即Hardy--Littlewood素数对信息。在使这类数据无条件可评估的高度平均之下，离线信号\emph{线性}于离线零点的个数：带宽$\lambda$处信号的缺席只能翻译为形如$N_{\mathrm{off}}(\delta,T)\ll T^{1-2\lambda\delta+\varepsilon}$的零点密度指数，在任何有限$\lambda$处都到不了空集——稀疏性使平均信号按比例减弱，平均世界内不存在稀疏性杠杆。黎曼猜想恰是$\lambda\to\infty$的角，等价地：亚间距分辨率的逐高度（非平均）素数数据。本文中没有任何陈述越过这堵墙，任何带内改进也不能；我们记录此地图，以免带内进展——包括密度一封顶——被误读为向它的逼近。

\medskip
\noindent\emph{跨越双矩屏障：消耗了什么、没消耗什么。}熟悉经典文献的读者会问：在Montgomery的无条件范围之内，这一切如何可能？Fourier支撑在$(-1,1)$内的配对相关信息，众所周知单凭自身不可能把比例推过熟知的天花板——本文同意，且量化了自己的实例（仅消耗带宽一数据的任何证书$\le 0.68185$，即$(\lambda{=}1,k{=}2)$角）。答案是：矩阶梯从不拓宽频带，它沿$k$攀升。Montgomery定理只喂养前两矩。每个更高的矩都是\emph{单独的、单独证明的、无条件的}算术输入：经显式公式与单元格分类账，它化为$\Lambda$的高阶自相关和（单例无关类的多线性素数和），由Siegel--Walfisz定理与Vaughan估计在模$\le(\log X)^{B}$内闭合（定理~\ref{t:D}、定理~\ref{t:winfty}、引理~\ref{l:pure}）。没有任何一步拓宽配对相关的Fourier支撑，没有引入任何广义假设，也没有奇迹式相消；屏障是$(\lambda,k)$平面上的一个点，不是横亘其上的一堵墙。

\medskip
\noindent 
每$k$的记账现已在\emph{两边}自动化：算术侧由定理~\ref{t:winfty}一劳永逸（每个进一步阶都是实例），模型侧由\S\ref{s:tt}的有限$N$路线以多项式成本完成。一条历史注记，保留它是因为其中藏着一个仍开放的结构问题：产出$b\le8$常数的位置空间轨道约化积分器，成本约每级$11$倍增长，在$b=10$附近到达经济边界；那里类的内在复杂度参数是图的\emph{交叉宽度}（非交叉扇区坍缩为\S\ref{s:conv}(v)的闭式，$C_8$项的$87\%$宽度$\ge4$）。该成本壁已被\emph{取代}而非翻越：定理~\ref{t:poly}与~\ref{t:parity}把每一阶改道为$\lfloor(b{+}1)/2\rfloor+1$次精确有限$N$求值，位置空间路线则作为独立第二方法存续（$C_9$面递归即其实例）。交叉宽度分析的意义仅存于(e)条的闭式问题。

\emph{(d) 超越矩。}矩是全局泛函，而原点质量由局部性质（排斥：间距密度$\sim x^{2}$）消除。接受局部测试函数的消耗机制——非多项式函数上具有算术可计算值的Christoffel型对偶——将绕过塔的递减收益；压缩矩阵框架是否允许一个，在我们看来是本文提出的最有意义的开放问题。任何未来此类机制若在\emph{单个}零点层面论证，必须通过我们在此记录为永久预注册门的判别器：它必须对Davenport--Heilbronn函数（共享函数方程和零点对称性但违反RH模拟）显式失败——一个无法区分$\zeta$与Davenport--Heilbronn的论证消耗了零算术并证明不了任何东西。

\emph{(e) 结构问题。}程序认证或识别的每个常数——$\tfrac13$，$\tfrac7{60}$，$\tfrac1{30}$，$-\tfrac1{60}$，$\tfrac{32}{105}$，$\tfrac1{36}$，$-\tfrac{23}{420}$，$-\tfrac1{126}$，$-\tfrac{17}{360}$——都是分母整除$2520=\mathrm{lcm}(1..9)$的有理数；在等级$8$该律失效（$3780\nmid2520$对$\{2^4\}$，轨道分母高达$90720$），在等级$10$出现素数$11$（$\{2^5\}=\tfrac{10531}{13860}$，$\{4,2,2,2\}=-\tfrac{208771}{498960}$，两者分母有$11$）——因此若存在正确不变量，它是等级依赖的，可能是$\mathrm{lcm}(1..b{+}1)$倍数的除数。\S\ref{s:conv}(v)的非交叉演算解释了机制：构建块是中心B样条值（$\int\mathrm{sinc}^{2m}=\beta_{2m-1}(0)$及其相关），其阶乘族分母随等级增长，正如观察。在九个精确矩上的结构探测（过确定P递归和低次代数拟合均被排除；无乘积结构的Hankel行列式；$\lambda$分母分裂为大素数）排除了谱测度本身最便宜的可积族。\S\ref{s:conv}的维数定理解释了有理性但未解释小分母；这些多面体体积的闭式理论（同时产生所有$C_b$）将折叠每个未来梯级的模型侧。

\emph{(f) 对角极限：已经达成。}本纲领的早期版本在此记录过一条归约：“百分之百需要每个梯级达到证明等级，加上$\nu$的原点质量。”该归约已被兑现——而且颇具启发性地，不是沿它当年预想的路线。预想路线用轨道积分器的成本壁为极限定价；实际的证明（\S\ref{s:dens}）根本\emph{不需要}任何计算常数：决定性（引理~\ref{l:growth}，且不对$\nu$作任何紧支撑假设）、精确对数律与经典Christoffel极限以软论证收束上确界，而有限梯级作为其有效实例存续。旧条目中存活下来的是它的边界句，我们重申如下：定理~\ref{t:dens1}是密度一陈述，非Riemann假设，非关于每个零点的陈述；剩余之墙的精确位置见(c)条的插值地图。

\section*{致谢}

本文的数学发展——\S\ref{s:frame}的框架、证书演算、引理D的谱证明、矩识别、卷积演算以及复现和Lean包——由Claude（Anthropic）——一个大型语言模型——在作者指导的研究程序内完成，采用\S\ref{s:verif}和附录~\ref{a:record}中描述的“先证伪”验证纪律。作者指定目标、仲裁决策、在自身工具链上编译和审计Lean包，并承担全部学术责任。本文基于同一研究程序的预印本三分之二定理\cite{C26}。

\appendix

\section{梳识别，附工作对偶}\label{a:comb}

在$\tau$框架中，$\rho(\tau)$的模型由孪生奇异级数组装：$\SSS$为Hardy--Littlewood级数$\SSS(h)=2C_2\prod_{p\mid h,\,p>2}\frac{p-1}{p-2}$（$h$偶），窗口重叠计数$V_m,V_n$，以及精确对角项，
\[
\mathrm{Model}(\tau)=\!\!\sum_{\substack{t\equiv0\ (b_2)\\
t\equiv\tau\ (b_1)\\ t\notin\{0,\tau\}}}\!\!
\SSS\!\Bigl(\frac{t}{b_2}\Bigr)\SSS\!\Bigl(\frac{t-\tau}{b_1}\Bigr)
V_m\!\Bigl(\frac{t}{b_2}\Bigr)V_n\!\Bigl(\frac{t-\tau}{b_1}\Bigr)
+\ \text{对角项}.
\]
CRT同余$t\equiv0\ (b_2)$，$t\equiv\tau\ (b_1)$是联合框架的\emph{完全}正交性：平移$t$贡献当且仅当它同时是$m$侧的$b_2$格孪生平移和$n$侧的$b_1$格孪生平移（偏移$\tau$）；无不完备和，因为$\tau$变换在整个周期上运行。$\SSS$在$p\mid b_1b_2$处的局部因子正是命题~\ref{p:local}在$b=4$处的$\kappa$值；对$p\nmid b_1b_2$为孪生因子。

\emph{工作对偶}$(b_1,b_2)=(5,3)$：$M=15$，对$\tau\equiv\tau_0$，幸存平移是单一剩余类$t\equiv t_0(\tau)\ (15)$，$t_0=3\cdot(2\tau\bmod5)\bmod15$（这里$3^{-1}\equiv2\ (5)$）。例如$\tau=1$给出$t\equiv6\ (15)$：模型和$\SSS(t/3)\SSS((t-1)/5)V_mV_n$在$t=6,21,36,\dots$和负分支上，加上$t=\tau$（当$3\mid\tau$）和$t=0$（当$5\mid\tau$）的对角项（此处两者均非）。这正好是\texttt{face\_dispersion}构造的；其在全部十六个族和八个高度上以$0.2$--$4\%$与测量轮廓一致，是引理~\ref{l:comb}的数值面。在循环长度$5$和$6$，相同的完全正交逻辑给出\S\ref{s:five}的乘积模型梳（聚合面在$0.1$--$2.5\%$通过）。

\section{来源、研究记录与差异登记表}\label{a:record}

\subsection*{研究记录}

本文数学在交互式研究程序（2026年8月）中发展，工作纪律为先证伪：每个结构声明在采纳前注册数值检查，每个触发在前向记录上转换，完整轨迹——轮次日志、备忘录、声明分类账、修正历史、脚本——保存在复现包附带的程序存档中。链的四阶段（\S\ref{s:cert}，\S\ref{s:fm}，\S\S\ref{s:trunc}--\ref{s:four}，\S\S\ref{s:five}--\ref{s:lp}）在\S\ref{s:verif}的认证审计冻结本链前，均经过对抗审计轮次关闭。

\subsection*{塑造本文本的转换}

注册表中触发并转换的六个检查直接塑造了数学：求值传输修复（包络从不传输到极限；r61）；缩减协方差的模型减法修正（r80）；混合大小预算（两次触发，r88）；$\tau$框架教训（固定截断弧统计从不自平均；Parseval恒等式在此进入；r90）；递推框架的不完全和失误（r95，注记~\ref{rem:prog}）；以及从内部梯级移除未证明的六矩上限（定理~\ref{t:main}中引用的中间值精确消耗各自链所证明）。第29次转换（r103）是\S\ref{s:local}的Lean草稿逆遗漏。

\subsection*{认证审计差异登记表（2026年8月16日）}

\begin{itemize}
\item[D1.] \S\ref{s:fm}组装段早期草稿打印$\Lambda_2(0)=0.156293$；$\bar R=0.0066$处闭式给出$0.154193$，定理常数$0.6916/0.8458$与后者一致。在\S\ref{s:fm}中修正；定理不受影响。\emph{对主定理影响：无。}
\item[D2.] 早期草稿将带$5.573\pm0.113$，$10.02\pm0.27$归因于“在$2\times10^{7}$个zeta零点上的测量”；复现包从GUE双窗口Richardson外推（\texttt{gue\_bands.py}，\texttt{m5m6\_bands.py}）精确导出这些数字，且记录中没有该大小的zeta零点数据集。本文将这些带归因于正弦模型/GUE测量\eqref{e:guebands}。若存在zeta零点测量，它不在记录中且未被消耗。\emph{影响：归属仅；常数不变。}
\item[D3.] 消耗LP带模型：相关 vs 反相关$\delta$记账（$0.79454$ vs $0.78972$）；相关形式正确且是\texttt{face\_lp\_full}实现的；两者均记录（\S\ref{s:lp}）。\emph{影响：反相关变体是双重保守的；标题不变。}
\item[D4.] $m_4$+$m_5$梯级：早期阶段引用$0.7295/0.8647$；审计独立LP在相同约束下给出$0.73144/0.86572$（有界支撑）。被D6取代：该梯级无论如何被降级。\emph{影响：对主链无影响（降级梯级未被消耗）。}
\item[D5.] 形式化语言：早期程序文档将\cite{C26}描述为“Lean形式化”；该陈述涉及\cite{C26}附带的包，其编译工件不存在于此仓库（\S\ref{s:verif}(f)）。本文的编译覆盖精确如(f)所分级。\emph{影响：等级语言仅。}
\item[D6.] 中间$m_4$+$m_5$梯级（$0.7295/0.8647$，早期阶段）在无界支撑上没有六矩信息时不成立：审计LP随支撑范围增长退化到$13/18$（$0.73144\to0.72257$在$R=8\to200$），且五矩下界的对偶证书必然消耗$x^6$信息。梯级降级并修复（注记~\ref{r:rung}）；主链不受影响（$y_6>0$在引理~\ref{l:dual}）。\emph{影响：移除中间梯级；标题不变。}
\item[D7.] 存档四循环锚点$t_{\mathrm{adj}}=0.11717$（轮次69--71；早期草稿和\texttt{m1\_suite}引用）带有$+0.43\%$矩形法则网格偏差：精确值为$\tfrac7{60}=0.11667$（命题~\ref{p:t222}）。下游总量不受影响（$m_5$/$m_6$有理层无锚）；饱和分裂三重$0.2677/0.0156/{-0.0177}$变为精确$\tfrac4{15}/\tfrac1{60}/{-\tfrac1{60}}$（注记~\ref{r:books}），与累积量引擎独立$-0.01663$更一致。\emph{影响：有理层无影响（无锚）；修正存档锚点。}
\item[D8.] 有理识别$T_0,\dots,T_3=\tfrac3{70},\tfrac1{90},\tfrac1{180},\tfrac1{70}$，在审计分类账中记录为猜想且未被消耗，现已证明（命题~\ref{p:t222}）；$\{2,2,2\}=\tfrac{131}{420}$被精确消耗且其带贡献退役。\emph{影响：收紧$m_6$；贡献标题。}
\item[D9.] 连续求积常数$c_0=-1.4228$（r60，“计算而非拟合”）以闭式确定：$c_0=\gamma-2$（命题~\ref{p:c0}）。审计定位了存档张量筛$\gamma$数组与原始Euler闭式之间的$d=2$差异（$\widetilde\gamma_2=C_2-1$ vs $\gamma_2=C_2$，常数偏移恰好$-2$）；这是存档的奇偶均值约定（$q=2$弧块计入$\bar D$），而非错误——原始约定常数是$\gamma$，两种约定现随转换记录。斜率$1$律（此前为测量$0.99996$）精确证明。\emph{影响：N2部分处理；主定理不受影响。}
\item[D10.] 先前连通常数背后的端点网格梯触发了其细化面：在$b\ge6$且粗步长时端点求值\emph{符号翻转}（$b=6$在$dv=0.4$：$+0.036$ vs 真$-0.010$），存档采用$C_5=0.0275(5)$，$\{4,2\}=-0.0548(10)$，$\{6\}=-0.0100(12)$继承系统偏移。分组\emph{中点}网格梯（\S\ref{s:conv}），在$b=4$处相对命题~\ref{p:t222}精确校准，单调稳定并给出修正$C_5=0.0278(1)$，$\{4,2\}=-0.0552(8)$，$\{6\}=-0.0078(8)$，现经精确认证（\S\ref{s:five:conn}）（注册表第$32$次转换；第$31$次是D7的锚点修正，由同一审计独立强制）。\emph{影响：直接——这些是标题消耗的常数；修正下定理值从$0.7946$移至$0.7947$。}
\item[D11.] Lean草稿文件声称其$\mathbb Q$陈述的核心可判定性（“在$\mathbb Q$上用\texttt{decide}”）并使用了\texttt{ring}；维护者编译显示两者在核心Lean~4.33中不可用（\texttt{ring}是\texttt{mathlib}；\texttt{instDecidableEqRat}不核归约）。五个证明用\texttt{grind}修复，陈述不变；构建成功（注册表第$33$次转换；\S\ref{s:verif}(f)）。\emph{影响：证明仅；陈述不变。}
\item[D12.] 第二次编译在\texttt{cert\_nonneg}触发：\texttt{grind}将平方$(\cdot)^2$视为不透明原子（其诊断猜测原子值），无法推导完全平方非负性；用核心\texttt{Lean.Grind.OrderedRing.sq\_nonneg}修复（注册表第$34$次转换）。审计还发现已淘汰种子\texttt{SOSCertificate.lean}在构建根外静默存在，同时声明相同名称的标题定理在退役常数上；该文件从包中移除（内容吸收；草稿保存在程序存档中），构建现在覆盖目录中的每个\texttt{.lean}文件。\emph{影响：证明和打包仅；陈述不变。}
\item[D13.] 与\cite{C26}的引用关系在前向记录上反转两次：中间修订将框架吸收到正文中（当时\cite{C26}尚未公开张贴），本修订恢复\cite{C26}作为引用基础并移除重复证明，仅保留陈述（\S\ref{s:frame}）。\emph{影响：归属和阐述；数学内容不变。}
\item[D14.] 连通常数作为尖锐窗口计数统计（长球谱路径）的谱识别被其预注册门在采纳前证伪（$\kappa_2=0.3442$ vs $\tfrac13$；$\kappa_3\ne0$）——第$35$次转换。修正结构（划分-循环项上的链图族，$b=2$字典在$1-\int\mathrm{sinc}^4=\tfrac13$精确闭合）为加速引擎设计提供信息。\emph{影响：常数无影响（采纳前捕获）；结构理解增益。}
\item[D15.] 识别运行：加速梯引擎（仅等价变换加速：精确零支撑剪枝，项trie共享，核融合，否定对称性；在五个门上相对参考引擎验证，跨三种架构逐位）将每个梯扩展到七个Romberg梯级，并识别$C_5=\tfrac1{36}$，$\{6\}=-\tfrac1{126}$（$C_7=-\tfrac{17}{360}$弱一阶），全部在先前带内。常数从$0.0278(1)\to\tfrac1{36}$，$-0.0078(8)\to-\tfrac1{126}$，定理从$0.7947\to0.7957$。\emph{影响：直接——三个已用常数中的两个带压缩了十一个量级。}
\item[D16.] $\{4,2\}$旁观者对变体（用三个粗略参考梯和预注册候选$-\tfrac{23}{420}$指定）移植到加速引擎，通过四个门（C$_4$机制；逐位参考一致；存档交叉检查$-0.0588$/$-0.0558$；带），其七阶梯识别$\{4,2\}=-\tfrac{23}{420}$，命中预注册候选，分量$-\tfrac2{315},-\tfrac1{1260},-\tfrac1{252}$。所有N1常数现被识别；定理从$0.7957\to0.7960$。\emph{影响：直接——最后常数带退役；仅传输允量保留。}
\item[D17.] 残余$m_6$允量（$0.0002$），以前标记为冻结槽传输允量，在审计分类账中追溯到混合奇块类$\{3,3\}$的采用带；该类被写出为消失（在两种协议下每个网格恒为零，所有放置——\S\ref{s:five}算术奇块消失的模型面），允量退役。$M_6$成为精确有理$\tfrac{12809}{1260}$，定理从$0.7960\to0.7962$。\emph{影响：直接——消耗输入现在是端到端的精确有理数。}
\item[D18.] N1的符号认证：直接4D/5D分段积分器逐项闭包$\{4,2\}$但在$b=6$投影到每项数百年（片数乘法增长）；\S\ref{s:five:conn}的多面体提升——每项一个有理多面体体积，精确面递归，在$b=6$处约$10^{9}$倍约化——将所有$1310$项作为精确分数求值，求和为$\tfrac1{36}$，$-\tfrac{23}{420}$，$-\tfrac1{126}$精确。所有门通过（$\Phi_4=-\tfrac1{60}$通过两种方法；3D回归匹配命题~\ref{p:t222}；每项双方法一致；$\{4,2\}$逐分量重组）；预注册证伪者F-SYM-N1未触发；符号和数值值逐位一致。\emph{影响：等级仅——N1从识别移至证明；常数和定理不变。}
\item[D19.] $\{2,2,2\}$分类修正（第$38$次转换；注记~\ref{r:d19}）：已用$\tfrac{131}{420}$按交叉数分类六边形匹配；嵌套零交叉类从未单独求值，等于$\tfrac{17}{420}\ne\tfrac3{70}$（三条独立精确路径）。修正束$\{2,2,2\}=\tfrac{32}{105}$；$M_6=\tfrac{640}{63}$。由八矩配对运行的预注册$b=6$校准门捕获；两个模型侧估计量独立检测（相对旧值$55$--$110\sigma$和$11$--$21\sigma$偏离，相对新值$\sim1\sigma$）。\emph{影响：直接且有利——上约束$M_6$减小，定理升至$\tfrac{2025}{2519}$；已发表$0.7962/0.8981$消耗了较弱但为真的界，仍有效。教训：交叉数不是循环弦图的完全不变量；每个束现按二面体放置类求值。}
\item[D20.] 分类账枚举器去重错误（第$39$次转换）：仅在等大小块间有效的首部固定商被跨不等大小应用，低估$\{5,2\}$（$21\to15$）等；在计算任何值前由不变量\emph{轨道大小整除群阶}（$5\nmid|D_7|=14$）捕获，现随每个枚举报告发出。第二个独立枚举器确认所有清单。\emph{影响：常数无影响（计算前捕获）；该检查已制度化。}
\item[D21.] 梯带校准（第$40$次转换）：$\{2^4\}$梯的预注册$h^4$验收带（$3\times10^{-7}$）乐观了$3$倍（相对精确值实现偏差$8.1\times10^{-7}$）；所有未来F-V门准则修订为$3\times10^{-6}$。同一审计中发现等级$6$分母律（$2520$）在等级$8$失效（$3780\nmid2520$），解释了小分母重构在那里失败。\emph{影响：验证校准仅；全部$17$个精确值在修订带内确认。}
\item[D22.] $C_8$候选偏差（第$41$次转换）：预注册$C_8=\tfrac7{180}$（来自模型侧$\Sigma_8$的小分母重构）被精确运行偏离，$C_8=\tfrac{157}{4032}$，偏差$4.96\times10^{-5}$——在容差门内，但不在候选上。角点移动，证书以闭式重新优化（$\tfrac{17230713}{20284508}\to \tfrac{138058329}{162540559}$，即$0.849452\to0.849378$）；伴随$C_7$运行精确命中其注册值。\emph{影响：标题$-7\times10^{-5}$；教训——在等级$8$，小分母重构不可靠（见D21），只有符号运行确定候选。}
\item[D23.] 声称的$\Sigma_2$平坦窗口极小性引理被其一阶探针在进入任何草稿前证伪（温和锥形降低$m_2$）；$\lambda$最优性平坦窗口仅作为经验陈述保留。\emph{教训：最优性声明在命名前接受变分探针。}
\item[D24.] 截断小数在八处被写成四舍五入（$1-2\lambda_7=0.910460410551\ldots$排印为“$0.910460411\ldots$”，使显示不等式在该精度下为假）；全部修正为带显式省略号的截断。\emph{教训：引用小数需声明舍入或截断；不等式仅在截断方向安全。}
\item[D25.] 活动备忘录误述$b=14$表范围（$11$点，余量$3$），而存档持有$9$点（N=2..10），余量$1$；前向修正，结论不变。此外，提升系统扫描中的过度合并对称折叠（反射置换块，因此必须置换$c$列）由计数对账相对独立枚举捕获（$11/61$ vs 真$13/73$）并通过移除折叠修复。\emph{教训：每当扩展坐标时，对称去重必须重新证明；证伪扫描可能高计数，从不低计数。}
\item[D26.] $b=7$提升系统扫描的成本标定低估了七倍：探针取
枚举顺序的\emph{前}三个矩阵，而该顺序把结构简单的分拆排在
前面；另外，一次派工报告$100$核而实际仅启动$24$核。重新分片
后重跑。\emph{教训：标定探针必须随机抽样——``取前$N$个''不是
抽样——且报告的核数必须与实际派工核对。}
\item[D27.] 一份计算指令书写明装配门``带引理~\ref{l:lift}
系数''而未给出可执行公式；计算侧的问询在产出任何错误结果之前
抓住了该缺口。随后精确的系数律（合并分拆、类上循环序、符号
$\prod_B(-1)^{|Q_B|-1}$、Fubini项数$75/541/4683$于$b=4/5/6$）
被推导、在双侧逐位验证并记录于引理~\ref{l:lift}。
\emph{教训：``带标准系数''不是规格；每份装配指令书都携带显式
公式与预注册自检值。}
\item[D28.] 最终打包时发现一份回执表格把$b=5$提升系统扫描列为通过，而其工件文件不在包内（该次运行早于一次会话切换，输出从未写盘）；扫描在验收前重跑并归档。\emph{教训：日志或回执中的数字不是凭据；表格的每一行都必须能解析到已归档的工件——且此纪律必须对回执本身执行。}
\end{itemize}

\noindent 注册表第$30$次转换发生在认证审计期间：W1解析梳面的第一个版本使用$b_1b_2$的原始倍数作为同余类并触发；修正面实现每坐标CRT类（\S\ref{s:writeouts}）。前向记录在\texttt{writeouts\_verification.py}中。第$31$--$35$次是D7，D10，D11，D12和D14。

\subsection*{复现}

完整包——论文、复现管道、认证层、GPU引擎及Lean源码和构建日志——公开于
\url{https://github.com/JoshuaHKU/zeta-density-one-reproduction}。
本文附有统一复现检查表（\texttt{REPRODUCTION.md}）：环境、每套件单命令、预期输出（2026年8月16日重跑值），以及每个引用常数到其再生脚本的映射。本修订的计算资产组织为自包含包（\texttt{repro/}：引擎，每常数轨道清单和值，内容规范键带每值来源和第二路径字段，测量原始数据，每个证伪者门一个可执行脚本带断言匹配计数，以及仅追加日志存档；规范在\texttt{compute/REPRO\_SPEC.md}中）。包冻结为v1：十三个常数带每轨道记录，$126$个文件在散列清单下，十一个门脚本全部通过，记录双面验收审计（\texttt{repro/ACCEPTANCE\_r141.md}）；其后的冻结 v3（迹路线）与 v4（分支相等战役）以只增不改方式将其扩展到当前的十八门套件（\texttt{repro/REPRO\_V4\_FREEZE.md}）。全集总运行时间在一台CPU上小于十五分钟（会话包可选长门小于三十；精确$\{2,2,2\}$认证小于两小时；N1多面体认证：$\{4,2\}$和$C_5$数分钟，$\{6\}$在$9.1$单核小时内，项并行）。

\section{精确认证层}\label{a:cert}

\subsection*{A1. 精确积分算法}

每个配对层常数是积分
$\int_{[-1,1]^{d}}(1-\mathrm{spread}(p_1,\dots,p_b))_+\,
\prod_i|x_i|\,dx$，其中游走位置$p_j$是弦变量的$\{0,1\}$组合，且$0$是位置（因此支撑强制$|x_i|\le1$，$C_2$因子的$\min(|\cdot|,1)$帽从不绑定）。被积函数是分段多项式；其扭结集包含在超平面族
\[
\mathcal H=\bigl\{p_i-p_j=c\bigr\}_{i<j,\ c\in\{0,\pm1\}}
\cup\bigl\{x_i=0,\pm1\bigr\},
\]
所有系数在$\{-1,0,1\}$中，且在最内变量中首项系数为单位。迭代积分：内层积分在枚举的内层断点之间精确计算；由于首项系数为单位，断点是外层变量的小整数线性型，部分积分的片结构仅在（i）直接扭结，（ii）两个内层断点碰撞，（iii）断点-边界碰撞时改变——每一层显式枚举的结式族。在每个无扭结片上，（部分）被积函数是次数至多$2/4/8$（内/中/外变量）的多项式，通过阶$5/7/9$的闭Newton--Cotes公式在有理节点处精确积分，全程精确有理算术。\emph{自检}：每片在同一规则下对整片及其两半积分；对次数低于规则阶的多项式两者都是精确积分，因此两个有理数必须\emph{相同}，枚举遗漏的扭结会在包含片上导致不一致。在任何运行的任何级别未发生不一致。独立交叉检查：闭式卦限求值（A2），以及存档的浮点网格梯（相对认证有理数$10^{-6}$一致）。

\subsection*{A2. 闭式交叉检查}

$T_0$（位置$\{0,v,w,u\}$）：在全正卦限$\int_{[0,1]^3}(1-\max(v,w,u))vwu=\tfrac1{56}$（排序并积分）；一个负卦限给出$\int_0^1t(1-t)^5/20\,dt=\tfrac1{840}$；对轨道求和，$T_0=2\cdot\tfrac1{56}+6\cdot\tfrac1{840}=\tfrac3{70}$。
$t_{\mathrm{opp}}$（位置$\{0,v,v{+}w,w\}$）：相同分解给出$\tfrac1{30}$精确，确认存档精确锚点。
$t_{\mathrm{adj}}$：两个同号卦限各给$\tfrac1{20}$，两个混合卦限各给$\tfrac1{120}$：$t_{\mathrm{adj}}=\tfrac{14}{120}=\tfrac7{60}$（非存档网格值$0.11717$；注册D7）。

\subsection*{A3. 认证值}

\[
t_{\mathrm{adj}}=\tfrac7{60},\quad
t_{\mathrm{opp}}=\tfrac1{30},\quad
T_0=\tfrac3{70},\quad
T_1=\tfrac1{90},\quad
T_2=\tfrac1{180},\quad
T_3=\tfrac1{70},\quad
\{2,2,2\}=2T_0+3T_0'+6T_1+3T_2+T_3=\tfrac{32}{105},
\qquad T_0'=\tfrac{17}{420}\ (\text{嵌套；D19}),
\]
且$\Phi_4=-\tfrac1{60}$，饱和$\{2,2\}=\tfrac4{15}$，平台$=\tfrac1{60}$，连通模型$=-\tfrac1{60}$（注记~\ref{r:books}）。代码：
\texttt{certification/exact\_t222.py}（类分别运行；双方案逐片相同有理数检查）。

\subsection*{A4. 连续求积常数的闭式（N2，部分处理）}

\begin{proposition}[$c_0=\gamma-2$]\label{p:c0}
令$\gamma_d$表示部分和恒等式（引理~\ref{l:PS}）在$(1,1)$类处的自由系数。则$\gamma_d$具有精确Euler闭式
\[
\gamma_{2m}=C_2\prod_{p\mid m}\frac1{p(p-2)}
\quad(m\ \text{奇平方自由}),\qquad
\gamma_d=0\ \text{否则},
\]
且在原始约定下，
$\sum_{d\le M}d\gamma_d=\log M+\gamma+o(1)$；在存档约定（$q=2$弧块计入奇偶均值，$\widetilde\gamma_2=\gamma_2-\beta_2=C_2-1$）下，
\[
c_0\;=\;\gamma-2\;=\;-1.4227843\ldots
\]
\end{proposition}

\begin{proof}
(i)~闭式：由弧恒等式（命题~\ref{p:S}）在$b_1=b_2=1$处，$\beta_q=\mu(q)^2/\varphi(q)^2$，且$\gamma_d=\sum_e\mu(e)\beta_{de}$逐素数分解，局部权重$f_0(p)=1-(p-1)^{-2}$，$f_1(p)=(p-1)^{-2}$，$f_{v\ge2}=0$；$f_0(2)=0$强制$2\,\|\,d$，乘积组装为显示（$C_2$为孪生素数常数）。
(ii)~常数：$\sum_{d\le M}d\gamma_d =2C_2\sum_{m\le M/2}\prod_{p\mid m}\tfrac1{p-2}$，其Dirichlet级数为$\zeta(s+1)E(s)$，$E(0)=\prod_p(1-\tfrac1p)\prod_{p>2}(1+\tfrac1{p-2}) =\tfrac1{2C_2}$（斜率$2C_2E(0)=1$精确），且$E'(0)/E(0)=\sum_p\tfrac{\log p}{p-1} -\sum_{p>2}\tfrac{\log p}{p-1}=\log2$（$p$和除$p=2$外成对相消）；二重极点留数给出$2C_2[E(0)(\log\tfrac M2+\gamma)+E'(0)]=\log M+\gamma$。
(iii)~约定偏移：存档张量筛（\texttt{tail\_bound}）与闭式在除$d=2$外每个$d$一致，在$d=2$存档类$d$替换记$\widetilde\gamma_2=C_2-1$（$q=2$块移至奇偶均值）；这使$\sum d\gamma_d$偏移恰好$-2$。数值验证（\texttt{n2\_c0\_certification.py}）：原始和收敛到$\gamma$，在$M=10^4/10^5/2\times10^7$处偏差$9.4\times10^{-4}/6.3\times10^{-5}/9.2\times10^{-6}$；存档测量$-1.42279$在$M=1.6\times10^{7}$处匹配$\gamma-2$至$10^{-5}$；部分和恒等式以其$d>M$尾部闭合（在$M=2\times10^5$，$D=2\times10^{7}$处$\mathrm{LHS}-\mathrm{RHS}-\mathrm{tail}=-0.010$，与剩余截断一致）。
\end{proof}

\noindent 因此引理~\ref{l:CL}的“计算常数”简化为：斜率$=1$（证明），$c_0=\gamma-2$（闭式），已认证几何尾部$T_\Gamma$（存档r60，比$0.500$），以及$c^{*}$内部的唯一剩余计算部分$\overline{\mathrm{osc}}$。N2因而部分处理；剩余的是$\overline{\mathrm{osc}}$识别（或认证包络）和普适性坍缩速率写出。

\subsection*{A5. N1证书族及其收益曲线}

精确对偶证书针对任何假设认证包络$C_5\in[c_-,c_+]$，$\{4,2\}\le u_{42}$，$\{6\}\le u_6$重新优化（相关记账：单个未知$C_5$同时进入$m_5$和$m_6$；最坏情况是仿射族的角点）。精确有理证书（\texttt{certificate\_family.py}，原子按场景重新优化）：

\begin{center}
\begin{tabular}{lcc}
场景（按消耗信息分级） &
$1-2w_0\ge$ & $1-w_0\ge$\\
\hline
$k\le6$精确认证（推论~\ref{c:2025}） &
$2025/2519$ & $2272/2519$\\
$k=8$层以所述等级（引理~\ref{l:dual8}） &
$0.84937$ & $0.92468$\\
包络宽度$0.01$（历史） & $0.77418$ & $0.88709$\\
包络宽度$0.04$（历史） & $0.74299$ & $0.87149$\\
仅符号：$C_5\ge0$，$\{4,2\}\le0$，$\{6\}\le0$ &
$0.74031$ & $0.87016$\\
无连通信息 & $13/18$ & $31/36$\\
\end{tabular}
\end{center}

\noindent 已交付证书的敏感性（精确）：
$\partial(\text{标题})/\partial c_-=+7.32$，
$\partial/\partial c_+=-5.61$，
$\partial/\partial(u_{42}{+}u_6)=-0.93$。特别地，三个连通常数的\emph{仅符号}认证——比包络粗得多——已认证$0.740/0.870$，且精确计算（现已完成，落在顶行）实现完整收益。（将直接分段运行置于簇尺度放置的成本模型，正是\S\ref{s:five:conn}的多面体提升所绕过的：认证运行在单工作站核心上数小时完成。）

\subsection*{A6. 精确对偶证书及历史原子来源}

引理~\ref{l:dual}的度$6$证书是$P=(Q_3^*)^2$，有理核系数显示；其角点求值$\lambda_3=\tfrac{247}{2519}$和引理~\ref{l:dual8}的度$8$角点$\lambda_4=\tfrac{12241115}{162540559}$是有限有理计算（\texttt{certification/certify84.py}；Lean排队）。$Q_3^*$的分子正比于
$1932x^3-7368x^2+8232x-2519$，其根$0.50151\ldots$，$1.27817\ldots$，$2.03397\ldots$是Christoffel极值测度的证书“原子”——证明从不需要。\emph{历史原子来源}：在被替代$M_6=\tfrac{12809}{1260}$下相同构造返回根$0.53234\ldots$，$1.31220\ldots$，$2.05855\ldots$，其四位有理化正是早期修订的原子$(\tfrac{5323}{10000},\tfrac{6561}{5000},\tfrac{10293}{5000})$——LP提取一直在找核多项式根。记录显示，该早期证书展开（分子，幺正）为
\[
x^6-\tfrac{39031}{5000}x^5+\tfrac{2422533313}{100000000}x^4
-\tfrac{4746141355593}{125000000000}x^3
+\tfrac{39293368568783721}{1250000000000000}x^2
-\tfrac{20200560698400422913}{1562500000000000000}x
+(abc)^2,
\]
其中$(abc)^2=\tfrac{129222147277358919747441}{62500000000000000000000}$，故$y_0=1$，$y_5=-\tfrac{39031}{5000}(abc)^{-2}<0$，$y_6=(abc)^{-2}>0$。在被替代角点给出$w_0=\tfrac{829278553005924403328783}{8140995278473611944088783}$及已发表标题$\tfrac{7962}{10000}$/$\tfrac{8981}{10000}$——当时有效且现在仍有效（消耗较弱但为真的$M_6$），被引理~\ref{l:dual}和引理~\ref{l:dual8}的精确最优值取代。在Christoffel公式下，求解器、网格、原子提取和有理化步骤不再以任何形式存在；整个消耗层是引理~\ref{l:christoffel}的有理线性代数。验证：
\texttt{certification/certify84.py}；Lean：
\texttt{lean/RhGate/Certificate.lean}（编译，先前角点）和\texttt{Certificate84.lean}（排队）：展开恒等式、非负性、角点求值、符号模式、标题；不含任何未完成的证明占位（\texttt{sorry}）。

\begin{thebibliography}{99}

\bibitem{Akh} N.~I.~Akhiezer, \emph{The Classical Moment
Problem and Some Related Questions in Analysis}, Oliver \&
Boyd, 1965.

\bibitem{BGST} S.~A.~C.~Baluyot, D.~A.~Goldston,
A.~I.~Suriajaya, and C.~L.~Turnage-Butterbaugh, \emph{Pair
correlation of zeros of the Riemann zeta function I: proportions
of simple zeros and critical zeros}, arXiv:2501.14545 (2025).

\bibitem{GS} D.~A.~Goldston and A.~I.~Suriajaya, \emph{Zeta zeros
on the critical line}, arXiv:2511.20059 (2025).

\bibitem{BCY11} H.~M.~Bui, J.~B.~Conrey, and M.~P.~Young,
\emph{More than 41\% of the zeros of the zeta function are on the
critical line}, Acta Arith.\ \textbf{150} (2011), 35--64.

\bibitem{C26} Claude (Anthropic), \emph{More than two thirds
of the zeros of the Riemann zeta function lie on the critical
line}, preprint, August 2026,
\url{https://www-cdn.anthropic.com/564f962e60643842f5fcb4a17c9dbc8f608f1c37.pdf}.

\bibitem{Con89} J.~B.~Conrey, \emph{More than two fifths of the
zeros of the Riemann zeta function are on the critical line},
J.\ reine angew.\ Math.\ \textbf{399} (1989), 1--26.

\bibitem{Dav} H.~Davenport, \emph{Multiplicative Number Theory},
3rd ed., Springer GTM 74, 2000.

\bibitem{Fen12} S.~Feng, \emph{Zeros of the Riemann zeta function
on the critical line}, J.\ Number Theory \textbf{132} (2012),
511--542.

\bibitem{HR} H.~Halberstam and H.-E.~Richert, \emph{Sieve
Methods}, Academic Press, 1974.

\bibitem{IK} H.~Iwaniec and E.~Kowalski, \emph{Analytic Number
Theory}, AMS Colloquium Publ.\ 53, 2004.

\bibitem{Lev74} N.~Levinson, \emph{More than one third of zeros
of Riemann's zeta-function are on $\sigma=1/2$}, Adv.\ Math.\
\textbf{13} (1974), 383--436.

\bibitem{MRT} K.~Matom\"aki, M.~Radziwi\l{}\l, and T.~Tao,
\emph{Correlations of the von Mangoldt and higher divisor
functions I. Long shift ranges}, Proc.\ London Math.\ Soc.\ (3)
\textbf{118} (2019), 284--350.

\bibitem{Mon73} H.~L.~Montgomery, \emph{The pair correlation of
zeros of the zeta function}, Proc.\ Sympos.\ Pure Math.\
\textbf{24} (1973), 181--193.

\bibitem{Mon75} H.~L.~Montgomery, \emph{Distribution of the zeros
of the Riemann zeta function}, Proc.\ ICM Vancouver 1974, vol.~1,
379--381.

\bibitem{PRZZ} K.~Pratt, N.~Robles, A.~Zaharescu, and
D.~Zeindler, \emph{More than five-twelfths of the zeros of
$\zeta$ are on the critical line}, Res.\ Math.\ Sci.\ \textbf{7}
(2020), no.~2.

\bibitem{Sim98} B.~Simon, \emph{The classical moment problem
as a self-adjoint finite difference operator}, Adv.\ Math.\
\textbf{137} (1998), 82--203.

\bibitem{Sos00} A.~Soshnikov, \emph{Determinantal random
point fields}, Russian Math.\ Surveys \textbf{55} (2000),
923--975.

\bibitem{Sel42} A.~Selberg, \emph{On the zeros of Riemann's
zeta-function}, Skr.\ Norske Vid.-Akad.\ Oslo I (1942), no.~10.

\bibitem{Tit86} E.~C.~Titchmarsh, \emph{The Theory of the
Riemann Zeta-Function}, 2nd ed.\ (revised by D.~R.~Heath-Brown),
Oxford University Press, 1986.

\bibitem{Wei52} A.~Weil, \emph{Sur les ``formules explicites''
de la th\'eorie des nombres premiers}, Comm.\ S\'em.\ Math.\
Univ.\ Lund (1952), 252--265.

\end{thebibliography}
\end{document}